\chapter{信息抽取与知识构建}
\label{chap:nlp-information-extraction}

读出通知中的“4月18日”，还不能直接填写活动日期：原文也可能说“原定4月18日，现已取消”。找出两个组织名称，也不能据此建立合作关系。上一章的词法、句法、语义和篇章分析帮助我们辨认语言结构，本章进一步问：怎样把这些结构转成可查询、可核验，并能随新材料更新的记录？

沿用本部分的教学通知：活动编号为\code{ML-0418}，通知发布于2025年4月15日，时间均为北京时间。编号和发布日期是给定元数据，不是从下列正文中猜出的字段。
\begin{quote}
机器学习读书会原定于4月17日举行，现改至4月18日（周五）14:00，地点为海棠楼201。校内师生可直接参加，校外来宾须在4月16日18:00前提交报名信息。
\end{quote}
本章所有通知、知识库和数值例子均为教学构造，不表示真实活动或模型实测结果；需要额外的组织、人物或业务状态时，会在当地明确给出补充材料。

\term{信息抽取}（information extraction, IE）按照给定的类型、角色和字段规范，把原文表达的对象与事实恢复为结构化结果。这份规范通常称为\term{模式}（schema），它规定本次要抽取哪些实体类型、关系和事件角色；规定对象类型、关系含义和角色约束的概念体系称为\term{本体}（ontology），模式通常取本体的一部分。模式可以要求输出地点，却不能替原文补出地点；模型判断“改至”指向新日期时，还要留下它依据哪个句子、哪个位置。

从这段通知到一组可用的记录，要依次跨过几类困难。表~\ref{tab:nlp-extraction-challenges}把它们列成清单，并给出各节的方法谱系。第一组困难出在实体本身：实体可以嵌套或不连续，中文没有词边界，新类型缺少标注，同一名字又可能指不同对象。第二组出在实体之间：候选实体对的数量随实体数平方增长，多个三元组共享实体时互相重叠，远程监督自动生成的关系标签带噪，关系的证据又常分散在不同句子。第三组出在事件：触发词与论元要绑定到同一次事件，同一文档的多个事件共享论元，事件之间的时间与因果关系还需全局一致。最后，多来源的记录需要对齐、消解冲突，并按时间更新。

\begin{table*}[htbp]
  \centering
  \small
  \begin{tabularx}{\linewidth}{@{}>{\raggedright\arraybackslash}p{40mm}>{\raggedright\arraybackslash}X>{\raggedright\arraybackslash}p{14mm}@{}}
    \toprule
    挑战 & 方法谱系 & 所在节\\
    \midrule
    实体嵌套与不连续；中文缺少词边界 & BIO--CRF与分层标注$\to$词汇增强（Lattice-LSTM、FLAT）$\to$跨度与阅读理解式$\to$超图与词对网格$\to$指针与结构化生成 & \ref{sec:nlp-extraction-entities}\\
    新类型标注稀缺 & 原型网络$\to$近邻标注与抽象转移$\to$大语言模型提示与实体级投票 & \ref{sec:nlp-extraction-entities}\\
    同名实体身份歧义；库外实体 & 别名先验$\to$双编码召回与交叉编码重排$\to$文档全局一致性$\to$生成式名称预测 & \ref{sec:nlp-extraction-linking}\\
    实体对平方增长；三元组重叠 & 流水线$\to$统一标注$\to$复制生成$\to$级联指针$\to$词对表格填充 & \ref{sec:nlp-extraction-relations}\\
    远程监督标签带噪 & 多示例学习与句子级注意力 & \ref{sec:nlp-extraction-relations}\\
    关系与事件论元跨句分布 & 文档级图与上下文池化；文档级事件表填充 & \ref{sec:nlp-extraction-relations}、\ref{sec:nlp-extraction-events}\\
    同一文档多事件共享论元 & 流水线$\to$联合打分与束搜索$\to$约束生成$\to$文档级路径扩展 & \ref{sec:nlp-extraction-events}\\
    事件间时间与因果需全局一致 & 局部分类$\to$全局约束推断与联合学习 & \ref{sec:nlp-extraction-event-relations}\\
    多来源记录的对齐、冲突与更新 & 记录匹配$\to$合并与冲突处理$\to$知识图谱嵌入补全$\to$核验与版本更新 & \ref{sec:nlp-extraction-records}\\
    \bottomrule
  \end{tabularx}
  \caption{本章的挑战清单与方法谱系。箭头表示每一步针对前一方法的一项限制，不表示后一方法在所有条件下都更好。所在节一列给出主要讨论的位置。}
  \label{tab:nlp-extraction-challenges}
\end{table*}

各节的方法沿一条共同的主线推进：流水线把任务拆成依次执行的子任务，简单但误差逐级传递；联合抽取让实体、关系或事件共享表示并统一解码；生成式统一抽取把各类结构写成同一种目标序列；少样本方法与大语言模型提示则应对标注稀缺。这条主线贯穿实体、关系与事件三节，读者可以在各节对照同一种思路在不同对象上的形式。

这些方法的通用机制已在第\ref{sec:nlp-analysis-paradigms}节的五种结构预测范式中写出，本章只写任务特有的部分。标注类方法（BIO实体、触发词标注、统一标注的关系抽取）沿用第\ref{sec:nlp-analysis-paradigm-tagging}小节的合法转移掩码与CRF；跨度类方法（双仿射实体、阅读理解式实体、实体对与论元跨度分类）沿用第\ref{sec:nlp-analysis-paradigm-span}小节的跨度表示与冲突解码；超图、词对网格和表格填充属于第\ref{sec:nlp-analysis-paradigm-graph}小节的图方法，结构得分分解到节点、边或网格单元上；指针生成、结构化生成、生成式实体链接与事件生成属于第\ref{sec:nlp-analysis-paradigm-generation}小节的约束生成，其约束解码的掩码计算见第\ref{sec:nlp-inference-grammar}小节。范式只说明输出怎样表示和解码；同一任务常有多种范式的实现，比较时应固定编码器与评价口径。

\section{实体识别}
\label{sec:nlp-extraction-entities}

实体识别从原文中找出指称对象的片段并给出类型，例如通知中的“海棠楼201”是地点，“海棠大学计算机学院”是组织。它的输出是本章其余任务的输入：实体链接为提及确定知识库身份，关系抽取在提及之间判断关系，事件抽取把提及填入论元角色。本节先说明实体结构有哪些形态、各路方法能表达哪些形态，再按结构与数据限制依次讨论序列标注、跨度、超图与词对网格、生成、少样本与提示等路线。

\subsection{实体结构的挑战与方法谱系}
\label{sec:nlp-extraction-mentions}

“海棠楼201”涉及三种不同的产出。它在原文中占据一段字符，这是\term{实体提及}（entity mention）；补充句中的“该教室”可能与它指向同一对象，这是文档内共指（第\ref{sec:nlp-analysis-span-coreference}小节）；把它对应到校园场所库中的某个编号，才是实体链接（第\ref{sec:nlp-extraction-linking}节）。图~\ref{fig:nlp-extraction-identity}用一段补充通知和一个小型知识库画出三者的关系。本节只处理第一种产出。

\begin{figure*}[htbp]
  \centering
  % Adapted from academic-drawing's TikZ template; project fonts and palette prevail.
% Final canvas: 167 mm wide. Ports: text east -> spans west; spans east
% -> cluster west; cluster east -> KB west. Side branches use a routing lane.
\begin{tikzpicture}[
  x=1mm,y=1mm,
  font=\figurefont\fontsize{9}{11}\selectfont,
  box/.style={rectangle,align=left,inner sep=2mm,line width=0.85pt,
    draw=BookRule,fill=white},
  source/.style={box,text width=31mm,draw=FigureInput,fill=FigureInput!8},
  span/.style={box,text width=27mm,draw=FigureInput,fill=FigureInput!8},
  cluster/.style={box,text width=23mm,draw=FigureModel,fill=FigureModel!8},
  target/.style={box,text width=33mm,draw=FigureOutput,fill=FigureOutput!8},
  locate/.style={draw=FigureInput,line width=0.85pt,rounded corners=1.2mm},
  coref/.style={draw=FigureModel,line width=0.85pt,dashed,rounded corners=1.2mm},
  link/.style={draw=FigureOutput,line width=1.1pt,-{Latex[length=2mm]}},
  head/.style={font=\figurefont\bfseries\fontsize{9}{11}\selectfont,
    text=BookInk,align=center}
]
  \node[head] at (18,37) {补充原文};
  \node[head] at (62,37) {原文跨度};
  \node[head] at (104,37) {文档共指簇};
  \node[head] at (146,37) {教学知识库};
  \node[source] (text) at (18,0) {
    地点为海棠楼201。\\
    该教室旁设临时服务点。\\[3pt]
    字符从0开始计数
  };
  \node[span] (m1) at (62,23) {海棠楼201\\$[3,9)$};
  \node[span] (m2) at (62,0) {该教室\\$[10,13)$};
  \node[span] (m3) at (62,-23) {临时服务点\\$[15,20)$};
  \node[cluster] (c1) at (104,11.5) {同一教室\\$\{m_1,m_2\}$};
  \node[cluster] (c2) at (104,-23) {服务点\\$\{m_3\}$};
  \node[target] (kb1) at (146,11.5) {ROOM-201\\海棠楼201};
  \node[target] (kb2) at (146,-23) {NIL\\无可确认库中项};
  \draw[locate] (text.east) -- (42,0);
  \draw[locate] (42,23) -- (42,-23);
  \draw[locate] (42,23) -- (m1.west);
  \draw[locate] (42,0) -- (m2.west);
  \draw[locate] (42,-23) -- (m3.west);
  \draw[coref] (m1.east) -- (84,23) -- (84,11.5) -- (c1.west);
  \draw[coref] (m2.east) -- (84,0) -- (84,11.5);
  \draw[coref] (m3.east) -- (c2.west);
  \draw[link] (c1.east) -- (kb1.west);
  \draw[link] (c2.east) -- (kb2.west);
  \draw[locate] (4,-38) -- (14,-38);
  \node[anchor=west] at (16,-38) {定位};
  \draw[coref] (49,-38) -- (59,-38);
  \node[anchor=west] at (61,-38) {共指归组};
  \draw[link] (111,-38) -- (121,-38);
  \node[anchor=west] at (123,-38) {身份选择};
\end{tikzpicture}

  \caption{实体提及、共指与身份链接。教学补充材料为“地点为海棠楼201。该教室旁设临时服务点”。原文到跨度的实线表示定位，跨度到共指簇的虚线表示同指归组，簇到知识库的箭头表示身份选择。临时服务点的NIL表示当前知识库没有可确认的对应项，不表示该对象不存在。}
  \label{fig:nlp-extraction-identity}
\end{figure*}

\term{命名实体识别}（named entity recognition, NER），本书简称实体识别，要同时回答“哪里是提及”和“属于什么类型”。沿用第\ref{chap:nlp-introduction}章的位置约定，原文字符区间写成从零开始的半开区间$[s,e)$；模型内部的词元从1开始编号，词元跨度$(i,j)$包含两端。若词元$i$对应原文$[a_i,b_i)$，连续跨度$(i,j)$恢复为$[a_i,b_j)$。子词切分和全半角转换都必须保留这张偏移表，否则模型给出的索引不能落回原文。

实体有四种结构形态，它们决定了输出空间。\emph{平面}实体互不重叠，一句话的实体集合可以写成一组不相交的区间。\emph{嵌套}实体允许一个区间包含另一个：“海棠大学计算机学院”是组织，其中“海棠大学”也是组织，这是\term{嵌套提及}（nested mention）。\term{不连续提及}（discontinuous mention）由几个不相邻的片段共同表达一个对象：在“左和右肺均清晰”中，若标注规范把省略了中心词的“左……肺”记为一个身体部位，它就由“左”和“肺”两段组成，输出应是有序片段列表，而不是覆盖中间文字的大区间。第四种困难与结构形式无关，而与输入单位有关：中文没有空格，字符模型不受分词错误影响，却丢失了“海棠楼”是一个词这类信息；词模型有词信息，但分词一旦切错，实体边界就无法恢复。

表~\ref{tab:nlp-extraction-ner-lineage}把本节讨论的方法按能表达的结构排开。从上到下，输出表示从逐位置标签变为区间、超图、词对网格，最后变为可线性化的序列，能表达的结构逐步变多；推断代价则从线性升到平方，生成方法还要承担串行解码与原文对齐。范式一列给出第\ref{sec:nlp-analysis-paradigms}节中的归类。

\begin{table*}[htbp]
  \centering
  \small
  \begin{tabularx}{\linewidth}{@{}>{\raggedright\arraybackslash}p{27mm}>{\raggedright\arraybackslash}p{15mm}*{3}{>{\centering\arraybackslash}p{10mm}}>{\centering\arraybackslash}p{12mm}>{\raggedright\arraybackslash}X@{}}
    \toprule
    方法 & 范式 & 平面 & 嵌套 & 不连续 & 词典信息 & 推断与代价\\
    \midrule
    BIO--CRF & 序列标注 & 能 & 否 & 否 & 不用 & 带类型掩码的Viterbi，$O(nL^2)$\\
    分层标注 & 序列标注 & 能 & 限层数 & 否 & 不用 & 逐层Viterbi，层数随嵌套深度增加\\
    Lattice-LSTM、FLAT & 序列标注 & 能 & 否 & 否 & 用 & Lattice-LSTM难并行；FLAT注意力为$O((n+m)^2)$\\
    双仿射跨度分类 & 跨度 & 能 & 能 & 否 & 不用 & $O(n^2)$候选，平面任务需冲突解码\\
    阅读理解式 & 跨度 & 能 & 能 & 否 & 不用 & 每个类型单独编码一次\\
    超图 & 图方法 & 能 & 能 & 扩展后有限 & 不用 & 超图inside递推，$O(n|T|)$；有结构歧义\\
    W2NER词对网格 & 图方法 & 能 & 能 & 能 & 不用 & $n\times n$网格分类，路径解码\\
    BARTNER、UIE & 约束生成 & 能 & 能 & 能 & 不用 & 串行解码，输出需解析并对齐原文\\
    \bottomrule
  \end{tabularx}
  \caption{实体识别的方法谱系。$n$为词元数，$L$为标签数，$|T|$为实体类型数，$m$为词典匹配到的词数。“词典信息”指是否显式接入中文词典词，不用词典的方法以字或子词为单位同样可以运行；“扩展后有限”指原始超图只表示连续提及，加入间隙状态的扩展可表示部分不连续提及。}
  \label{tab:nlp-extraction-ner-lineage}
\end{table*}

评价采用严格实体级F1：预测集合$\hat{\mathcal E}$与金标准集合$\mathcal E^\star$的元素是“片段列表加类型”，两者完全相同才算命中，$P=|\hat{\mathcal E}\cap\mathcal E^\star|/|\hat{\mathcal E}|$，$R=|\hat{\mathcal E}\cap\mathcal E^\star|/|\mathcal E^\star|$，$F_1=2PR/(P+R)$。把“海棠楼201”预测成“海棠楼”或预测成组织，都不计真阳性。细粒度和少样本设定另有基准：Few-NERD标注了8个粗类与66个细类，并按类型划分训练与测试，用来检验模型能否在少量样本下识别训练时没见过的细类\scite{ding2021fewnerd}。

\subsection{基于序列标注的实体识别}
\label{sec:nlp-extraction-ner-tagging}

序列标注是实体识别最早成熟的路线。输出表示、合法转移掩码、CRF训练与Viterbi解码已在第\ref{sec:nlp-analysis-paradigm-tagging}小节写出，BiLSTM-CRF的逐层运算见第\ref{sec:nlp-analysis-crf-chunk}小节；本小节只写实体任务特有的标签集、掩码与中文词汇增强。这条路线训练和解码都是线性时间，但一层标签不能表示嵌套与不连续，中文字符模型又缺少词边界信息；词汇增强针对后一问题。

\subsubsection{BIO--CRF与BiLSTM-CRF实体标注}
\label{sec:nlp-extraction-bio}

实体标注与分词、组块的差别在于标签带类型。设实体类型集合为$T$，BIO标签集为$\{\mathrm B\text{-}X,\mathrm I\text{-}X:X\in T\}\cup\{\mathrm O\}$，共$L=2|T|+1$个；BIOES再加E-$X$与S-$X$，共$4|T|+1$个。类型使掩码多了一条约束：延续标签必须与它所接的开头同类型。BIO规范下的转移掩码与起始掩码为
\begin{equation}
  \bm M_{a,b}=
  \begin{cases}
    -\infty, & b=\mathrm I\text{-}X\ \text{且}\ a\notin\{\mathrm B\text{-}X,\mathrm I\text{-}X\},\\
    0, & \text{其他},
  \end{cases}
  \qquad
  m^{\mathrm{start}}_{\mathrm I\text{-}X}=-\infty,\quad X\in T,
  \label{eq:nlp-extraction-crf}
\end{equation}
其余起始项与全部结束项为0。把它代入式~\eqref{eq:nlp-analysis-legal-mask}，整句得分、CRF负对数似然和Viterbi解码不需要任何改动：O之后不能直接接I-ORG，B-LOC之后也不能接I-ORG，这两类序列在配分函数和$\argmax$中都自动消失。掩码与标签集一样是规范的一部分，训练和解码必须使用同一份。

发射分数的来源经历了三代：特征CRF使用字形、词典与上下文窗口特征；BiLSTM-CRF用词嵌入、字符级表示与双向LSTM取代人工特征\scite{nlp-analysis-lample2016,v9a-ma2016}；预训练编码器再以上下文表示取代前两者，按首子词对齐收集词级表示（第\ref{sec:nlp-classical-bert-heads}小节）。三代模型的CRF层完全相同，改变的只是$\bm h_i$从哪里来。

以教学切分“海棠楼／201／开放”为例，正确标签是B-LOC、I-LOC、O，恢复“海棠楼201”$[0,6)$。若某位置的发射分数使逐位置取最大得到O、I-LOC、O，第二个标签不合法；掩码下的Viterbi只能在B-LOC、I-LOC、O与B-LOC、O、O等合法序列中选择。B-ORG、I-ORG、O边界正确而类型错误，B-LOC、O、O类型正确却漏了房间号，二者在严格F1中都不计命中。若输入只被切成一个词元“海棠楼201开放”，任何标签都表示不出正确边界，这时需要字符级输入或下一小节的词汇增强。

一层标签不能表示嵌套。Ju等的分层模型把多个平面标注层叠起来：第一层BiLSTM-CRF识别最内层实体，把每个识别出的实体内的词表示取平均、合并成一个单元，作为下一层的输入，再识别包含它的外层实体，直到某层不再输出实体为止\scite{ju2018layered}。它保留了线性链解码，代价是层间误差会传递，嵌套深度受层数限制，同一跨度属于多个类型的情形也难以表示。

\subsubsection{中文词汇增强：Lattice-LSTM与FLAT}
\label{sec:nlp-extraction-lattice}

中文实体识别有两种输入单位可选。词模型先分词再标注，若“海棠楼201开放”被切成“海棠楼／201开放”，实体的右边界落在第二个词内部，任何词级标签都恢复不出“海棠楼201”；字符模型没有这种传播错误，却看不到“海棠楼”是一个已知地名。第\ref{sec:nlp-analysis-pretrained-lattice}小节介绍了词格的思想：把词典匹配到的词作为额外单元接入字符序列，让语境决定采用哪个候选词。这里写出两种实现的完整计算。

第一种实现是Zhang与Yang的Lattice-LSTM，它在字符LSTM上加入词单元\scite{nlp-analysis-zhang2018lattice}。设句子有$n$个字，第$j$个字的嵌入为$\bm x^{\mathrm c}_j\in\mathbb R^{d_{\mathrm c}}$，隐藏维度为$d_{\mathrm h}$；词典$D$中从第$b$个字到第$e$个字的匹配词记为$w_{b,e}$，嵌入为$\bm x^{\mathrm w}_{b,e}\in\mathbb R^{d_{\mathrm w}}$；以第$j$个字结尾的匹配词起点集合为$B_j=\{b:w_{b,j}\in D\}$。字符单元先按普通LSTM算出各门与候选更新：
\begin{equation}
  \begin{bmatrix}\bm i^{\mathrm c}_j\\ \bm o^{\mathrm c}_j\\ \bm f^{\mathrm c}_j\\ \tilde{\bm c}^{\mathrm c}_j\end{bmatrix}
  =\begin{bmatrix}\sigma\\ \sigma\\ \sigma\\ \tanh\end{bmatrix}
  \left(\bm W^{\mathrm c}\begin{bmatrix}\bm x^{\mathrm c}_j\\ \bm h^{\mathrm c}_{j-1}\end{bmatrix}+\bm b^{\mathrm c}\right),
  \qquad \bm W^{\mathrm c}\in\mathbb R^{4d_{\mathrm h}\times(d_{\mathrm c}+d_{\mathrm h})}.
  \label{eq:nlp-extraction-lattice-char}
\end{equation}
每个匹配词$w_{b,e}$有自己的词单元。它以词嵌入和起点字的隐藏状态为输入，没有输出门，因为词单元不直接产生隐藏状态，只向词尾字输送记忆：
\begin{equation}
  \begin{bmatrix}\bm i^{\mathrm w}_{b,e}\\ \bm f^{\mathrm w}_{b,e}\\ \tilde{\bm c}^{\mathrm w}_{b,e}\end{bmatrix}
  =\begin{bmatrix}\sigma\\ \sigma\\ \tanh\end{bmatrix}
  \left(\bm W^{\mathrm w}\begin{bmatrix}\bm x^{\mathrm w}_{b,e}\\ \bm h^{\mathrm c}_{b}\end{bmatrix}+\bm b^{\mathrm w}\right),
  \qquad
  \bm c^{\mathrm w}_{b,e}=\bm f^{\mathrm w}_{b,e}\odot\bm c^{\mathrm c}_{b}+\bm i^{\mathrm w}_{b,e}\odot\tilde{\bm c}^{\mathrm w}_{b,e},
  \label{eq:nlp-extraction-lattice-word}
\end{equation}
其中$\bm W^{\mathrm w}\in\mathbb R^{3d_{\mathrm h}\times(d_{\mathrm w}+d_{\mathrm h})}$。在词尾字$j$处，每个词单元再得到一个额外输入门$\bm i^{\mathrm c}_{b,j}=\sigma(\bm W^{\mathrm l}[\bm x^{\mathrm c}_j;\bm c^{\mathrm w}_{b,j}]+\bm b^{\mathrm l})$，它衡量这个词对当前字有多大贡献。字符单元的记忆由词单元与自身候选按门控加权合并：
\begin{equation}
  \bm c^{\mathrm c}_j=\sum_{b\in B_j}\bm\alpha^{\mathrm c}_{b,j}\odot\bm c^{\mathrm w}_{b,j}+\bm\alpha^{\mathrm c}_j\odot\tilde{\bm c}^{\mathrm c}_j,
  \qquad
  \bm\alpha^{\mathrm c}_{b,j}=\frac{\exp(\bm i^{\mathrm c}_{b,j})}{\exp(\bm i^{\mathrm c}_j)+\sum_{b'\in B_j}\exp(\bm i^{\mathrm c}_{b',j})},
  \label{eq:nlp-extraction-lattice-merge}
\end{equation}
$\bm\alpha^{\mathrm c}_j$的分子换成$\exp(\bm i^{\mathrm c}_j)$，指数与除法都逐维进行，因而每一维上全部权重之和为1。最后$\bm h^{\mathrm c}_j=\bm o^{\mathrm c}_j\odot\tanh(\bm c^{\mathrm c}_j)$。若$B_j$为空，字符单元退回普通LSTM的$\bm c^{\mathrm c}_j=\bm f^{\mathrm c}_j\odot\bm c^{\mathrm c}_{j-1}+\bm i^{\mathrm c}_j\odot\tilde{\bm c}^{\mathrm c}_j$。注意有匹配词时式中不再出现$\bm c^{\mathrm c}_{j-1}$：前文信息改由词单元从起点$b$带过来。$\bm h^{\mathrm c}_j$之上接发射层与带类型掩码的CRF，整网用式~\eqref{eq:nlp-analysis-legal-mask}的负对数似然端到端训练。

以“海棠楼201”与教学词典$\{$海棠，海棠楼$\}$为例，匹配词为$w_{1,2}$与$w_{1,3}$，两个词单元都从第1个字的状态出发，分别在“棠”和“楼”处并入。看“楼”处的某一维，设$i^{\mathrm c}_3=0.3$、$i^{\mathrm c}_{1,3}=0.9$，则$\alpha^{\mathrm c}_{1,3}=e^{0.9}/(e^{0.3}+e^{0.9})\approx0.6457$，$\alpha^{\mathrm c}_3\approx0.3543$；若词单元在这一维的记忆为0.8、字符候选为$-0.5$，合并结果为$0.6457\times0.8-0.3543\times0.5\approx0.3394$，主要由“海棠楼”的词信息决定。这个算例也显示出式~\eqref{eq:nlp-extraction-lattice-merge}的一个性质：门值经sigmoid后落在$(0,1)$，再取指数归一化，两个候选的权重只能在$1/(1+e)\approx0.27$与$e/(1+e)\approx0.73$之间，任何一方都不会被完全关掉。

Lattice-LSTM的代价来自结构本身。每个字的输入数随匹配词数变化，词单元又要等起点字算完才能计算，整句只能按字串行推进，难以在批内对齐并行；词典的覆盖与质量也直接决定可用的词信息。

第二种实现针对这一代价。Li等的FLAT（flat-lattice Transformer）保留词格的信息，换掉串行结构\scite{li2020flat}。它把字与匹配词放进同一个序列：$n$个字在前，$m$个匹配词在后，每个单元$u$记录它在原句中的头位置$\mathrm{head}[u]$与尾位置$\mathrm{tail}[u]$；字的头尾相同，词的头尾是首字与末字的位置。“海棠楼201”展平后共8个单元：6个字的头尾为$(1,1),\ldots,(6,6)$，“海棠”为$(1,2)$，“海棠楼”为$(1,3)$。序列中的先后次序不再表示原句位置，位置信息全部由头尾坐标给出。

两个单元之间的关系由四个相对距离刻画：
\begin{equation}
  \begin{aligned}
    d^{(\mathrm{hh})}_{u,v}&=\mathrm{head}[u]-\mathrm{head}[v],&\quad
    d^{(\mathrm{ht})}_{u,v}&=\mathrm{head}[u]-\mathrm{tail}[v],\\
    d^{(\mathrm{th})}_{u,v}&=\mathrm{tail}[u]-\mathrm{head}[v],&\quad
    d^{(\mathrm{tt})}_{u,v}&=\mathrm{tail}[u]-\mathrm{tail}[v].
  \end{aligned}
  \label{eq:nlp-extraction-flat-distance}
\end{equation}
四个数合起来能区分包含、相交与相邻。“海棠楼”与“棠”的四个距离为$(-1,-1,1,1)$，头在其左、尾在其右，表示包含；“海棠”与“楼”为$(-2,-2,-1,-1)$，尾距为$-1$，表示前者紧接在后者之前。每个距离$d$先映射为正弦位置向量$\bm p_d$（$p_d^{(2k)}=\sin(d/10000^{2k/d_{\mathrm{model}}})$，$p_d^{(2k+1)}$取余弦），再融合为相对位置表示$\bm R_{u,v}=\operatorname{ReLU}(\bm W_{\mathrm r}[\bm p_{d^{(\mathrm{hh})}};\bm p_{d^{(\mathrm{th})}};\bm p_{d^{(\mathrm{ht})}};\bm p_{d^{(\mathrm{tt})}}])$。注意力分数采用Transformer-XL式的相对位置写法，设$\bm e_u$为单元嵌入，$\bm q_u=\bm W_{\mathrm q}\bm e_u$，$\bm k_v=\bm W_{\mathrm k}\bm e_v$：
\begin{equation}
  A_{u,v}=\bm q_u^\top\bm k_v+\bm q_u^\top\bm W_{\mathrm R}\bm R_{u,v}+\bm a^\top\bm k_v+\bm c^\top\bm W_{\mathrm R}\bm R_{u,v},
  \label{eq:nlp-extraction-flat-attention}
\end{equation}
四项依次是内容与内容、内容与相对位置、全局内容偏置与全局位置偏置，$\bm a,\bm c$为可学习向量。分数在全部$n+m$个单元上做softmax，其余与标准自注意力相同；自注意力的一般计算可参考\BookChapterRef{chap:transformer}{卷五的Transformer章}。最后只取前$n$个字单元的输出接CRF，词单元只在注意力中提供信息。

与Lattice-LSTM相比，FLAT把“词在哪里、和哪个字是什么关系”从网络拓扑移到了位置编码里。拓扑一旦不再随句子变化，整句可以一次并行计算；每个字也能直接注意到所有匹配词，而不仅是以自己结尾的词。代价是注意力的计算与内存为$O((n+m)^2)$，匹配词多时序列明显变长；两者都依赖词典，词典里没有的新名称得不到词信息。预训练编码器普及后，字级BERT已经从大规模语料中学到不少词边界知识，词汇增强在领域词典质量高、标注语料小的场景中更有价值。

\subsection{基于跨度的实体识别}
\label{sec:nlp-extraction-ner-span}

为表示嵌套实体，跨度路线直接对起止对分类，跨度表示与冲突解码见第\ref{sec:nlp-analysis-paradigm-span}小节。它允许候选相互重叠，代价是$O(n^2)$个候选和平面任务中额外的冲突选择。

\subsubsection{跨度分类与双仿射起止评分}
\label{sec:nlp-extraction-biaffine}

候选集合为$C=\{(i,j):1\leq i\leq j\leq n\}$，每个候选取一个实体类型或“非实体”$\varnothing$。最直接的打分器把式~\eqref{eq:nlp-analysis-span-representation}的跨度表示送入前馈网络。Yu等改用依存分析中的双仿射打分\scite{yu2020ner}：编码器（BERT、词与字符表示加BiLSTM）给出$\bm h_i$后，两个前馈网络分别把它投影为起点表示$\bm p_i=f_{\mathrm s}(\bm h_i)$与终点表示$\bm q_j=f_{\mathrm e}(\bm h_j)\in\mathbb R^{k}$，类型$t$的分数为
\begin{equation}
  z_{i,j,t}=
  \bm p_i^\top\bm U_t\bm q_j+
  \bm w_t^\top[\bm p_i;\bm q_j]+b_t,
  \qquad
  p(t\mid i,j,x)=
  \frac{\exp z_{i,j,t}}{\sum_{t'\in T\cup\{\varnothing\}}\exp z_{i,j,t'}},
  \label{eq:nlp-extraction-biaffine}
\end{equation}
其中$\bm U_t\in\mathbb R^{k\times k}$、$\bm w_t\in\mathbb R^{2k}$、$b_t\in\mathbb R$。双线性项刻画两端是否配对，仿射项保留单独起点和终点的倾向。把全部类型叠在一起，$\bm U$是$k\times(|T|+1)\times k$的张量，一次运算就得到$n\times n\times(|T|+1)$的分数张量，只取上三角$i\leq j$。训练时金标准跨度取其类型，其余候选取$\varnothing$，最小化$L_{\mathrm{span}}=-\sum_{(i,j)\in C}\log p(t^{*}_{i,j}\mid i,j,x)$；$\varnothing$远多于正例，可降采样或降权，但评价必须在完整候选集上进行。

\begin{example}[嵌套候选怎样保留]
\label{ex:nlp-extraction-nested}
教学词元序列为“海棠／大学／计算机／学院／发布／通知”，只用一个实体类别ORG，非实体分数均为0。候选及其ORG分数为$(1,4):4.0$、$(1,2):3.0$、$(2,5):2.0$、$(3,4):1.5$。取$(1,4)$的$\bm p_1=(1,0)^\top$、$\bm q_4=(2,1)^\top$、$\bm U_{\mathrm{ORG}}=\bm I_2$、$\bm w_{\mathrm{ORG}}=(0,0,1,0)^\top$、$b_{\mathrm{ORG}}=0$，式~\eqref{eq:nlp-extraction-biaffine}给出$2+2=4$，实体概率为$\exp(4)/(1+\exp(4))\approx0.9820$，该候选损失约为0.0181。

按分数降序扫描，先保留$(1,4)$。嵌套规范保留其内部的$(1,2)$和$(3,4)$，排除与它部分交叉的$(2,5)$，因为$1<2\leq4<5$；平面规范只保留$(1,4)$。对应字符区间为$[0,9)$、$[0,4)$与$[4,9)$。贪心扫描不保证总分最大，平面规范下的精确解是第\ref{sec:nlp-analysis-paradigm-span}小节的加权区间调度。
\end{example}

与独立的起点、终点指针相比，双仿射给每个起止组合单独评分。若只预测起点集合$\{1,3\}$和终点集合$\{2,4\}$，无法判断$(1,4)$是否成立；双仿射直接对$(1,4)$打分，消除了这种配对歧义。它仍有两点限制：候选为平方级，长文档常需设最大宽度；每个$(i,j)$仍是连续区间，表示不了不连续提及。

\subsubsection{查询式阅读理解NER}
\label{sec:nlp-extraction-mrc-ner}

类型标签“组织”本身携带的信息很少，模型只能从标注样本里猜它的边界。Li等把标注指南改写为查询，例如“找出学校、学院等组织名称”，与正文拼接后送入编码器，把实体识别改成抽取式阅读理解：查询指定找哪类对象，答案是原文中的跨度\scite{li2020mrcner}。每个类型提一次问，不同查询的答案可以重叠，于是不同类型之间的嵌套自然得到保留。

设类型$t$的查询为$c_t$，编码$[\mathrm{CLS}]\,c_t\,[\mathrm{SEP}]\,x\,[\mathrm{SEP}]$后取正文部分的表示$\bm H^{(t)}\in\mathbb R^{n\times d}$。一篇文本可能有多个同类实体，因此模型不是在全文中各选一个起点和终点，而是对每个位置独立做“是否起点”“是否终点”的二分类，两个$d\times2$的分类头给出逐行概率。起点集合与终点集合之间还要配对，匹配头为
\begin{equation}
  p_{i,j}^{(t)}=\sigma\!\left(
  \bm v^{\top}[\bm h^{(t)}_i;\bm h^{(t)}_j]\right),
  \qquad
  L=\alpha L_{\mathrm{s}}+\beta L_{\mathrm{e}}+\gamma L_{\mathrm{match}},
  \label{eq:nlp-extraction-mrc}
\end{equation}
其中$\bm v\in\mathbb R^{2d}$，三个损失分别为起点、终点与配对的二元交叉熵，权重在开发集上选择。没有该类型实体的文本提供全负的起止与配对标签。推理时在预测的起点与终点中枚举$i\leq j$的组合，经匹配头筛选后恢复字符偏移。组织查询应返回“海棠大学”和“海棠大学计算机学院”，地点查询在这句话上返回空集。

查询中的定义和例词让同一个编码器在不同类型间共享知识，适合类型多、需要说明含义的场景。代价是$|T|$个类型要编码$|T|$遍；同一类型内部的嵌套仍取决于匹配头能否正确配对，查询写法也会影响结果，比较时应固定查询模板。

\subsection{嵌套与不连续实体}
\label{sec:nlp-extraction-ner-nested}

不连续实体既不能用BIO表示，也不能用连续跨度表示。“左和右肺均清晰”中，“左……肺”与“右肺”共享中心词“肺”，前者的两段之间隔着“和右”。处理这类结构需要一种输出表示，既能表达“哪些词属于同一个实体”，又不要求这些词相邻。下面的超图与词对网格是两种做法。

\subsubsection{超图表示}
\label{sec:nlp-extraction-hypergraph}

一句话中可能的嵌套实体组合随长度指数增长，逐个枚举不可行。Lu与Roth提出\term{提及超图}（mention hypergraph），用多项式规模的节点与超边共同编码这些组合\scite{lu2015hypergraph}。超图中的一条\term{超边}（hyperedge）从一个父节点指向一组子节点，表示父节点的状态由这组子节点共同展开。下面按其思想给出一种简化写法。对每个位置$k$与类型$t$设五类节点：$A_k$表示“处理位置$k$及其后的所有提及”；$E_k$表示“从$k$开始的提及”；$T_k^t$表示“从$k$开始的类型$t$提及”；$I_k^t$表示“位置$k$在某个类型$t$提及内部”；终止节点$X$表示“到此结束”。超边规定节点可以怎样展开：$A_k\to(E_k,A_{k+1})$；$E_k\to(T_k^{1},\ldots,T_k^{|T|})$；$T_k^t\to X$表示没有类型$t$的提及从$k$开始，$T_k^t\to I_k^t$表示有；$I_k^t$有三种展开，$\to I_{k+1}^t$表示提及延续，$\to X$表示提及在$k$结束，$\to(I_{k+1}^t,X)$表示一个提及在$k$结束、另一个更长的同时延续。

一个输入的标注结构是从$A_1$出发、每个访问到的节点各选一条超边形成的子超图，提及就是从某个$T_k^t$经若干$I^t$节点到达$X$的路径。以“海棠／大学／计算机／学院”的三个ORG提及$(1,4)$、$(1,2)$、$(3,4)$为例：$T_1\to I_1$，$I_1\to I_2$，$I_2\to(I_3,X)$表示$(1,2)$在第2个词结束而$(1,4)$继续，$T_3\to I_3$，$I_3\to I_4$，$I_4\to X$，其余$T_k\to X$。从$T_1$出发有两条到$X$的路径，读出$(1,2)$与$(1,4)$；从$T_3$出发读出$(3,4)$。$I_3$被两个提及共享，这正是超图紧凑的来源。

模型是超图上的对数线性模型。每条超边$\varepsilon$有特征$\bm f(\varepsilon,x)$，结构$y$的得分是所选超边得分之和，条件概率按全部合法子超图归一化。配分函数用超图上的inside递推计算：
\begin{equation}
  \begin{aligned}
    Z(X)&=1,\qquad
    Z(v)=\sum_{\varepsilon\in\mathrm{out}(v)}\exp\bigl(\bm w^\top\bm f(\varepsilon,x)\bigr)\prod_{u\in\mathrm{child}(\varepsilon)}Z(u),\\
    \log p(y\mid x)&=\sum_{\varepsilon\in y}\bm w^\top\bm f(\varepsilon,x)-\log Z(A_1),
  \end{aligned}
  \label{eq:nlp-extraction-hypergraph-inside}
\end{equation}
$\mathrm{out}(v)$为从$v$出发的超边，$\mathrm{child}(\varepsilon)$为其子节点。节点按位置从右向左计算，每个位置有$O(|T|)$个节点和常数条超边，总代价为$O(n|T|)$，与句长成线性；把求和换成取最大即得解码。这与第\ref{sec:nlp-analysis-paradigm-graph}小节的区间动态规划同属一类，只是递推定义在超图而非区间表上。后续工作把超边特征换成LSTM表示，训练目标不变。

线性复杂度的代价是\emph{结构歧义}：一个子超图可能对应不止一个提及集合。设同类型提及为$(1,3)$与$(2,4)$，超图需要$T_1\to I_1$、$I_1\to I_2$、$T_2\to I_2$、$I_2\to I_3$、$I_3\to(I_4,X)$、$I_4\to X$，从中可读出$(1,3)$、$(1,4)$、$(2,3)$、$(2,4)$四条路径，与“恰有$(1,3)$和$(2,4)$”无法区分。解码时只能约定一种读法，例如输出全部路径，这在同类型交叉提及上会产生多余实体。原始超图只表示连续提及；Muis与Lu在超图中为提及内部的词加入“属于实体”与“间隙”两种状态，从而能表示部分不连续实体，但歧义也随之增多\scite{muis2016discontiguous}。

\subsubsection{W2NER词元关系图}
\label{sec:nlp-extraction-w2ner}

超图把结构编码在节点状态里，需要为不同结构精心设计状态。W2NER换一个角度：既然不连续实体的困难在于“哪些词属于同一实体”，就直接预测词与词之间的关系\scite{li2022w2ner}。它定义两类有向关系：NNW（next-neighboring-word）从实体中的一个词指向该实体中次序上的下一个词，这里的“下一个”指实体内部的次序，不要求原文相邻；THW-$t$（tail-head-word）从实体的尾词指回首词，并携带类型$t$。对$n$个词，输出是$n\times n$的关系网格$\bm G$，每格取$\{\mathrm{NONE},\mathrm{NNW}\}\cup\{\mathrm{THW}\text{-}t:t\in T\}$之一。由于实体内的词按原文顺序排列，NNW总在上三角$i<j$，THW总在下三角或对角线$i\geq j$，单词实体就是对角线上的THW。

“左／和／右／肺／均／清晰”中，“左……肺”由第1、4个词组成，“右肺”由第3、4个词组成。正标签为$G_{1,4}=G_{3,4}=\mathrm{NNW}$、$G_{4,1}=G_{4,3}=\mathrm{THW}\text{-}\mathrm{BODY}$，其余为NONE，如图~\ref{fig:nlp-extraction-w2ner-grid}。两个实体共享尾词“肺”，在网格中分别占用不同的格，不会互相覆盖；平面与嵌套实体同样只是网格中的若干NNW路径和THW格。

\begin{figure}[htbp]
  \centering
  % 第九卷 10.1.4.2 节：W2NER词对关系网格（教学句“左／和／右／肺／均／清晰”）。
% 行为i、列为j；上三角(j>i)放NNW，下三角(i>j)放THW-类型，对角线可放单词实体的THW。
% 正标签：G(1,4)=G(3,4)=NNW，G(4,1)=G(4,3)=THW-BODY，其余为NONE。
\begin{tikzpicture}[
  x=1mm,y=1mm,figure picture,
  cellbox/.style={draw=FigureGrid,line width=.4pt,minimum size=8mm,inner sep=0pt},
  upper/.style={cellbox,fill=white},
  lower/.style={cellbox,fill=FigurePaper},
  nnw/.style={cellbox,fill=FigureModelFill,draw=FigureModel,line width=.7pt,
    font=\figurefont\scriptsize},
  thw/.style={cellbox,fill=FigureInputFill,draw=FigureInput,line width=.7pt,
    font=\figurefont\tiny,align=center},
  lab/.style={font=\figurefont\footnotesize,text=FigureInk},
  note/.style={font=\figurefont\scriptsize,text=FigureMuted,anchor=west,align=left}
]
  % 网格底色
  \foreach \i in {1,...,6}{
    \foreach \j in {1,...,6}{
      \pgfmathtruncatemacro{\up}{\j>\i ? 1 : 0}
      \pgfmathtruncatemacro{\dn}{\i>\j ? 1 : 0}
      \ifnum\up=1
        \node[upper] at ({(\j-1)*8},{-(\i-1)*8}) {};
      \else\ifnum\dn=1
        \node[lower] at ({(\j-1)*8},{-(\i-1)*8}) {};
      \else
        \node[cellbox,fill=white] at ({(\j-1)*8},{-(\i-1)*8}) {};
      \fi\fi
    }
  }
  % 行列标签
  \foreach \k/\w in {1/左,2/和,3/右,4/肺,5/均,6/清晰}{
    \node[lab] at ({(\k-1)*8},7) {\w};
    \node[font=\figurefont\scriptsize,text=FigureMuted] at ({(\k-1)*8},11.5) {$j{=}\k$};
    \node[lab,anchor=east] at (-6,{-(\k-1)*8}) {\w};
    \node[font=\figurefont\scriptsize,text=FigureMuted,anchor=east] at (-13,{-(\k-1)*8}) {$i{=}\k$};
  }
  % 正标签
  \node[nnw] at ({(4-1)*8},{-(1-1)*8}) {NNW};
  \node[nnw] at ({(4-1)*8},{-(3-1)*8}) {NNW};
  \node[thw] at ({(1-1)*8},{-(4-1)*8}) {THW\\BODY};
  \node[thw] at ({(3-1)*8},{-(4-1)*8}) {THW\\BODY};
  % 区域说明
  \node[note] at (50,0) {上三角（$j>i$）：NNW，\\实体内部的次序相邻词};
  \node[note] at (50,-14) {下三角（$i>j$）：THW-类型，\\由尾词指回首词};
  \node[note] at (50,-28) {解码：THW(4,1)与路径1$\to$4\\得“左…肺”$[0,1)\cup[3,4)$；};
  \node[note] at (50,-38) {THW(4,3)与路径3$\to$4\\得“右肺”$[2,4)$};
\end{tikzpicture}

  \caption{W2NER在“左和右肺均清晰”上的词对关系网格。第$i$行第$j$列表示从词$i$到词$j$的关系，未标注的格为NONE。上三角的NNW连接同一实体内次序相邻的词，下三角的THW-BODY由尾词指回首词并给出类型。两条NNW路径$1\to4$与$3\to4$共享尾词“肺”，分别由THW$(4,1)$与THW$(4,3)$确认。标注规范为教学设定。}
  \label{fig:nlp-extraction-w2ner-grid}
\end{figure}

网格中每一格需要一个词对表示，W2NER的逐层运算如下。
\begin{enumerate}
  \item 编码：BERT输出经首子词对齐与BiLSTM，得到$\bm H\in\mathbb R^{n\times d_{\mathrm h}}$。
  \item \term{条件层归一化}（conditional layer normalization, CLN）：用词$i$生成归一化的缩放与平移参数，作用在词$j$上，
  \[
    \bm V_{i,j}=\bm\gamma_{i}\odot\frac{\bm h_j-\mu_j}{\sigma_j}+\bm\lambda_{i},\qquad
    \bm\gamma_i=\bm W_{\gamma}\bm h_i+\bm b_{\gamma},\quad
    \bm\lambda_i=\bm W_{\lambda}\bm h_i+\bm b_{\lambda},
  \]
  $\mu_j,\sigma_j$是$\bm h_j$各维的均值与标准差。普通层归一化的缩放与平移是全局参数，这里由另一个词决定，于是$\bm V_{i,j}$表示“以$i$为条件看$j$”，得到$\bm V\in\mathbb R^{n\times n\times d_{\mathrm h}}$。
  \item 距离与区域嵌入：把$|i-j|$分桶后查表得距离嵌入$\bm E^{\mathrm d}_{i,j}$，再用区域嵌入$\bm E^{\mathrm r}_{i,j}$区分上三角、下三角与对角线；三者拼接后经多层感知机降维，得到$\bm C\in\mathbb R^{n\times n\times d_{\mathrm c}}$。区域嵌入让模型知道某格只可能是NNW还是THW。
  \item 多粒度空洞卷积：在网格上做空洞率$l\in\{1,2,3\}$的二维卷积，$\bm Q^{(l)}=\operatorname{GELU}(\operatorname{DConv}_l(\bm C))$，拼接得$\bm Q\in\mathbb R^{n\times n\times3d_{\mathrm c}}$。\term{空洞卷积}（dilated convolution）在卷积核元素之间留出间隔，不增加参数就扩大感受野；在这里它让一格能看到同行、同列附近的格，例如判断$G_{4,1}$时参考$G_{1,4}$是否为NNW。
  \item 联合预测：双仿射分支以编码器输出为输入，$\bm s_i=\operatorname{MLP}_{\mathrm s}(\bm h_i)$、$\bm o_j=\operatorname{MLP}_{\mathrm o}(\bm h_j)$，$\bm y'_{i,j}=\bm s_i^\top\bm U\bm o_j+\bm W[\bm s_i;\bm o_j]+\bm b$；多层感知机分支以网格表示为输入，$\bm y''_{i,j}=\operatorname{MLP}(\bm Q_{i,j})$。两者相加后逐格softmax，$p(\cdot\mid i,j)=\operatorname{softmax}(\bm y'_{i,j}+\bm y''_{i,j})\in\mathbb R^{|T|+2}$。
\end{enumerate}
训练对全部$n^2$个格做交叉熵，NONE格也参与。双仿射分支直接看两个词本身，网格分支看邻近格的关系模式，两者互补。

解码把网格还原为实体，步骤见算法~\ref{alg:nlp-extraction-w2ner-decode}。NNW边只从小下标指向大下标，构成有向无环图，从首词到尾词的路径可以用深度优先搜索枚举。

\Needspace{10\baselineskip}
\begin{algorithm}[W2NER网格解码]
  \label{alg:nlp-extraction-w2ner-decode}
  \begin{algorithmic}[1]
    \Require 预测网格$\hat{\bm G}\in\mathcal R^{n\times n}$，$\mathcal R=\{\mathrm{NONE},\mathrm{NNW}\}\cup\{\mathrm{THW}\text{-}t\}$；词到字符的偏移表
    \Ensure 实体集合$\hat{\mathcal E}$，每个实体为片段列表加类型
    \State 建有向图$\mathcal D$：对每个$i<j$且$\hat G_{i,j}=\mathrm{NNW}$，加边$i\to j$
    \State $\hat{\mathcal E}\gets\varnothing$
    \For{每个$i\geq j$且$\hat G_{i,j}=\mathrm{THW}\text{-}t$} \Comment{$i$为尾词，$j$为首词}
      \If{$i=j$}
        \State 加入单词实体$(\{j\},t)$
      \Else
        \For{$\mathcal D$中每条从$j$到$i$的路径$j=k_1\to k_2\to\cdots\to k_r=i$}
          \State 加入实体$(\{k_1,\ldots,k_r\},t)$
        \EndFor
      \EndIf
    \EndFor
    \State 把每个实体中下标连续的词合并为片段，按偏移表转为字符区间
  \end{algorithmic}
\end{algorithm}

在教学例子中，THW$(4,1)$要求从1到4的路径，图中只有边$1\to4$，得到$\{1,4\}$，两段字符区间为$[0,1)$与$[3,4)$；THW$(4,3)$得到路径$3\to4$，合并为一段$[2,4)$。“左和右肺”整体不会被当作一个实体，因为图中没有$1\to2$或$2\to3$的边。

W2NER用一套标签统一了三种结构，解码也只是图上找路径。它的弱点在两处。其一，网格大小为$n^2$，长文本的计算与内存压力明显，通常按句处理。其二，路径组合可能产生标注中不存在的实体：若两个同类型实体的NNW边交错，THW确认的首尾之间可能有多条路径，算法会把它们都输出。不同语料对省略形式是否算不连续实体的规定不一，比较结果前需要统一标注规范。

\subsection{基于生成的实体识别}
\label{sec:nlp-extraction-ner-generation}

前几小节的每种方法都为特定结构设计了专用标签和解码器：BIO管平面，跨度管嵌套，网格管不连续。生成路线把实体序列化为目标序列，一个序列到序列模型同时覆盖三种结构，并能直接利用预训练的编码器—解码器；线性化与约束解码的通用讨论见第\ref{sec:nlp-analysis-paradigm-generation}小节。

\subsubsection{BARTNER指针式生成}
\label{sec:nlp-extraction-bartner}

若让模型直接生成实体字符串，同一字符串在原文中出现多次时无法确定位置，模型还可能写出原文没有的字。Yan等的BARTNER让目标序列只含两种符号：指向输入位置的指针和类型标记\scite{yan2021bartner}。每个实体写成其各片段起止位置的指针序列，后接一个类型标记；多个实体按原文位置依次排列。“左和右肺均清晰”的目标为
\[
  \underbrace{1\ 1\ 4\ 4\ \langle\mathrm{BODY}\rangle}_{\text{左……肺}}\quad
  \underbrace{3\ 4\ \langle\mathrm{BODY}\rangle}_{\text{右肺}}\quad\langle\mathrm{eos}\rangle,
\]
第一个实体有两段$[1,1]$与$[4,4]$，故有四个指针。连续实体只有一对起止指针，嵌套实体就是指针范围相互包含的两个实体，三种结构用同一种序列表示。

解码器每步的输出空间是$n$个输入位置与$|C|$个类型标记（含结束符）的并集。设BART编码器输出$\bm H^{\mathrm e}\in\mathbb R^{n\times d}$，输入词元嵌入为$\bm E^{\mathrm e}\in\mathbb R^{n\times d}$，类型标记的嵌入为$\bm G\in\mathbb R^{|C|\times d}$。已生成的目标符号先换回词元：指针$y_u\leq n$换成输入的第$y_u$个词元，类型标记换成对应的类型词，这样解码器的输入仍是普通词元。第$u$步的解码器状态为$\bm g_u\in\mathbb R^{d}$，输出分布为
\begin{equation}
  \begin{aligned}
    \bar{\bm H}^{\mathrm e}&=\kappa\operatorname{MLP}(\bm H^{\mathrm e})+(1-\kappa)\bm E^{\mathrm e},\\
    p_\theta(y_u\mid\bm y_{<u},x)&=\operatorname{softmax}\bigl([\bar{\bm H}^{\mathrm e}\bm g_u;\ \bm G\bm g_u]\bigr)\in\mathbb R^{n+|C|},
  \end{aligned}
  \label{eq:nlp-extraction-bartner}
\end{equation}
前$n$维是指向各输入位置的概率，后$|C|$维是各类型标记的概率，$\kappa\in[0,1]$为超参数，混合上下文表示与原始词嵌入。这是第\ref{sec:nlp-analysis-paradigm-generation}小节“只含位置指针和类型标记”的彻底复制形式：输出不经过词表，因此不会生成原文没有的字。

训练使用教师强制，即每步以正确前缀为条件，最小化目标序列$\bm y^*$的负对数似然：
\begin{equation}
  L_{\mathrm{gen}}=
  -\sum_{u=1}^{|\bm y^{*}|}
  \log p_\theta(y_u^{*}\mid \bm y_{<u}^{*},x,c),
  \label{eq:nlp-extraction-generation}
\end{equation}
$c$为可选的任务条件，BARTNER中为空，下一小节的UIE中为模式提示。推理用束搜索逐步生成，再按规则解析：读取指针直到遇到类型标记，把指针两两配成片段。指针个数为奇数、片段起点大于终点或后一段起点不大于前一段终点时，该实体不合法。这些规则也可以写成解码时的掩码，例如类型标记之前必须有偶数个指针。

指针序列的位置信息可以直接还原；不足在于实体之间本无顺序，按原文位置排列是额外约定，模型生成正确集合的另一种排列也会受交叉熵惩罚，逐步生成也慢于网格或跨度方法的一次并行预测。

\subsubsection{UIE结构化生成与模式提示}
\label{sec:nlp-extraction-uie}

BARTNER只处理实体。若实体、关系和事件都要用同一个模型抽取，就需要一种能表达类型、片段及其关联的统一输出语言，以及一种告诉模型“这次要抽什么”的输入方式。Lu等的UIE为此设计了两部分\scite{lu2022uie}。

\term{结构化抽取语言}（structured extraction language, SEL）把抽取结果写成括号表达式。一条记录由若干“定位”组成，每个定位写出类型名与原文片段，下面可挂若干“关联”，每个关联写出角色名与片段：
\[
  \begin{aligned}
    \langle\text{记录}\rangle&::=\texttt{(}\ \{\langle\text{定位}\rangle\}\ \texttt{)},\\
    \langle\text{定位}\rangle&::=\texttt{(}\ \text{类型名}\ \texttt{:}\ \text{片段}\ \{\langle\text{关联}\rangle\}\ \texttt{)},\\
    \langle\text{关联}\rangle&::=\texttt{(}\ \text{角色名}\ \texttt{:}\ \text{片段}\ \texttt{)},
  \end{aligned}
\]
花括号表示可重复零次或多次。
实体只用定位，关系是带一个关联的定位，事件是触发词定位下挂多个论元关联。\term{结构模式提示}（structural schema instructor, SSI）把本次要识别的类型和关联名放在输入前面，用特殊标记\code{[spot]}引出类型名、\code{[asso]}引出关联名、\code{[text]}引出原文。对通知抽取地点与组织时，输入为“\code{[spot]} 地点 \code{[spot]} 组织 \code{[text]} 机器学习读书会原定于……地点为海棠楼201。”，目标为\code{((地点: 海棠楼201))}。模式提示换成关系或事件的名称，同一模型就转为关系或事件抽取。

UIE以T5为骨架，编码器读入$[c;x]$，其中$c$为SSI，解码器生成SEL。下游微调使用式~\eqref{eq:nlp-extraction-generation}的教师强制损失。为了让模型在没见过的模式上也会用这套语言，UIE先在三类数据上预训练：文本与结构配对的数据$\mathcal D_{\mathrm{pair}}$，来自知识库与百科文本的远程对齐；只有结构的数据$\mathcal D_{\mathrm{record}}$；以及无标注文本$\mathcal D_{\mathrm{text}}$。三项损失相加：
\begin{equation}
  \begin{aligned}
    L_{\mathrm{pre}}&=L_{\mathrm{pair}}+L_{\mathrm{record}}+L_{\mathrm{text}},\qquad
    L_{\mathrm{pair}}=-\sum_{(x,y)\in\mathcal D_{\mathrm{pair}}}\log p_\theta(y\mid x,c_{\mathrm{meta}}),\\
    L_{\mathrm{record}}&=-\sum_{y\in\mathcal D_{\mathrm{record}}}\sum_{u}\log p_\theta(y_u\mid y_{<u}),\qquad
    L_{\mathrm{text}}=-\sum_{x\in\mathcal D_{\mathrm{text}}}\log p_\theta(x''\mid x').
  \end{aligned}
  \label{eq:nlp-extraction-uie-pretrain}
\end{equation}
$L_{\mathrm{pair}}$教模型从文本到结构的映射，其提示$c_{\mathrm{meta}}$除了样本中真实出现的类型，还随机加入未出现的负类型，让模型学会对提示中的类型说“没有”；$L_{\mathrm{record}}$只训练解码器，把它当作结构语言的语言模型，学会生成合法的SEL；$L_{\mathrm{text}}$是T5式的跨度损坏去噪（第\ref{sec:nlp-llm-objective-t5-bart}小节），$x'$为遮去若干片段的文本，$x''$为被遮去的片段，用来保持通用语言能力。微调时UIE还在目标中随机插入\code{(类型: [NULL])}一类负定位，进一步训练拒绝输出。

模式提示只是条件，不能保证输出合法。推理时可以维护SEL的文法状态并施加约束：类型名位置只允许SSI中列出的名称，片段位置只允许原文的子串（用原文的前缀树检查），括号必须配对；机制见第\ref{sec:nlp-inference-grammar}小节。生成结束后还要把片段对齐回原文：若“海棠楼201”出现两次，字符串本身不能确定是哪一次，需要按上下文选择或保留全部候选位置。

与BARTNER的位置指针相比，UIE输出文字片段，换来跨任务的统一接口，代价正是这一步原文对齐。

\subsection{少样本与提示式实体识别}
\label{sec:nlp-extraction-ner-fewshot}

新类型往往只有几个标注样本，前述监督模型无法直接训练。度量学习路线用少量样本构造类别表示，对查询词元按距离分类，仍需在源领域上训练表示；大语言模型提示路线把类型定义和示例写进输入，不更新参数，但要处理格式解析、原文对齐与输出不稳定。

\subsubsection{原型网络与近邻标注}
\label{sec:nlp-extraction-prototype}

少样本实体识别通常按\emph{$N$类$K$样本}设定评价：每个测试任务给出$N$个新类型，每类只有约$K$个标注实体，组成\term{支持集}（support set）$\mathcal S$；模型要标注另一批句子组成的\term{查询集}（query set）。Few-NERD按句子采样任务，保证每类实体在支持集中出现$K$到$2K$次，并区分两种划分：训练与测试的粗类不同，或粗类相同而细类不同\scite{ding2021fewnerd}。词元级的标签空间是$N$个类型再加O。

\term{原型网络}（prototypical network）用每类支持样本的均值作为该类的代表\scite{ket-snell2017-prototypical}。设编码器把词元$x$映射为$f_\theta(x)\in\mathbb R^{d}$，$\mathcal S_t$为支持集中标签为$t$的全部词元，类$t$的原型和查询词元的类别分布为
\begin{equation}
  \bm c_t=\frac{1}{|\mathcal S_t|}\sum_{x'\in\mathcal S_t}f_\theta(x'),\qquad
  p(y=t\mid x)=\frac{\exp\bigl(-\lVert f_\theta(x)-\bm c_t\rVert_2^2\bigr)}{\sum_{t'}\exp\bigl(-\lVert f_\theta(x)-\bm c_{t'}\rVert_2^2\bigr)}.
  \label{eq:nlp-extraction-prototype}
\end{equation}
训练采用\term{元训练}（meta-training）：在源领域的已知类型上反复采样与测试同形的$N$类$K$样本任务，以查询词元的交叉熵更新$\theta$，使编码器学会把同类词元聚在一起；测试时冻结$\theta$，只用新类型的支持集计算原型。原型网络的一般讨论可参考\BookChapterRef{chap:ket-transfer-adaptation}{卷八的知识迁移章}。

用于实体识别时，O类暴露出原型的弱点。O包含标点、动词、普通名词等各种词元，表示空间里是多个分散的簇，均值可能落在一片空白区域。设二维教学表示中O的两个支持词元为$(0,1)$与$(0,-1)$，原型为$(0,0)$；LOC的支持词元为$(1,0)$，原型就是它本身。一个LOC查询词元位于$(0.45,0)$，到O原型的平方距离为$0.2025$，到LOC原型为$0.3025$，于是被判为O。若改看最近的单个支持词元，它到两个O词元的平方距离都是$1.2025$，到LOC词元为$0.3025$，判为LOC。

Yang与Katiyar据此提出NNShot：先在源领域上按普通监督方式训练一个实体识别模型，取其表示作为$f$，测试时把查询词元归到最近的支持词元所属的类，$\hat y(x)=\argmin_t\min_{x'\in\mathcal S_t}\lVert f(x)-f(x')\rVert_2^2$\scite{yang2020structshot}。它不需要元训练，近邻规则也不假设每类只有一个簇。逐词元分类仍可能输出不合法或不连贯的标签，StructShot再加入标签转移：新领域的类型与源领域不同，具体类型间的转移无法估计，于是把标签抽象为O、“与前一个同类型的I”和“与前一个不同类型的I”等少数几类，在源领域上估计这些\emph{抽象转移}的概率，再展开到新类型上。解码用Viterbi求
\begin{equation}
  \hat{\bm y}=\argmax_{\bm y}\sum_{i=1}^{n}\Bigl[\log p_{\mathrm{NN}}(y_i\mid x_i)+\lambda\log p_{\mathrm{abs}}(y_i\mid y_{i-1})\Bigr],
  \label{eq:nlp-extraction-structshot}
\end{equation}
$p_{\mathrm{NN}}$由各类最近支持词元的距离经softmax得到。原文把发射概率与转移概率相乘，由于两者分别估计、尺度不一致，再用温度超参数重新归一化转移概率；式中的$\lambda$是本书对这一尺度调节的对数形式写法，不是原文的参数。它与第\ref{sec:nlp-extraction-bio}小节的CRF形式相同，区别在于发射来自近邻距离、转移来自源领域的抽象统计，两者都不在新类型上训练。

三者的取舍在于怎样用几个样本定义类别：原型网络适合类内紧凑的情形；NNShot适应多峰的O类，但对支持集中的个别错误标注更敏感；StructShot补回标签依赖，前提是源领域与目标领域的标签转移模式相近。

\subsubsection{大语言模型提示式抽取}
\label{sec:nlp-extraction-llm-ner}

大语言模型不需要元训练，也不需要支持集上的编码器。上下文学习（第\ref{sec:nlp-inference-icl-shots}小节）把任务说明与少量示例写进提示，模型在不更新参数的情况下给出输出。提示式实体识别的提示$P$由三部分组成：指令$\mathcal I$写明类型定义、边界规则和“没有就返回空”等要求；$k$个示例$(x^{(1)},y^{(1)}),\ldots,(x^{(k)},y^{(k)})$展示输入与输出的对应；最后是待标注文本$x$与输出格式说明。模型按$\hat y\sim p_{\mathrm{LM}}(\cdot\mid P(\mathcal I,\{(x^{(r)},y^{(r)})\},x))$生成，再由解析程序把$\hat y$转成实体集合。

示例与输出格式是两个主要设计点。示例数受上下文长度限制，常按与$x$的表示相似度检索最近的$k$个训练句，让示例覆盖相近的边界与类型模式。输出格式有两类。一类是结构化列表，例如JSON数组\code{[\{"text": "海棠楼201", "type": "地点"\}]}，便于解析，可用JSON Schema约束解码保证格式合法（第\ref{sec:nlp-inference-grammar}小节）；但它只给字符串，需要在原文中查找位置，同一字符串出现多次时还要消歧。另一类是标记改写，让模型把原文抄写一遍并在实体两侧加记号，例如GPT-NER用\code{@@海棠楼201\#\#}标出地点，每次只问一种类型\scite{wang2023gptner}；输出与原文逐字对齐即可恢复偏移，代价是要抄写全文，长文本上成本高，抄写偏差也会导致对齐失败。

单次生成的结果不稳定：同一提示多次采样，可能漏掉或多出某个实体，边界也可能漂移。可以借用自一致性（第\ref{sec:nlp-execution-self-consistency}小节）的思路做实体级投票。采样$R$次得到实体集合$\hat{\mathcal E}^{(1)},\ldots,\hat{\mathcal E}^{(R)}$，对每个候选“片段加类型”$e$计票，保留票数达到阈值$\kappa$的实体：
\begin{equation}
  \operatorname{vote}(e)=\sum_{r=1}^{R}\mathbb 1\bigl[e\in\hat{\mathcal E}^{(r)}\bigr],\qquad
  \hat{\mathcal E}=\{e:\operatorname{vote}(e)\geq\kappa\}.
  \label{eq:nlp-extraction-entity-vote}
\end{equation}
与第\ref{sec:nlp-execution-self-consistency}小节对单个答案投票不同，这里每个实体单独投票，输出仍是集合。教学示例：$R=5$，“海棠楼201／地点”出现4次，边界不全的“海棠楼／地点”出现1次，“机器学习读书会／组织”出现2次；取$\kappa=3$时只保留第一项。读书会是否算组织取决于指令中的类型定义，票数分散正反映了模型对定义的不确定，补充定义比提高采样次数更直接。另一种做法是GPT-NER采用的自我核验：对每个候选实体再问一次“原文中的‘海棠楼201’是不是地点”，只保留肯定回答，用来抑制模型把非实体标成实体的倾向。无论哪种做法，输出不合格式、片段不是原文子串、类型不在指令列出的集合中，都应计为错误，不能静默修补；原文中没有的“李老师”即使格式合法，也是无依据的新增。

提示式抽取几乎不需要标注，换类型只需改指令，适合冷启动与类型频繁变化的场景；代价是长提示与多次采样的成本，严格边界上的准确率通常也低于在充足标注上训练的专用模型。两者可以衔接：用提示式抽取产生初始标注，经人工校对后训练第\ref{sec:nlp-extraction-ner-span}小节的跨度模型；或者把多种抽取任务整理为指令数据，对模型做指令微调（第\ref{chap:nlp-supervised-finetuning}章），使它在抽取格式与边界上更稳定。

\section{实体链接}
\label{sec:nlp-extraction-linking}

识别出提及以后，还要确定它对应知识库中的哪个对象。\term{实体链接}（entity linking, EL）把提及映射到给定版本知识库中的条目；没有可确认的对应项时输出NIL。为了比较本节的四种方法，构造两个同名“海棠大学”：$e_1$是教学大学，描述含“计算机学院”；$e_2$是教学小说中的虚构学校，描述含“故事人物”。“海棠大学计算机学院发布通知”应链接$e_1$，“小说中的海棠大学迎来新角色”应链接$e_2$。

实体链接有四项挑战。别名多：同一对象有全称、简称和俗称；同名多：一个名字对应多个对象；知识库大：候选对象可达数百万，不能对每个都做精细比较；库外实体：提及的对象可能尚未登记。表~\ref{tab:nlp-extraction-linking-methods}按这些挑战对照本节的四种方法：别名字典与先验、BLINK召回与重排、文档全局一致性，以及GENRE生成式名称预测。

\begin{table*}[htbp]
  \centering
  \small
  \begin{tabularx}{\linewidth}{@{}>{\raggedright\arraybackslash}p{22mm}>{\raggedright\arraybackslash}p{30mm}>{\raggedright\arraybackslash}X>{\raggedright\arraybackslash}p{30mm}>{\raggedright\arraybackslash}p{34mm}@{}}
    \toprule
    方法 & 候选来源 & 使用的证据 & 新增实体 & 主要代价\\
    \midrule
    别名字典与先验 & 登记名称与锚文本计数 & 链接先验、上下文相似、类型兼容 & 需补别名与计数 & 漏别名则无法召回\\
    BLINK & 实体向量近邻检索 & 提及上下文与实体描述的语义匹配 & 写好描述即可编码入库 & 索引维护；重排需逐候选联合编码\\
    全局一致性 & 前一阶段的候选 & 局部分数加实体间兼容 & 依赖实体表示 & 联合推断需近似；跨领域文档易误伤\\
    GENRE & 名称前缀树 & 上下文条件下的名称生成概率 & 加入名称即合法 & 名称须唯一；串行解码\\
    \bottomrule
  \end{tabularx}
  \caption{实体链接的方法谱系。前两种方法常组成“召回—重排”流水线，全局一致性作用在重排之后，GENRE以一个生成模型同时完成召回与选择。}
  \label{tab:nlp-extraction-linking-methods}
\end{table*}

评价要分两段看。候选召回问正确身份是否进入候选集合$C(m)$，常用前$k$召回率；重排问在集合中是否选对。正确身份一旦在召回阶段遗漏，后续再准确的重排也无法恢复。评价还须声明知识库快照、提及是否给定、候选规模和NIL比例。

\subsection{别名字典与实体先验排序}
\label{sec:nlp-extraction-alias}

最直接的候选来源是登记名称、人工别名和历史链接记录。百科文本中的锚文本是天然的链接标注：若“海棠大学”这段锚文本有$c(m,e)$次指向条目$e$，则\term{链接先验}（link prior）为$p(e\mid m)=c(m,e)/\sum_{e'}c(m,e')$，它表示不看语境时这个名字指向$e$的频率。别名表把“海棠大学”映射到$\{e_1,e_2\}$，再加入上下文与类型信息组成线性评分：
\begin{equation}
  s(m,e)=\log p(e\mid m)+
  \lambda\,\operatorname{sim}(x_m,d_e)+
  \eta\,\operatorname{type}(m,e),
  \label{eq:nlp-extraction-alias}
\end{equation}
其中$x_m$为提及上下文，$d_e$为实体描述，$\operatorname{sim}$为二者的相似度（如词袋余弦），$\operatorname{type}$为提及类型与实体类型的兼容分数，$\lambda,\eta$在开发集上选择。先验只能从训练期的计数估计，罕见别名需要平滑。

若教学先验为$(0.8,0.2)$，小说语境的相似度为$(0,2)$，取$\lambda=1$且类型项相等，两候选的分数为$\log0.8\approx-0.2231$与$\log0.2+2\approx0.3906$，选$e_2$，语境推翻了先验的首选；在学院语境中相似度换成$(2,0)$，则选$e_1$。这个方法查表即可、速度很快，先验本身也是强基线；但候选质量完全取决于别名覆盖，若别名表漏掉$e_2$，任何语境都选不到它。低分拒绝为NIL的阈值需要在含库外提及的开发集上校准。

\subsection{BLINK双编码召回与交叉编码重排}
\label{sec:nlp-extraction-blink}

别名表漏候选、只用浅层语境，两项不足都来自“按名字找”。Wu等的BLINK改为“按意思找”：用语义向量在全库中检索候选，再用更精细的模型重排\scite{wu2020blink}。它的设计目标是\emph{零样本链接}：测试时的实体在训练中从未出现，只有一段文字描述。

第一阶段是\term{双编码器}（bi-encoder）。提及一侧的输入是带边界标记的上下文“$[\mathrm{CLS}]$ 左文 $[\mathrm{M_s}]$ 提及 $[\mathrm{M_e}]$ 右文 $[\mathrm{SEP}]$”，实体一侧的输入是“$[\mathrm{CLS}]$ 名称 $[\mathrm{ENT}]$ 描述 $[\mathrm{SEP}]$”，两个独立的BERT编码器各取$[\mathrm{CLS}]$位置的输出，得到$\bm u_m,\bm v_e\in\mathbb R^{d}$，相似分数为内积$s_{\mathrm{bi}}(m,e)=\bm u_m^\top\bm v_e$。训练使用第\ref{chap:nlp-classification-matching-organization}章式~\eqref{eq:nlp-matching-infonce}的对比损失，温度取1，对正确实体$e^+$与负实体集合$N_m$，
\begin{equation}
  L_{\mathrm{bi}}=
  -\log\frac{\exp s_{\mathrm{bi}}(m,e^+)}
  {\sum_{e\in\{e^+\}\cup N_m}\exp s_{\mathrm{bi}}(m,e)}.
  \label{eq:nlp-extraction-blink}
\end{equation}
负例先取同一批中其他提及的正确实体，这样一批$B$个提及只需编码$B$个实体，就得到$B\times B$的分数矩阵；第二轮再用训练好的双编码器检索出分数高的错误实体作为\term{难负例}（hard negative）继续训练。两侧编码相互独立，因此全部实体向量可以离线算好、建近邻索引，查询时只编码提及一次。

第二阶段是\term{交叉编码器}（cross-encoder）。它把提及上下文与候选实体的名称和描述拼成一个输入，送入同一个BERT，取$[\mathrm{CLS}]$输出$\bm g_{m,e}\in\mathbb R^{d}$，打分$s_{\mathrm{cross}}(m,e)=\bm w^{\top}\bm g_{m,e}$。训练时在双编码器召回的前$k$个候选上做softmax交叉熵。联合编码让提及中的每个词元都能注意到描述中的每个词元，比较更细致，但每个候选都要重新编码一遍，只能用于少量候选。

学院语境的教学计算：令$\bm u_m=(1,0.5)^\top$，$\bm v_{e_1}=(0.8,0.4)^\top$，$\bm v_{e_2}=(0.9,0.5)^\top$，内积为1.00与1.15，双编码器把错误的$e_2$排在前面。若只召回一项，$e_1$已经丢失；召回两项后，设交叉编码分数为$(2.4,0.2)$，则$p(e_1\mid C(m),x)=e^{2.4}/(e^{2.4}+e^{0.2})\approx0.9002$，重排纠正了召回的次序。这里的概率只在候选集合内归一化，不代表“库中没有对应项”的概率，NIL需要另设阈值或拒绝分类器。

两段的代价不同。精确检索每个提及为$O(|\mathcal K|d)$，$|\mathcal K|$为库中实体数，近似近邻索引以少量召回损失换取查询速度；重排对$k$个候选各做一次完整编码。新实体只需写好描述、编码后加入索引，这是它能做零样本链接的原因；编码器更新后则要重算全部实体向量。

\subsection{文档全局一致性重排}
\label{sec:nlp-extraction-global-linking}

逐个提及独立链接，忽略了同一文档中的实体往往彼此相关。“海棠大学”与“计算机学院”出现在同一通知中，若后者明确链接到校园的学院，前者也更可能是校园的大学。把这种相关性写进目标，就得到全局链接。设文档有提及$m_1,\ldots,m_q$，联合赋值为$\bm e=(e_1,\ldots,e_q)$，$\Psi_i(e_i)$为第$i$个提及的局部分数，$\Phi(e_i,e_j)$为两实体的兼容分数：
\begin{equation}
  \widehat{\bm e}=\argmax_{\bm e\in\prod_i C(m_i)}
  \left[\sum_{i}\Psi_i(e_i)+
  \frac{2}{q-1}\sum_{i<j}\Phi(e_i,e_j)\right].
  \label{eq:nlp-extraction-global-linking}
\end{equation}
这是一个以提及为节点、候选实体为取值的全连接条件随机场，一元势为局部分数，成对势为兼容分数；系数$2/(q-1)$使成对项的总量不随提及数增长。

Ganea与Hofmann给出了一种端到端的实现\scite{ganea2017deep}。局部分数用实体嵌入$\bm x_e$对上下文词做注意力：每个上下文词$w$先按它与候选实体的最大相关度打分，只保留得分最高的若干词，归一化为权重$\beta(w)$，局部分数为$\Psi(e)=\bm x_e^\top\bm B\sum_w\beta(w)\bm x_w$，$\bm B$为对角矩阵；再用一个小型网络把它与$\log p(e\mid m)$组合。成对势为$\Phi(e,e')=\bm x_e^\top\bm C\bm x_{e'}$，$\bm C$也是对角矩阵。全连接图上精确求解式~\eqref{eq:nlp-extraction-global-linking}的代价为$O(k^q)$，$k$为每个提及的候选数，因此改用截断的\term{环置信传播}（loopy belief propagation）：在有环图上照常迭代传递消息，固定轮数后停止，以此近似每个提及的边缘分布。可参考\BookChapterRef{chap:pgm-variational}{卷六的变分推断与确定性近似章}，那里说明了消息传递在有环图上为何只是近似、何时可能不收敛。第$t$轮从提及$i$传给提及$j$的最大积消息为
\begin{equation}
  \mu^{(t+1)}_{i\to j}(e)=\max_{e'\in C(m_i)}\Bigl[\Psi_i(e')+\tfrac{2}{q-1}\Phi(e',e)+\sum_{k\neq i,j}\mu^{(t)}_{k\to i}(e')\Bigr],
  \label{eq:nlp-extraction-lbp}
\end{equation}
它表示“在$m_j$取$e$的条件下，$m_i$及其上游能贡献的最大分数”。Ganea与Hofmann还对消息做阻尼：把本轮消息归一化后与上一轮消息按系数$\delta\in(0,1]$加权平均，以稳定迭代。$T$轮后，提及$i$的置信为$b_i(e)=\Psi_i(e)+\sum_{k\neq i}\mu^{(T)}_{k\to i}(e)$，经softmax后再与先验组合成最终分数$\rho_i(e)$。整个展开的消息传递过程可微，训练以排序间隔损失$\sum_i\sum_{e\neq e_i^*}\max(0,\gamma-\rho_i(e_i^*)+\rho_i(e))$端到端更新嵌入、注意力和组合网络。

教学数值：大学的局部分数为校园0.6、小说0.8，学院为校园0.9、小说0.4；成对项（已含系数$2/(q-1)$）同领域取0.7，跨领域取0。四种组合的总分为校园—校园2.2、校园—小说1.0、小说—校园1.7、小说—小说1.9，全局选择校园组合，改变了大学的局部首选。用式~\eqref{eq:nlp-extraction-lbp}计算，学院传给大学的消息为：大学取校园时$\max(0.9+0.7,\ 0.4+0)=1.6$，取小说时$\max(0.9+0,\ 0.4+0.7)=1.1$；大学的置信为校园$0.6+1.6=2.2$、小说$0.8+1.1=1.9$，选校园。两个节点的图没有环，一轮消息即得到与穷举相同的答案；提及增多后图中有环，消息传递就只是近似。

Le与Titov进一步认为实体之间的兼容不止一种：两个实体可能同属一个机构，也可能是人物与其所在城市。他们引入$K$种隐关系，成对势写成$\Phi(e_i,e_j)=\sum_{r=1}^{K}a_{i,j,r}\,\bm x_{e_i}^\top\bm R_r\bm x_{e_j}$，每种关系有自己的对角矩阵$\bm R_r$，权重$a_{i,j,r}$由两个提及的上下文表示计算并归一化，无需关系标注\scite{le2018latent}。

全局一致性依赖“同一文档的实体相互关联”这一假设。一篇专门比较现实大学与小说学校的文章本来就跨领域，强制一致会把正确的跨领域链接改错；成对势的权重需要在包含这类文档的验证集上调整，NIL候选也应参与联合推断。

\subsection{生成式实体名称预测}
\label{sec:nlp-extraction-generative-linking}

BLINK把实体压成向量，需要为数百万实体存储并维护索引，还需要负例采样训练。De Cao等的GENRE换一种思路：每个实体有唯一的规范名称（如百科标题），链接就是在上下文条件下生成这个名称\scite{decao2021genre}。模型是BART，输入是用$[\mathrm{START\_ENT}]$与$[\mathrm{END\_ENT}]$标出提及的上下文$x_m$，输出名称的词元序列$\bm y_e$，评分为$\log p_\theta(\bm y_e\mid x_m)=\sum_u\log p_\theta(y_{e,u}\mid\bm y_{e,<u},x_m)$，用已链接提及按式~\eqref{eq:nlp-extraction-generation}训练。生成的名称必须是库中存在的名称，GENRE为此把全部名称建成一棵前缀树，在束搜索（第\ref{sec:nlp-generation-beam}小节）的每一步只允许能延续成某个合法名称的词元，机制与第\ref{sec:nlp-inference-grammar}小节的约束解码相同。

教学知识库把两个对象规范化为“海棠大学（校园）”和“海棠大学（小说）”。生成“海棠大学（”之后，前缀树只允许“校园）”与“小说）”两条延伸；若两条完成路径的概率分别为0.12与0.08，就选前者，再查名称表恢复$e_1$。输入中的别名可以相同，输出名称必须一一对应身份，知识库中的重名标题需要消歧后缀。

与双编码器相比，GENRE的存储只是一棵前缀树，不需要负例采样；生成时每个词元都能与提及上下文充分交互，却不必像交叉编码器那样逐个候选编码。它的局限是串行解码，束宽决定了能比较的候选数；新增实体只要加入前缀树就成为合法输出，但模型未必学会在新名称之间作出选择。无对应项时需要另设NIL输出或拒绝判断。

\section{实体关系抽取}
\label{sec:nlp-extraction-relations}

找全组织和活动名称以后，还要判断谁举办了什么。\term{关系抽取}（relation extraction, RE）
把文本中的实体与它们之间的语义关系绑定起来，输出有序的关系三元组
\[
  (s,r,o),\qquad r\in R\cup\{\mathrm{NA}\},
\]
其中$s$为\emph{主体}（subject），$o$为\emph{客体}（object），$R$是本体预先给定的关系类型集合，
$\mathrm{NA}$表示两者之间没有$R$中的任何关系。三元组是有序的：
$(\text{海棠大学},\text{主办},\text{读书会})$表示前者主办后者，交换两端就成了另一个命题。
本节用到的组织、人物和句子都是补充的教学材料，不由章首通知推断。
二元关系有两个参与者；“海棠大学于4月18日在海棠楼201主办读书会”这样的多元事实
还带有时间和地点。若把它拆成几条二元三元组存储，需要给它们一个共同的实例键，
让时间、地点都挂到同一次“主办”上；否则另一句“晨光书院于4月19日在松柏楼举办研讨会”
的日期和地点可能被拼进这条记录。多元事实的完整处理留给第~\ref{sec:nlp-extraction-events}节的事件抽取，
事件实例之间的先后与因果见第~\ref{sec:nlp-extraction-event-relations}节。

句内或文档级说明允许读取的上下文，封闭模式或开放关系说明标签是否预先枚举，
二元或多元说明参与结构。它们是可以组合的任务条件，不是互斥的模型类别。

\subsection{关系抽取的挑战与三元组重叠类型}
\label{sec:nlp-extraction-relation-challenges}

关系抽取的困难可以归为五项。其一是候选规模：一个文本中有$M$个实体时，
不含自关系的有向实体对有$M(M-1)$个，4个实体就有12对，其中绝大多数是NA，
类别极不平衡。其二是三元组之间共享实体，一个实体可能同时参与多条关系，
逐词元只给一个标签的输出表示装不下这种结构。其三是监督稀缺：逐句标注关系代价高，
借助知识库自动对齐得到的标签又带有大量噪声。其四是证据跨句：一个实体在文档中有多个提及，
关系可能要结合几句话才能判断。其五是类型不封闭：预先定义的$R$总会遗漏新出现的关系表达。
后续各小节分别回应这些困难：流水线方法和PURE处理候选规模与实体信息的注入，
远程监督与多示例学习处理监督稀缺，联合抽取处理实体共享，文档级方法处理跨句证据，
开放关系抽取处理不封闭的类型。

第二项困难需要先说清“共享”有哪几种。设句子的金标准三元组集合为$T$，
记$\operatorname{ent}(t)=\{s,o\}$为三元组$t=(s,r,o)$涉及的两个实体。
\citeauthor{zeng2018copyre}按实体共享方式把句子分为三类\cite{zeng2018copyre}：
\begin{itemize}
  \item \emph{无重叠}（Normal）：任意两条三元组不共享实体，即对所有$t\neq t'$有
    $\operatorname{ent}(t)\cap\operatorname{ent}(t')=\varnothing$；
  \item \term{实体对重叠}（entity pair overlap, EPO）：存在两条三元组共享同一对实体，
    即$\operatorname{ent}(t)=\operatorname{ent}(t')$而关系不同；
  \item \term{单实体重叠}（single entity overlap, SEO）：存在两条三元组恰好共享一个实体，
    即$|\operatorname{ent}(t)\cap\operatorname{ent}(t')|=1$。
\end{itemize}
同一句话可以同时属于EPO和SEO。图~\ref{fig:nlp-extraction-overlap}给出三类句子的实例。
“林舟出生并工作于海棠市”有$(\text{林舟},\text{出生地},\text{海棠市})$与
$(\text{林舟},\text{工作地},\text{海棠市})$，两条三元组的实体对完全相同，属于EPO。
“海棠大学和晨光书院共同主办读书会”有两条“主办”三元组，主体不同、客体都是读书会，
属于SEO；“林舟创办的晨光书院位于海棠市”中，晨光书院在一条三元组里是客体、
在另一条里是主体，也属于SEO。两个主体共享客体只说明它们各自主办该活动，
并不蕴含两个主体之间存在“合作”关系；后者若需要，应作为另一条有原文依据的三元组。

\begin{figure*}[htbp]
  \centering
  % 三元组重叠的三种类型：Normal、EPO、SEO。节点为实体，有向边为关系，边上标注关系类型。
\begin{tikzpicture}[
  x=1mm,y=1mm,figure picture,
  ent/.style={figure node,minimum height=7mm,minimum width=13mm,
    draw=FigureInput,fill=FigureInputFill,font=\figurefont\scriptsize},
  shared/.style={ent,draw=FigureOutput,fill=FigureOutputFill,line width=1pt},
  rel/.style={figure flow},
  rel label/.style={figure edge label,font=\figurefont\scriptsize},
  sent/.style={font=\figurefont\scriptsize,text=FigureMuted,align=center}
]
  % 三个分区
  \figurePanel{0}{0}{44}{40}{Normal}{FigurePanel}
  \figurePanel{50}{0}{44}{40}{EPO}{FigurePanel}
  \figurePanel{100}{0}{48}{40}{SEO}{FigurePanel}

  % Normal：两个三元组不共享实体（此处只画一条）
  \node[ent] (n1) at (9,19) {海棠大学};
  \node[ent] (n2) at (35,19) {读书会};
  \draw[rel] (n1) -- node[rel label,above=1mm] {主办} (n2);
  \node[sent] at (22,6) {海棠大学主办读书会};

  % EPO：同一有序实体对上有两种关系
  \node[shared] (e1) at (59,19) {林舟};
  \node[shared] (e2) at (85,19) {海棠市};
  \draw[rel] (e1.north east) to[bend left=28] node[rel label] {出生地} (e2.north west);
  \draw[rel] (e1.south east) to[bend right=28] node[rel label] {工作地} (e2.south west);
  \node[sent] at (72,4) {林舟出生并工作于海棠市};

  % SEO：两个三元组只共享一个实体
  \node[ent] (s1) at (110,26) {海棠大学};
  \node[ent] (s2) at (110,12) {晨光书院};
  \node[shared] (s3) at (139,19) {读书会};
  \draw[rel] (s1.east) -- node[rel label,above=.5mm,pos=.45] {主办} (s3.north west);
  \draw[rel] (s2.east) -- node[rel label,below=.5mm,pos=.45] {主办} (s3.south west);
  \node[sent] at (124,4) {海棠大学和晨光书院共同主办读书会};
\end{tikzpicture}

  \caption{三元组重叠的三种类型。方框为实体，有向边从主体指向客体，边上为关系类型，
  加粗方框为被多条三元组共享的实体，方框下方为对应的教学句。无重叠的句子中各三元组互不共享实体；
  EPO中同一对实体之间有两种关系；SEO中两条三元组只共享一个实体（此处为共同的客体）。
  三类划分依据Zeng等（2018）的定义，例句为教学构造。}
  \label{fig:nlp-extraction-overlap}
\end{figure*}

重叠决定了输出表示的最低要求。若每个词元只能带一个标签，EPO中的实体对要同时表达两种关系，
SEO中的共享实体要同时扮演两个角色，都装不下。若改为对每个有序实体对分别判断，
SEO自然可以表示；EPO则还要求同一实体对允许多个关系，即多标签而不是单标签输出。
表~\ref{tab:nlp-extraction-relation-methods}按这一标准对照本节的方法，
各行的具体机制在后续小节展开。

\begin{table*}[htbp]
  \centering
  \small
  \begin{tabularx}{\linewidth}{@{}>{\raggedright\arraybackslash}p{24mm}
    >{\raggedright\arraybackslash}X
    >{\centering\arraybackslash}p{16mm}>{\centering\arraybackslash}p{8mm}
    >{\raggedright\arraybackslash}X@{}}
    \toprule
    方法 & 输出表示与核心机制 & EPO & SEO & 主要代价 \\
    \midrule
    实体对分类（流水线、PURE） & 先识别实体，再对每个有序实体对分类 & 多标签时能 & 能 & 实体错误向下传播；每对单独编码 \\
    PCNN、选择性注意力 & 同一实体对的句子组成包，包级分类 & 多标签时能 & 能 & 依赖远程监督假设；只处理给定实体对 \\
    NovelTagging & 每词一个“位置—关系—角色”组合标签 & 否 & 否 & 无法表示重叠 \\
    CopyRE & 每三步生成关系并复制两端实体 & 能 & 能 & 只复制实体末词；三元组数目与次序固定 \\
    CasRel & 主体首尾指针，再按关系条件标客体首尾 & 能 & 能 & 主体漏检丢失其全部关系；训练与推断条件不一致 \\
    表格填充 & $n\times n$词对表，每格一个标签 & 否 & 能 & 一格只能存一种关系；填表顺序影响结果 \\
    TPLinker & 每种关系一组首首、尾尾词元对连接 & 能 & 能 & 输出规模$O(|R|n^2)$ \\
    GAIN、ATLOP、EIDER & 文档级实体对多标签分类 & 能 & 能 & 实体与提及须先给定；长文档编码 \\
    ReVerb & 词性模式匹配关系短语，再找左右论元 & 能 & 能 & 规则面向英语；关系不归一 \\
    \bottomrule
  \end{tabularx}
  \caption{关系抽取方法谱系。EPO与SEO两列表示输出表示能否容纳该重叠类型，
  不表示实际抽取的准确率。“多标签时能”指同一实体对允许输出多个关系时才能表示EPO。}
  \label{tab:nlp-extraction-relation-methods}
\end{table*}

\subsection{流水线关系抽取}
\label{sec:nlp-extraction-relation-pipeline}

最直接的做法分两步。第一步用第~\ref{sec:nlp-extraction-entities}节的实体识别得到实体集合
$\widehat{E}$；第二步对每个有序实体对$(e_i,e_j)$，$i\neq j$，计算
\[
  \hat r_{i,j}=\argmax_{r\in R\cup\{\mathrm{NA}\}}p(r\mid x,e_i,e_j).
\]
这种\emph{流水线}（pipeline）结构的好处是两个模型可以分别设计、分别训练，
实体识别还能借用现成的标注数据和模型。局限也来自分步。

第一个局限是\term{误差传播}（error propagation）：一条三元组正确，要求两个实体的边界与类型都对，
关系也对。设实体识别对每个金标准实体的召回率为$\rho_e$，两端识别近似独立，
关系分类器在实体正确的前提下判对的概率为$\rho_r$，则三元组召回率约为$\rho_e^{2}\rho_r$。
取教学数值$\rho_e=\rho_r=0.9$，三元组召回率只剩$0.9^{3}\approx0.729$；
第一步漏掉的实体，第二步无论多准都无法补回。错误的实体还会制造虚假候选，
降低精确率。第二个局限是训练与推断的条件不一致：关系模型通常在金标准实体上训练，
推断时面对的却是带错误的预测实体。第三个局限是代价：每个实体对若要单独编码，
一句话有12个候选对就要做12次前向计算。联合抽取（第~\ref{sec:nlp-extraction-joint}小节）
针对第一、二项，用一个模型同时输出实体和关系；PURE则说明，只要把实体信息正确地注入关系模型，
分开的两个模型仍有很强的竞争力，并给出处理第三项的近似计算。

\subsubsection{句内实体对分类与依存路径特征}
\label{sec:nlp-extraction-sentence-re}

给定句子和一个有向提及对$(m_s,m_o)$，关系分类要回答：这句话是否表达了$m_s$到$m_o$的
某种关系。早期方法手工构造特征向量$\bm{\phi}$，包括两个实体的词与类型、两者的相对距离、
两者之间的词，以及两实体在依存树上的\term{最短依存路径}（shortest dependency path）。
\citeauthor{bunescu2005shortest}的出发点是：依存树中连接两个实体的最短路径集中了判断关系所需的
大部分信息，路径以外的修饰成分多半与关系无关\scite{bunescu2005shortest}
依存树的构造见第~\ref{sec:nlp-analysis-dependency}小节。
补充句“海棠大学主办读书会”的路径为
“海棠大学$\xleftarrow{\text{nsubj}}$主办$\xrightarrow{\text{obj}}$读书会”，
弧的方向和标签告诉模型谁是施事、谁是受事，这正是区分$(s,r,o)$与$(o,r,s)$的依据。
分数由线性层给出，$\bm{z}=\bm{W}\bm{\phi}+\bm{b}\in\mathbb{R}^{|R|+1}$，
额外一维对应NA，再经softmax得到各关系的概率。

预训练编码器出现后，特征改由模型自己学。常用做法是在输入中插入\term{实体标记}（entity marker）：
在主体前后加\code{[E1]}、\code{[/E1]}，在客体前后加\code{[E2]}、\code{[/E2]}，
编码后取两个起始标记位置的向量拼接，
$\bm{\phi}=[\bm{h}_{\texttt{[E1]}};\bm{h}_{\texttt{[E2]}}]\in\mathbb{R}^{2d}$，
再以$\bm{W}\in\mathbb{R}^{(|R|+1)\times2d}$分类\scite{baldinisoares2019mtb}
标记把“哪两个跨度是当前候选”写进了输入，编码器的每一层都能据此调整上下文表示；
同一句话换一个实体对，就是一个不同的输入。两种标记区分了主客体，
交换实体顺序会得到另一个输入，模型不默认关系对称。
单标签设定用softmax交叉熵训练；若同一对允许多个关系（EPO），则改为逐关系的二分类，
或第~\ref{sec:nlp-extraction-atlop}小节的自适应阈值目标。

举一个教学数值。设三个类别“主办、支持、NA”的得分为$(2,0,-1)$，则
\[
  p(\text{主办})=\frac{\mathrm{e}^{2}}{\mathrm{e}^{2}+\mathrm{e}^{0}+\mathrm{e}^{-1}}
  \approx\frac{7.389}{8.757}\approx0.8438,
\]
正确标签的负对数似然约为0.1698，解码输出有向的“主办”关系。
依存路径特征的弱点可以用一个反例看出：若句子变成“海棠大学并不主办读书会”，
“并不”修饰“主办”，却不在两实体之间的最短路径上，只看路径的特征与原句完全相同，
分类器会给出同样的“主办”。路径之外的否定、情态和条件修饰必须另作特征或进入编码器表示；
带实体标记的编码器读入整句，可以利用这些成分，但是否学会利用仍取决于训练数据中有无此类样本。

\subsubsection{PURE实体标记与近似加速}
\label{sec:nlp-extraction-pure}

联合抽取的一个常见论点是：共享编码器能让实体和关系互相提供信息，从而克服流水线的缺陷。
\citeauthor{zhong2021pure}提出的PURE反其道而行，用两个\emph{不共享}参数的编码器分别做实体识别和关系分类，
结论是关系模型的关键在于尽早获得实体的边界与类型\scite{zhong2021pure}
它的实体模型是跨度分类器：枚举宽度不超过$W$的跨度$(i,j)$，以
$\bm{g}_{i,j}=[\bm{h}_i;\bm{h}_j;\bm{\phi}(j-i)]$表示，即式~\eqref{eq:nlp-analysis-span-representation}
去掉内部注意力汇聚的版本，再经前馈网络在“实体类型$\cup\{\varepsilon\}$”上分类，
$\varepsilon$表示不是实体。为了利用跨句信息，输入在句子左右各扩展若干上下文词元。

关系模型把实体模型的输出写进输入。对一个候选对，在主体两端插入带类型的标记
\code{<S:ORG>}、\code{</S:ORG>}，在客体两端插入\code{<O:EVT>}、\code{</O:EVT>}，
这里ORG、EVT是实体模型预测的类型。例如
\[
  \texttt{<S:ORG>}\ \text{海棠大学}\ \texttt{</S:ORG>}\ \text{主办}\
  \texttt{<O:EVT>}\ \text{读书会}\ \texttt{</O:EVT>}.
\]
编码后取两个起始标记的输出拼接，
$\bm{u}_{s,o}=[\hat{\bm{x}}_{\texttt{<S:ORG>}};\hat{\bm{x}}_{\texttt{<O:EVT>}}]\in\mathbb{R}^{2d}$，
经线性层在$R\cup\{\varepsilon\}$上做softmax，训练时实体取金标准。
与上一小节不带类型的标记相比，类型标记让关系模型直接得知“这是组织对活动”，
可以排除大量类型不相容的关系。若把实体模型的跨度向量拼接到关系分类器上，
实体信息只在最后一层才进入，编码器在计算上下文表示时并不知道当前是哪一对；
PURE的实验对照显示，这种后期融合不如在输入层插入标记有效。

代价是每个实体对都要重新编码一次。设句子有$n=30$个词元、$M=4$个实体，
共12个有序对，精确做法要做12次前向计算，每次长度约34。
PURE的近似做法把所有实体对的标记移到句末，标记与它所标跨度的首、尾词元共享位置编号，
并用注意力掩码规定：原文词元只看原文词元，标记看原文词元和同一实体对的其他标记。
这样原文词元的表示与实体对无关，只需计算一次；每对的四个标记各自读取原文并彼此交互。
一句话只需一次前向，长度为$30+4\times12=78$。
近似的代价是原文表示不再随实体对变化，标记信息不能反过来影响上下文，
论文报告推断加速8到16倍，F1仅小幅下降。

PURE说明，误差传播并不必然使流水线劣于联合模型：实体模型足够准、实体信息又在关系模型
输入层注入时，分开训练的两个模型可以与当时的联合模型相当或更好。但流水线的结构性问题仍在，
实体模型漏掉的实体不会再有关系预测，关系模型也仍在金标准实体上训练，
推断时面对的预测实体类型若有错，带类型的标记反而会把错误类型直接告诉关系模型。

\subsection{远程监督与多示例学习}
\label{sec:nlp-extraction-distant}

逐句标注关系的代价很高。\term{远程监督}（distant supervision）利用已有知识库自动产生训练标签：
若知识库中有三元组$(e_1,r,e_2)$，就把所有同时提及$e_1$与$e_2$的句子都标为关系$r$
\scite{knowledge-acquisition-mintz2009}
这一假设过强。设知识库记有$(\text{海棠大学},\text{主办},\text{读书会})$，句子
“海棠大学主办的读书会改至周五”确实表达了主办，而“海棠大学图书馆张贴了读书会海报”只是共同提及，
也会被标为“主办”。反方向上，知识库不完整，库中没有的关系会让表达了该关系的句子被标为NA。
常用的NYT数据集就是把Freebase中的关系与《纽约时报》语料对齐得到的\cite{knowledge-acquisition-riedel2010}，两类噪声都普遍存在。

\citeauthor{knowledge-acquisition-riedel2010}把假设放宽为“至少一次表达”：
同时提及两实体的句子中，至少有一句表达了库中的关系\scite{knowledge-acquisition-riedel2010}
这正是\term{多示例学习}（multiple-instance learning）的设定：训练数据由若干\emph{包}组成，
每个包是一组示例，标签只给到包，不给到其中的示例；正包中至少有一个正示例，
负包中全部是负示例。监督约束施加在整组对象上，示例级标签始终未知。
这一设定的一般讨论可参考\BookChapterRef{chap:machine-learning-paradigms}{卷三的范式简介章}。
在关系抽取中，同一实体对$(e_1,e_2)$的全部句子组成包$B=\{\bm{x}_1,\ldots,\bm{x}_m\}$，
包标签取知识库中的关系$y$。学习目标是包级的：
\[
  \max_{\bm{\theta}}\sum_{k}\log p_{\bm{\theta}}(y_k\mid B_k),\qquad
  p_{\bm{\theta}}(y\mid B)=f\bigl(\bm{x}_1,\ldots,\bm{x}_m\bigr),
\]
不同方法的区别在于聚合函数$f$怎样从句子得到包：PCNN只选一句，选择性注意力对全部句子加权。
两种方法都以给定实体对为前提，属于流水线的第二步；评价常用留出的知识库三元组作答案，
库中没有的正确预测会被算作错误，因此得到的精确率是保守估计。

\subsubsection{PCNN分段池化}
\label{sec:nlp-extraction-pcnn}

先看句子怎样编码。卷积网络用于关系分类时，常对整句的卷积特征做一次最大池化，
得到与句长无关的定长向量；但一次池化只保留每个滤波器在全句的最强响应，
丢掉了这个响应出现在两个实体之前、之间还是之后。对关系而言，这一位置信息很要紧：
“海棠大学主办读书会”中，动词出现在两个实体之间。
\citeauthor{zeng2015pcnn}的\emph{分段卷积网络}（piecewise convolutional neural network, PCNN）
按两个实体的位置把特征切成三段，分别池化\scite{zeng2015pcnn}

设句子有$n$个词元，主体与客体分别位于位置$p_1<p_2$。第$i$个位置的输入是词嵌入与
两个相对位置嵌入的拼接，
$\bm{x}_i=[\bm{w}_i;\bm{d}^{(1)}_{i-p_1};\bm{d}^{(2)}_{i-p_2}]\in\mathbb{R}^{d_w+2d_p}$，
两个位置嵌入分别编码该词到主体、客体的距离。宽度为$w$的第$k$个滤波器
$\bm{f}_k\in\mathbb{R}^{w(d_w+2d_p)}$在每个位置给出响应
$c_{k,i}=\bm{f}_k^{\top}\bm{x}_{i-w+1:i}+b_k$，
$\bm{x}_{i-w+1:i}$是相邻$w$个输入向量的拼接，句首句尾按需补零。
\term{分段最大池化}（piecewise max pooling）把位置分成三段
$I_1=\{i\leq p_1\}$、$I_2=\{p_1<i\leq p_2\}$、$I_3=\{i>p_2\}$，分别取最大：
\begin{equation}
  q_{k,\ell}=\max_{i\in I_\ell}c_{k,i},\quad \ell=1,2,3,\qquad
  \bm{g}=\tanh\bigl([q_{1,1},q_{1,2},q_{1,3},\ldots,q_{K,3}]\bigr)\in\mathbb{R}^{3K}.
  \label{eq:nlp-extraction-pcnn}
\end{equation}
$K$个滤波器得到$3K$维句子表示，经dropout和线性层在$R\cup\{\mathrm{NA}\}$上做softmax。
原文的设置为$d_w=50$、$d_p=5$、$w=3$、$K=230$，句子表示为690维。
举一个教学数值：$n=7$、$p_1=2$、$p_2=5$，某滤波器的响应为
$(0.2,0.9,0.1,-0.3,0.4,0.6,0.8)$。整句最大池化只得到0.9；
分段池化得到$(0.9,0.4,0.8)$，记录了实体之间也有中等响应、客体之后响应较强，
后续线性层可以据此区分不同的关系模式。

包级训练采用“至少一次表达”的最直接实现：对第$k$个包，先用当前模型找出最可能表达包标签的句子，
\[
  j^{*}=\argmax_{j}p(y_k\mid\bm{x}_{k,j}),\qquad
  J(\bm{\theta})=\sum_{k}\log p(y_k\mid\bm{x}_{k,j^{*}}),
\]
只用这一句计算梯度。例如一个“主办”包有三句，模型给出的$p(\text{主办})$为$(0.2,0.7,0.1)$，
则选第二句，该包的损失为$-\log0.7\approx0.3567$，另外两句不贡献梯度。
推断时，只要包内有一句被判为某关系，就给该实体对输出这一关系。
第~\ref{sec:nlp-extraction-dmcnn}小节事件检测中的DMCNN用的是同一类动态多池化，
切分点换成触发词和候选论元。PCNN的分段池化代价极低，却给卷积编码补上了实体相对位置的结构；
包级只选一句则浪费了其余句子中的有效信息，且训练初期模型选中的句子本身可能就是噪声句。

\subsubsection{句子级选择性注意力}
\label{sec:nlp-extraction-selective-attention}

只选一句的做法有两个问题：一个包中常有多句表达同一关系，其余各句的信息被丢掉；
选择是硬的，一旦选错，梯度全部落到噪声句上。
\citeauthor{lin2016selective}改用\term{选择性注意力}（selective attention），
为包内每句计算一个与关系相关的权重，再加权汇聚成包表示\scite{lin2016selective}
设包内句子表示为$\bm{x}_1,\ldots,\bm{x}_m\in\mathbb{R}^{d_c}$（例如PCNN的$3K$维输出），
关系$r$有查询向量$\bm{r}\in\mathbb{R}^{d_c}$，$\bm{A}$为可学习的对角矩阵，
\begin{equation}
  \begin{aligned}
    e_i&=\bm{x}_i^{\top}\bm{A}\bm{r},\qquad
    \alpha_i=\frac{\exp e_i}{\sum_{j=1}^{m}\exp e_j},\\
    \bm{s}&=\sum_{i=1}^{m}\alpha_i\bm{x}_i,\qquad
    p(r\mid B)=\operatorname{softmax}(\bm{M}\bm{s}+\bm{d})_r.
  \end{aligned}
  \label{eq:nlp-extraction-selective-attention}
\end{equation}
$e_i$衡量第$i$句与关系$r$的匹配程度，$\bm{M}\in\mathbb{R}^{(|R|+1)\times d_c}$是关系分类矩阵。
训练时用包标签对应关系的查询向量，目标是$\sum_k\log p(y_k\mid B_k)$；
推断时对每个候选关系分别用自己的查询计算包表示和分数，因为不同关系应关注不同的句子。

设一个三句包对“主办”的匹配分数为$(1,3,0)$，则
\[
  \bm{\alpha}=\frac{(\mathrm{e}^{1},\mathrm{e}^{3},\mathrm{e}^{0})}
  {\mathrm{e}^{1}+\mathrm{e}^{3}+\mathrm{e}^{0}}\approx(0.114,0.844,0.042),
\]
包表示主要来自第二句，但第一句仍贡献约11\%。这一权重可以看成在两个极端之间：
PCNN的选一句相当于把最大者的权重设为1，其余为0；对全部句子取平均则让噪声句与有效句等权。
注意力是硬选择的平滑版本，梯度能同时流向多个句子，噪声句的权重则随训练降低。
它仍建立在“至少一次表达”的假设上：若包内没有一句真正表达该关系，
注意力也只能在噪声句之间分配权重；知识库缺失造成的假NA包同样无法纠正。
远程监督数据中只含一句的包很多，此时注意力退化为该句本身，聚合带来的好处有限。

\subsection{实体与关系的联合抽取}
\label{sec:nlp-extraction-joint}

\term{联合抽取}（joint extraction）用一个模型同时输出实体与关系，
目的是避免流水线中实体错误无法纠正、实体与关系模型训练条件不一致的问题。
它的难点在于输出表示：既要装下实体边界，又要装下有向关系，还要允许重叠。
下面四种方法依次修补前一方法的缺陷。NovelTagging把三元组编码进一层序列标签，
却装不下重叠；CopyRE改为逐个生成三元组，同一实体可以复制多次，
但只复制实体末词，三元组的数目和次序也被固定；CasRel把关系看作从主体到客体集合的映射，
以首尾指针完整恢复多词实体并处理重叠，但推断依赖先预测出的主体；
表格填充与TPLinker把全部实体与关系同时标在词元对上，一步完成解码。
四种方法沿用同一个SEO教学句，词元依次为
“海棠大学／和／晨光书院／共同／主办／读书会”，金标准为$(1,\text{主办},6)$和$(3,\text{主办},6)$，
编号表示单词元提及。

\subsubsection{NovelTagging统一标注}
\label{sec:nlp-extraction-noveltagging}

把关系抽取改写成序列标注，就可以直接沿用第~\ref{sec:nlp-analysis-paradigm-tagging}小节的标注模型。
\citeauthor{zheng2017noveltagging}为此设计了一套组合标签\scite{zheng2017noveltagging}
每个标签由三部分组成：词在实体中的位置（B、I、E、S），该实体参与的关系类型，
以及它在关系中的角色（1为主体，2为客体）；不属于任何三元组的词标O。标签总数为$2\times4\times|R|+1$。
“海棠大学主办读书会”标为“S-主办-1，O，S-主办-2”。
解码时先按位置标签恢复实体，再把关系类型相同、角色互补的实体配成三元组；
同类实体有多个时，按就近原则配对。
模型是双向LSTM编码器加一个逐位置的LSTM解码器，训练目标在交叉熵上加大非O标签的权重：
\[
  L=-\sum_{t=1}^{n}\Bigl[\mathbb{I}(y_t=\mathrm{O})\log p(y_t\mid x)
    +\alpha\,\mathbb{I}(y_t\neq\mathrm{O})\log p(y_t\mid x)\Bigr],\qquad\alpha>1,
\]
因为O标签占绝大多数，不加权时模型倾向于把关系词也标成O。

统一标注把实体与关系放进同一个输出，一次解码完成，免去了流水线的两步。
它的根本限制在表示上：每个词只有一个标签，一个实体在整句中只能有一种关系、一个角色。
EPO句“林舟出生并工作于海棠市”中，“林舟”需要同时带“出生地-1”和“工作地-1”；
链式SEO句“林舟创办的晨光书院位于海棠市”中，“晨光书院”需要同时带“创办-2”和“位于-1”，
这些都无法用一个标签表示。共享客体的SEO教学句虽然可以标成
“S-主办-1，O，S-主办-1，O，O，S-主办-2”，就近配对却只会把读书会与较近的晨光书院配成一条，
海棠大学的三元组随之丢失。重叠在真实语料中很常见，这一限制促使后续工作改换输出表示。

\subsubsection{CopyRE复制式三元组生成}
\label{sec:nlp-extraction-copyre}

既然一个词元装不下多个角色，就不再给词元贴标签，而是把三元组一条一条地“写”出来：
同一个实体在不同三元组中出现几次，就复制几次。CopyRE用编码器—解码器实现这一想法
\scite{zeng2018copyre}
编码器为双向LSTM，输出$\bm{o}^{E}_1,\ldots,\bm{o}^{E}_n\in\mathbb{R}^{d}$；
解码器每三步生成一条三元组，第$t$步状态为$\bm{o}_t$。三类步骤输出不同的对象：
\begin{equation}
  \begin{aligned}
    &\text{关系步}\ (t\equiv1):&
      p(r)&=\operatorname{softmax}\bigl([\operatorname{selu}(\bm{W}_r\bm{o}_t+\bm{b}_r);\ a_t]\bigr)_r,\\
    &\text{实体步}\ (t\equiv2,0):&
      p(i)&=\operatorname{softmax}\bigl([u_{t,1},\ldots,u_{t,n};\ a_t]\bigr)_i.
  \end{aligned}
  \label{eq:nlp-extraction-copyre}
\end{equation}
同余均按模3计，$u_{t,i}=\operatorname{selu}(\bm{w}_e^{\top}[\bm{o}_t;\bm{o}^{E}_i])$是复制源位置$i$的分数，
$a_t=\operatorname{selu}(\bm{w}_{\mathrm{NA}}^{\top}\bm{o}_t+b_{\mathrm{NA}})$是NA的分数，selu为缩放指数线性激活，
$\bm{W}_r\in\mathbb{R}^{|R|\times d}$，$\bm{w}_e\in\mathbb{R}^{2d}$。
第一步在$|R|$种关系和NA上选关系；后两步在源句的$n$个位置上选一个位置复制，
依次得到主体和客体，复制客体时用掩码排除刚复制的主体位置。
解码预设生成固定数目的三元组，句中三元组不足时以NA三元组补齐。
这里的复制就是式~\eqref{eq:nlp-analysis-copy}去掉生成门的特例：
关系步相当于$p_{\mathrm{gen}}=1$，而且“词表”换成关系集合；实体步相当于$p_{\mathrm{gen}}=0$，
输出完全由源位置上的分布决定。论文另有每条三元组用一个独立解码器的变体。

在SEO教学句上，参考输出序列为
“主办，复制1，复制6，主办，复制3，复制6”，位置6被复制两次；
EPO句则为“出生地，复制1，复制4，工作地，复制1，复制4”。
训练用教师强制，对每一步取负对数似然；若某条三元组三步的概率依次为0.8、0.7、0.9，
它对损失的贡献为$-\log(0.8\times0.7\times0.9)\approx0.6852$。
生成式输出第一次让同一实体自由参与多条三元组，EPO和SEO都能表示。

代价有三项。其一，每次复制只取一个位置，CopyRE以实体的最后一个词元代表整个实体：
若“海棠大学”被切成“海棠／大学”，模型只复制“大学”，多词实体的边界无法完整恢复，
“海棠大学”与“晨光大学”在输出中也难以区分。其二，三元组本是一个集合，
教师强制却按标注的固定次序计算损失，模型以另一种次序生成同样的集合也会受罚。
其三，三元组数目上限需预先设定。CasRel用首尾指针代替单位置复制，
并把输出改为无次序的条件集合，同时消除了这三项问题。

\subsubsection{CasRel级联主体--关系--客体抽取}
\label{sec:nlp-extraction-casrel}

CasRel先找主体，再在指定主体条件下，按每个关系分别寻找客体\scite{wei2020casrel}
关系不再是两个已选好的实体之间的分类标签，而是从主体到客体集合的映射$f_r(s)$；
映射的值可以为空，也可以有多个客体，因此不需要预设三元组数目，也没有生成次序。
相应地，句子$x$的三元组集合$T$的似然分解为
\[
  p(T\mid x)=\prod_{s\in T}p(s\mid x)
  \prod_{r\in R}p\bigl(\{o:(s,r,o)\in T\}\mid s,x\bigr),
\]
其中$s\in T$指在$T$中作主体的实体，未与$s$构成关系$r$的客体集合为空。
设BERT编码器输出$\bm{h}_i\in\mathbb{R}^{d}$。主体起止头为
$p_i^{\mathrm{sh}}=\sigma(\bm{w}_{\mathrm{sh}}^\top\bm{h}_i+b_{\mathrm{sh}})$与
$p_i^{\mathrm{st}}=\sigma(\bm{w}_{\mathrm{st}}^\top\bm{h}_i+b_{\mathrm{st}})$，
每个位置都是二分类，不要求一句话只有一个主体。对于主体跨度$(a,b)$，取
$\bm{v}_s=(\bm{h}_a+\bm{h}_b)/2\in\mathbb{R}^{d}$，关系$r$的客体头为
\begin{equation}
  p_{i,r}^{\mathrm{oh}\mid s}
  =\sigma\!\left(\bm{w}_{r,\mathrm{oh}}^\top
       (\bm{h}_i+\bm{v}_s)+b_{r,\mathrm{oh}}\right),
  \qquad
  p_{i,r}^{\mathrm{ot}\mid s}
  =\sigma\!\left(\bm{w}_{r,\mathrm{ot}}^\top
       (\bm{h}_i+\bm{v}_s)+b_{r,\mathrm{ot}}\right).
  \label{eq:nlp-extraction-casrel}
\end{equation}
所有$\bm{w}$均为$d$维；关系条件体现在每个关系各自的客体头参数中，共$2|R|$个客体头。
主体标签由全部金标准主体构造，客体标签由“当前主体及当前关系”的全部客体构造。
令$\operatorname{BCE}(y,p)=-y\log p-(1-y)\log(1-p)$，
训练对主体两类边界与选定主体条件下各关系客体两类边界的BCE求和，
原方法每个句子随机抽取一个金标准主体作为条件。

在教学句中，主体起点、终点的正位置均为$\{1,3\}$。给定$s=(1,1)$，
“主办”客体两端的正位置均为6；给定$s=(3,3)$也如此。
假设两次客体起止概率为$(0.9,0.8)$与$(0.85,0.9)$，教学阈值取0.5，
便恢复两条三元组。第一对正边界的损失为$-\log0.9-\log0.8\approx0.3285$；
这只是一对边界的贡献，完整损失还包含负位置及主体预测。
一个主体可以在多个关系头上都有客体，EPO由不同关系头标出同一客体来表示；
多个主体分别作条件，SEO也自然得到表示。

解码先为每个主体起点匹配其后最近的主体终点，再按主体逐关系恢复客体。
这一最近匹配规则简单，却表示不了同一起点对应多个嵌套终点的情形。
更根本的问题是级联：主体漏检会使其全部关系失去后续预测机会；
训练时客体头以金标准主体为条件，推断时却以预测主体为条件，两者分布不同，
这与生成模型的暴露偏差（exposure bias）同类。推断时每个主体还要单独再算一遍客体头。
下面的单阶段方法正是针对这两点。

\subsubsection{表格填充与TPLinker词元对联合标记}
\label{sec:nlp-extraction-tplinker}

要去掉级联，需要一种输出表示，让全部实体和关系同时预测、彼此不互为条件。
\term{表格填充}（table filling）提供了这种表示\scite{miwa2014table}
对长度为$n$的句子建一张$n\times n$的表，对角格$(i,i)$记录第$i$词的实体标签
（如BILOU加类型），非对角格$(i,j)$记录第$i$词与第$j$词所在实体之间的关系标签，
标签带方向。实体和关系在同一张表上，任何两个实体之间都有自己的格子，SEO不成问题。
原方法按固定次序逐格填写，用特征和束搜索组织决策，后填的格依赖先填的格。
表格表示的限制在于一格只存一个标签：EPO中同一对实体的两种关系落在同一组格子上，无法并存。

TPLinker保留“在词元对上标注”的思路，但给每种关系单独一组标注，
并只标少数几种连接，从而在一步之内解码全部三元组\scite{wang2020tplinker}
它标出三类连接：实体首词元到尾词元（EH--ET），某关系的主体首词元到客体首词元（SH--OH），
以及主体尾词元到客体尾词元（ST--OT）。首首、尾尾两类连接合起来确定了多词实体的两端，
不必像CopyRE那样只取末词。
对$i\leq j$，以
$\bm{g}_{i,j}=\tanh(\bm{W}[\bm{h}_i;\bm{h}_j]+\bm{b})\in\mathbb{R}^{k}$
构造词元对表示，其中$\bm{W}\in\mathbb{R}^{k\times2d}$。
上三角的$n(n+1)/2$个位置按固定顺序铺成一条序列，这种布局称为\term{握手标注}（handshaking tagging）；
它只是共享的索引方式，不是先后执行的推断链。

EH--ET头标记实体首尾，输出“无连接／有连接”。
每个关系另有SH--OH和ST--OT两个头。因为只存上三角，后两类头各含“无连接、正向、反向”三个标签；
反向标签记录主体位置大于客体位置的情形，方向不会因只存上三角而丢失。
共有$2|R|+1$个标记头，各头将$\bm{g}_{i,j}$线性投影到自己的标签空间，
对全部格作交叉熵训练。每种关系各有自己的格子，同一实体对可以在多个关系上都有连接，EPO得以表示。

六词元教学句有21个格。实体头的正格为$(1,1)$、$(3,3)$、$(6,6)$；
“主办”的首部连接和尾部连接均在$(1,6)$、$(3,6)$取正向标签。
推理先由实体头恢复跨度集合$E$。对每个关系$r$，枚举
$s=(a,b),o=(c,d)\in E$，仅当SH--OH给出有向$a\to c$，
且ST--OT给出同方向的$b\to d$时，才恢复$(s,r,o)$。
本例得到两个共享客体的三元组。若只有$(1,6)$的首部连接而没有对应的尾部连接，
就不输出这条关系。

首部和尾部两类连接为什么都需要，可以看一个多词实体。设分词为“海棠／大学／主办／读书会”，
实体头给出$(1,2)$“海棠大学”，同时给出嵌套的$(1,1)$“海棠”（例如标注规范要求标出地名）。
若只有首部连接$1\to4$，无法判断主办者是$(1,1)$还是$(1,2)$；
尾部连接$2\to4$把主体的末端定在位置2，才唯一确定主体为“海棠大学”。
首尾两端同时约束，多词实体和嵌套实体都能完整恢复，这正是CopyRE只复制末词时缺少的信息。

与CasRel相比，TPLinker的所有标签都只以输入为条件，训练与推断完全一致，一次前向即得全部结果。
代价是输出规模：每个关系两个$n(n+1)/2$格的头，总计$O(|R|n^2)$个分类，
长句和大关系集合下显存与计算都显著增加。三类连接独立预测，
在实体边界重复、重叠密集的句子中仍可能错误组合；
实体边界、关系连接和完整三元组的正确率应分别统计。

四种方法每一步扩大了可表示的结构，也增加了解码步骤或输出规模。联合建模并不自动带来更高的准确率：
第~\ref{sec:nlp-extraction-pure}小节的PURE表明，实体信息在输入层注入时，分开训练的流水线同样有竞争力；
选择还取决于实体与关系标注是否齐全，以及句长和关系数目带来的计算代价。

\subsection{文档级关系抽取}
\label{sec:nlp-extraction-docre}

句内方法默认关系的证据在一句话里。DocRED的统计显示，
维基百科文档中相当一部分关系需要跨句才能判断\scite{yao2019docred}
\term{文档级关系抽取}（document-level relation extraction）的输入是整篇文档和其中的实体，
每个实体有若干提及，任务是对每个有序实体对输出关系集合，并可同时给出支持该关系的证据句。
构造三个教学句：$S_1$“海棠大学设有机器学习中心”；$S_2$“该中心负责读书会的筹备”；
$S_3$“海棠楼201配有投影设备”。海棠大学与读书会从未同句出现，若本体与标注规范允许
“所属机构经下属单位参与筹备”这类推理，就要连接$S_1$与$S_2$，并经过“该中心”与“机器学习中心”的共指；
若规范只标直接陈述的关系，就不应补出这条关系。

由此可以列出文档级方法要解决的问题：怎样把一个实体的多个提及聚合为实体表示；
怎样建模经由第三个实体的跨句推理；一个实体对可能有多个关系，也可能没有，阈值怎样决定；
证据句怎样利用。GAIN显式构造文档图，在图上推理；ATLOP不建图，
用预训练编码器的注意力找实体对的上下文，并以可学习的阈值处理多标签；
EIDER在ATLOP的基础上加入证据句预测，并融合原文与证据两路推断。
三种方法都以文档中的实体与提及已经给定为前提。

\subsubsection{文档图推理：GAIN}
\label{sec:nlp-extraction-gain}

跨句推理可以看成图上的路径：海棠大学与机器学习中心在$S_1$中相连，
机器学习中心经共指到$S_2$又与读书会相连。若把这些连接显式建成图，
图神经网络的消息传递就能把信息沿路径送到原本相距很远的实体之间。
GAIN为此构造两层图\scite{zeng2020gain}

第一层是\emph{异构提及图}（heterogeneous mention-level graph），
节点为全部实体提及，外加一个文档节点。三类边各有含义：
同一实体的不同提及之间连\emph{同实体边}，体现共指；
不同实体但出现在同一句中的提及之间连\emph{同句边}，体现局部交互；
每个提及都与文档节点相连，使任意两个提及至多两步可达。
编码器输出的词元表示拼接实体类型嵌入和实体编号嵌入后，取提及内词元的平均作为节点初值$\bm{h}^{(0)}_u$，文档节点的初值取编码器给出的文档表示。
图卷积按边的类型使用不同参数：
\begin{equation}
  \bm{h}^{(l+1)}_u=\sigma\Bigl(\sum_{k\in\mathcal{K}}\sum_{v\in N_k(u)}
    \bm{W}^{(l)}_k\bm{h}^{(l)}_v+\bm{b}^{(l)}_k\Bigr),
  \label{eq:nlp-extraction-gain-gcn}
\end{equation}
$\mathcal{K}$是三种边类型，$N_k(u)$是经类型$k$的边与$u$相连的邻居。
这种按关系类型区分参数的图卷积可参考\BookChapterRef{chap:graph-networks}{卷五的图神经网络章}；
此处只需知道，每层让一个节点汇聚各类邻居的信息，$L$层后信息可传播$L$步。
最终的提及表示拼接各层输出，$\bm{m}_u=[\bm{h}^{(0)}_u;\ldots;\bm{h}^{(L)}_u]$。

第二层是\emph{实体级图}（entity-level graph）。同一实体的提及表示取平均得到实体表示$\bm{e}_i$；
两个实体只要有提及出现在同一句中就连边，边表示
$\bm{e}_{i,j}=\sigma(\bm{W}_q[\bm{e}_i;\bm{e}_j]+\bm{b}_q)$。
判断实体对$(e_h,e_t)$时，GAIN考察经过第三个实体$e_o$的两跳路径，路径表示为
$\bm{p}^{o}_{h,t}=[\bm{e}_{h,o};\bm{e}_{o,t};\bm{e}_{t,o};\bm{e}_{o,h}]$，
再以实体对为查询对路径做注意力：
\[
  \beta_o=\operatorname{softmax}_o\bigl(\sigma([\bm{e}_h;\bm{e}_t]^{\top}\bm{W}_l\,\bm{p}^{o}_{h,t})\bigr),
  \qquad
  \bm{p}_{h,t}=\sum_{o}\beta_o\,\bm{p}^{o}_{h,t}.
\]
分类特征拼接两端实体、它们的差与积、文档节点表示和路径表示，
$\bm{I}_{h,t}=[\bm{e}_h;\bm{e}_t;|\bm{e}_h-\bm{e}_t|;\bm{e}_h\odot\bm{e}_t;\bm{m}_{\mathrm{doc}};\bm{p}_{h,t}]$，
经两层前馈网络对每个关系输出sigmoid概率，用二元交叉熵训练，按开发集选定的阈值输出。
计算量上，提及图的节点数为提及数加一，图卷积的代价与边数成正比；
路径注意力对每个实体对枚举与两端都相连的中间实体，最多$M-2$个，
全部$M(M-1)$个实体对合计为$O(M^3)$，实体多的长文档中这一项不可忽略。

图~\ref{fig:nlp-extraction-gain}展示教学文档上的两层图。提及图中，
$m_2$与$m_3$由同实体边相连，$m_1$与$m_2$、$m_3$与$m_4$分别由同句边相连。
实体图中，海棠大学$e_1$与读书会$e_3$之间没有边，但存在经$e_2$的路径；
路径注意力把这条路径的表示交给$(e_1,e_3)$的分类器。$S_3$中的海棠楼201与三者都不同句，
只经文档节点相连，在路径注意力中不构成两跳路径。

\begin{figure*}[htbp]
  \centering
  % GAIN 的两层图：左为异构提及图（同实体边、同句边、文档边），右为实体级图与两跳路径。
\begin{tikzpicture}[
  x=1mm,y=1mm,figure picture,
  mention/.style={figure node,minimum height=9mm,minimum width=13mm,
    draw=FigureInput,fill=FigureInputFill,font=\figurefont\scriptsize},
  docnode/.style={figure node,minimum height=7mm,minimum width=14mm,
    draw=FigureModel,font=\figurefont\scriptsize},
  entity/.style={figure node,minimum height=9mm,minimum width=15mm,
    draw=FigureInput,fill=FigureInputFill,font=\figurefont\scriptsize},
  intra/.style={draw=FigureStroke,line width=1.1pt},
  inter/.style={draw=FigureStroke,line width=.8pt,dash pattern=on 3pt off 1.5pt},
  docedge/.style={draw=FigureMuted,line width=.6pt,dash pattern=on .8pt off 1.2pt},
  path/.style={-{Stealth[length=1.8mm,width=1.2mm]},draw=FigureOutput,line width=1.1pt},
  query/.style={draw=FigureMuted,line width=.8pt,dash pattern=on 2pt off 2pt},
  note/.style={font=\figurefont\scriptsize,text=FigureMuted,align=center}
]
  \figurePanel{0}{0}{80}{48}{异构提及图}{FigurePanel}
  \figurePanel{88}{0}{60}{48}{实体级图与路径推理}{FigurePanel}

  % 提及图
  \node[docnode] (d) at (40,32) {文档节点};
  \node[mention] (m1) at (10,15) {$m_1$\\海棠大学};
  \node[mention] (m2) at (30,15) {$m_2$\\机器学习\\中心};
  \node[mention] (m3) at (50,15) {$m_3$\\该中心};
  \node[mention] (m4) at (70,15) {$m_4$\\读书会};
  \foreach \m in {m1,m2,m3,m4}{\draw[docedge] (d) -- (\m);}
  \draw[inter] (m1) -- (m2);
  \draw[intra] (m2) -- (m3);
  \draw[inter] (m3) -- (m4);
  \node[note] at (20,4.5) {句$S_1$};
  \node[note] at (60,4.5) {句$S_2$};

  % 实体图
  \node[entity] (e1) at (100,14) {$e_1$\\海棠大学};
  \node[entity] (e2) at (118,30) {$e_2$\\机器学习中心};
  \node[entity] (e3) at (136,14) {$e_3$\\读书会};
  \draw[path] (e1) -- (e2);
  \draw[path] (e2) -- (e3);
  \draw[query] (e1) -- node[figure edge label,font=\figurefont\scriptsize] {?} (e3);
  \node[note] at (118,4.5) {目标实体对$(e_1,e_3)$经$e_2$的两跳路径};
\end{tikzpicture}

  \caption{GAIN在教学文档“$S_1$海棠大学设有机器学习中心；$S_2$该中心负责读书会的筹备”上的两层图。
  左：异构提及图，粗实线为同实体边（“机器学习中心”与“该中心”共指），
  虚线为同句边，点线为各提及到文档节点的边。右：实体级图，同实体提及合并为实体节点，
  同句共现的实体之间连边；目标实体对$(e_1,e_3)$没有直接边，其关系由经$e_2$的两跳路径表示支持。}
  \label{fig:nlp-extraction-gain}
\end{figure*}

显式图的好处是推理路径可见，跨句信息有明确的传递通道。代价在于边由启发式规则决定：
共指是否正确、同句是否意味着相关，都直接影响图的结构，规则漏掉的连接无法由模型补上。
LSR对此的回应是不预先规定边，而把文档图当作潜在结构，用结构化注意力推断边的权重，
并在多轮中迭代细化\scite{nan2020lsr}
它减少了对启发式规则的依赖，却使图的含义更难解释，计算也更重。

\subsubsection{ATLOP文档实体对预测}
\label{sec:nlp-extraction-atlop}

ATLOP的出发点是：预训练Transformer的自注意力已经在全部词元之间建立了连接，
可以看成一张带权的全连接图，长距离依赖不一定要靠另建的图来传递
\scite{zhou2021atlop}
它因此不建图，而用三个部件分别处理多提及、实体对上下文和多标签阈值。
设文档编码为$\bm{H}\in\mathbb{R}^{n\times d}$，每个实体提及前后插入标记“*”。
实体$e$有提及标记表示$\bm{h}_{e,1},\ldots,\bm{h}_{e,M_e}$，逐维对数和指数聚合为
\begin{equation}
  [\bm{h}_e]_u=\log\sum_{j=1}^{M_e}\exp[\bm{h}_{e,j}]_u,
  \qquad u=1,\ldots,d.
  \label{eq:nlp-extraction-atlop-entity}
\end{equation}
它是最大值的平滑近似：每一维主要由取值最大的提及决定，其他提及也有少量贡献。
两个同名对象若被误并为一个实体，聚合结果会直接受到污染。

不同实体对需要文档中不同的上下文。设最后一层注意力有$H_{\!a}$个头，
先在每个头内平均同一实体全部提及标记的注意力行，
得到$\bm{A}_s,\bm{A}_o\in\mathbb{R}^{H_{\!a}\times n}$。
\term{局部上下文池化}（localized context pooling）对两个实体都关注的位置赋权：
\begin{equation}
  q_i=\sum_{h=1}^{H_{\!a}}A_{s,h,i}A_{o,h,i},
  \qquad a_i=\frac{q_i}{\sum_{j=1}^{n}q_j},
  \qquad \bm{c}_{s,o}=\bm{H}^{\top}\bm{a}\in\mathbb{R}^{d}.
  \label{eq:nlp-extraction-atlop-context}
\end{equation}
这里先在同一个头内相乘，再对头求和；先把各头平均再相乘会引入不同头之间的交叉项，
不是同一个计算。该权重反映模型对实体对的关注，不是人工标注的证据概率。

融合后令$\bm{z}_s=\tanh(\bm{W}_s\bm{h}_s+\bm{W}_{cs}\bm{c}_{s,o})$，
$\bm{z}_o$类似，四个投影矩阵均为$d\times d$。
将两向量各分为$K$组，每组$d/K$维，关系$r$的分数为
\begin{equation}
  \ell_r=\sum_{g=1}^{K}
  \bm{z}_{s,g}^{\top}\bm{U}_{r,g}\bm{z}_{o,g}+b_r,
  \qquad
  \bm{U}_{r,g}\in\mathbb{R}^{(d/K)\times(d/K)}.
  \label{eq:nlp-extraction-atlop-score}
\end{equation}
分组把每个关系的双线性参数从$d^2$降到$d^2/K$。以$d=768$、$K=12$为例，每组64维，
参数从$768^2=589{,}824$降到$12\times64^2=49{,}152$。主客体用不同的投影，相同的局部上下文不会使关系自动对称。
除实际关系外，模型还输出一个阈值类别TH，其分数$\ell_{\mathrm{TH}}$也依赖当前实体对。

多标签输出需要阈值。对全部实体对用一个固定阈值，忽略了不同实体对的分数尺度不同。
\term{自适应阈值损失}（adaptive thresholding loss）让阈值成为一个可学习的类别。
令$P$为该对的正关系集合，$N=R\setminus P$为负关系集合，目标把正关系推到TH以上、负关系推到TH以下：
\begin{equation}
  \begin{aligned}
    L_{+}&=-\sum_{r\in P}
      \log\frac{\exp\ell_r}
      {\exp\ell_{\mathrm{TH}}+\sum_{r'\in P}\exp\ell_{r'}},\\
    L_{-}&=-\log\frac{\exp\ell_{\mathrm{TH}}}
      {\exp\ell_{\mathrm{TH}}+\sum_{r'\in N}\exp\ell_{r'}},
    \qquad L_{\mathrm{AT}}=L_{+}+L_{-}.
  \end{aligned}
  \label{eq:nlp-extraction-atlop-loss}
\end{equation}
$L_+$的每一项的分母包含全部正类和TH，不是对每个正类与TH分别做二分类。
$P$为空时$L_+=0$，$L_-$仍训练TH。推理取
$\widehat{R}_{s,o}=\{r\in R:\ell_r>\ell_{\mathrm{TH}}\}$，空集表示NA。

\begin{example}[从局部上下文到多关系输出]
\label{ex:nlp-extraction-atlop}
为单独检查计算，使用四个位置、一个注意力头和二维表示。
令$\bm{H}$的四行依次为$(1,0)$、$(0,1)$、$(2,1)$和$(1,2)$，
前两行是主体的两个提及标记表示，后两行是客体的提及标记表示。
主体和客体分别逐维聚合，得到
\[
  \begin{aligned}
    \bm{h}_s&=(\log(1+\mathrm{e}),\log(1+\mathrm{e})),\\
    \bm{h}_o&=(\log(\mathrm{e}+\mathrm{e}^{2}),
                \log(\mathrm{e}+\mathrm{e}^{2})).
  \end{aligned}
\]
两实体平均后的注意力行取
\[
  \bm{A}_s=(0.4,0.3,0.2,0.1),\qquad
  \bm{A}_o=(0.1,0.2,0.3,0.4).
\]
逐位置乘积为$(0.04,0.06,0.06,0.04)$，其元素之和为0.2。
归一化后，权重向量为$\bm{a}=(0.2,0.3,0.3,0.2)$。
再用同一$\bm{H}$得到上下文向量$\bm{c}_{s,o}=(1.0,1.0)$。

用固定参数使融合向量均为$(0.5,0.5)$：取实体投影为零，
上下文投影为$\operatorname{arctanh}(0.5)\bm{I}_2$。
取$K=1$、偏置为零，并令$r_1,r_2,r_3,\mathrm{TH}$的双线性矩阵依次为
$4\bm{I}_2,2\bm{I}_2,-2\bm{I}_2,\bm{0}$，
式~\eqref{eq:nlp-extraction-atlop-score}给出分数$(2,1,-1,0)$。
若$P=\{r_1,r_2\}$，则
$L_+=2\log(\mathrm{e}^{2}+\mathrm{e}^{1}+1)-3\approx1.8152$，
$L_-=\log(1+\mathrm{e}^{-1})\approx0.3133$，总损失约2.1285。
解码同时保留$r_1,r_2$，排除$r_3$。这里的简化参数只用来贯通计算，
不表示训练好的模型会忽略实体表示。
\end{example}

与GAIN相比，ATLOP的取舍很清楚。它不需要共指、同句共现等规则来定义边，
也没有额外的图网络层，实现和训练都更简单；实体对的上下文由预训练注意力直接给出，
不受启发式边的遗漏所限。放弃的是显式的推理结构：经由哪个中间实体、走了哪条路径，
在ATLOP中只隐含在注意力权重里，无法像GAIN的路径注意力那样逐条检查；
需要多步推理且中间实体与两端都相距很远时，单层注意力乘积未必能把它们联系起来。
文档超过编码器长度上限时，还要分段编码再合并注意力与表示。

\subsubsection{EIDER证据增强与推理融合}
\label{sec:nlp-extraction-eider}

ATLOP的局部上下文是模型自己学到的关注，没有直接训练“哪些句子足以支持关系”。
DocRED标注了\term{证据句}（evidence sentence），可以作为额外监督。
EIDER在共享的文档编码器上增加实体对条件的证据预测，并在推理时结合完整文档与证据伪文档的关系分数
\scite{xie2022eider}
设文档有$J$句，第$j$句的表示$\bm{s}_j\in\mathbb{R}^{d}$由句内词元向量逐维对数和指数池化获得。
给定实体对上下文$\bm{c}_{s,o}$，第$j$句是证据的概率为
\begin{equation}
  p_j^{s,o}=\sigma(\bm{s}_j^\top\bm{W}_v\bm{c}_{s,o}+b_v),
  \qquad \bm{W}_v\in\mathbb{R}^{d\times d}.
  \label{eq:nlp-extraction-eider-evidence}
\end{equation}
证据按实体对选取，而不是为每个关系类别分别选一套句子。

设$D^+$为至少有一个正关系的训练实体对集合，$y_j^{s,o}$为该对第$j$句的证据标签，
证据损失为二元交叉熵，与关系损失加权相加：
\begin{equation}
  \begin{aligned}
    L_{\mathrm{evi}}
      &=-\sum_{(s,o)\in D^+}\sum_{j=1}^{J}
      \left[y_j^{s,o}\log p_j^{s,o}
       +(1-y_j^{s,o})\log(1-p_j^{s,o})\right],\\
    L&=L_{\mathrm{AT}}+\lambda L_{\mathrm{evi}}.
  \end{aligned}
  \label{eq:nlp-extraction-eider-loss}
\end{equation}
证据目标只在有正关系的实体对上计算，NA实体对没有可定义的证据。
两个任务共同更新编码器，证据监督让编码器更关注真正支持关系的句子。
没有人工证据标注时，EIDER用启发式规则构造证据：依次尝试两实体同句共现的句子、
含共指提及的句子、以及经一个桥接实体分别与两端相连的句组。规则证据可能包含与目标关系无关的句子，
它是弱监督，不能当作人工核验。

推理分两路。对每个实体对按阈值选出证据句，保持原句序拼成伪文档，
同时重建提及位置以及伪文档位置到原文位置的映射；
关系模型分别在完整文档和伪文档上给出分数$\ell_r^O$与$\ell_r^E$，各自减去各自的TH分数，
再以一个在开发集上拟合的标量$\tau$融合：
\begin{equation}
  S_r^O=\ell_r^O-\ell_{\mathrm{TH}}^O,\qquad
  S_r^E=\ell_r^E-\ell_{\mathrm{TH}}^E,\qquad
  p_r^{\mathrm{fuse}}=\sigma(S_r^O+S_r^E-\tau).
  \label{eq:nlp-extraction-eider-fusion}
\end{equation}
以$p_r^{\mathrm{fuse}}>0.5$判正等价于$S_r^O+S_r^E>\tau$。
完整文档一路保留了所有背景，证据一路去掉了干扰句，两路分数相加让二者互相补充。
这种\term{推理阶段融合}（inference-stage fusion）不需要重新训练分类器。
选出的证据为空或缺少某个实体的提及时，应退回完整文档一路的结果，并记录回退。

沿用上面的教学文档。设三句的证据标签为$(1,1,0)$，
预测概率为$(0.8,0.7,0.2)$，证据损失之和为
$-\log0.8-\log0.7-\log0.8\approx0.8030$。
取0.5为选句阈值，伪文档包含$S_1,S_2$。假设某个允许的筹备关系在两路的“关系、TH”分数
分别为$(0.6,0.2)$和$(0.8,0.1)$，则$S^O=0.4$、$S^E=0.7$。取教学融合阈值$\tau=0.3$，
得到$\sigma(0.8)\approx0.6900$，输出该关系。
若证据预测漏掉$S_1$，伪文档中缺少海棠大学与中心的联系，第二路可能变为$S^E=-0.5$，
融合概率降为$\sigma(-0.4)\approx0.4013$，关系被漏掉。
可见证据一路的收益依赖证据预测的召回；证据遗漏、背景误选与关系误判应分别检查。

表~\ref{tab:nlp-extraction-docre-methods}汇总本小节的文档级方法：从GAIN到ATLOP，跨句信息的传递从显式图转到预训练注意力；
从ATLOP到EIDER，模型关注的上下文从无监督的注意力转到有监督的证据句。

\begin{table*}[htbp]
  \centering
  \small
  \begin{tabularx}{\linewidth}{@{}>{\raggedright\arraybackslash}p{12mm}
    *{3}{>{\raggedright\arraybackslash}X}@{}}
    \toprule
    方法 & 解决的问题 & 核心机制 & 新的代价 \\
    \midrule
    GAIN & 跨句推理缺少传递通道 & 异构提及图上按边类型的图卷积；实体级图上的两跳路径注意力 & 边依赖共指与同句规则；增加图网络层 \\
    LSR & 启发式边可能遗漏或出错 & 以结构化注意力推断潜在文档图，并迭代细化 & 图的含义难解释；计算更重 \\
    ATLOP & 多提及聚合、实体对上下文与多标签阈值 & 对数和指数聚合；同头注意力乘积池化；可学习阈值类TH & 推理路径隐含在注意力中 \\
    EIDER & 关注的上下文未经证据监督 & 实体对条件的证据预测；原文与证据伪文档两路分数融合 & 需证据标注或规则；推断需两次编码 \\
    \bottomrule
  \end{tabularx}
  \caption{文档级关系抽取方法对照。各方法均假设文档中的实体及其提及已经给定；
  “新的代价”指相对于前一行方法新增的依赖或计算。}
  \label{tab:nlp-extraction-docre-methods}
\end{table*}

\subsection{开放关系抽取：ReVerb}
\label{sec:nlp-extraction-reverb}

前面的方法都要求预先定义$R$。\term{开放信息抽取}（open information extraction, OpenIE）
不预设关系类型，直接从句子中抽取“论元—关系短语—论元”的三元组，关系用原文的\term{关系短语}
（relation phrase）表示。它适合类型无法穷举、需要先探索语料再决定本体的场合。
早期开放抽取系统常产生两类坏结果：关系短语不连贯，由不相邻、拼不成意义的词凑成；
或信息不足，例如把“made a deal with”只抽成“made”，丢掉了承载关系含义的名词。
ReVerb用句法约束抑制这两类错误，再用词汇约束排除只对个别论元成立的过度特定短语\scite{fader2011reverb}

句法约束要求关系短语匹配词性模式$V\mid VP\mid VW^{*}P$：
$V$为动词（可带小品词和副词），$P$为介词、小品词或不定式标记，
$W$为名词、形容词、副词、代词或限定词；多个匹配相邻时合并，并取最长匹配。
这一模式覆盖了“hosts”“is held in”“has a meeting with”等形式，又排除了不以动词为核心的片段。
词汇约束要求关系短语在大规模语料中与足够多不同的论元对共现，
论文的阈值为20对，用以排除只对某一对论元成立的过度特定短语。
找到关系短语后，在其左右各取最近的名词短语作论元，并排除关系代词、疑问词和存在句中的“there”等。
最后用逻辑回归给每个抽取打置信度，$p(\text{正确}\mid\bm{\phi})=\sigma(\bm{w}^{\top}\bm{\phi})$，
特征描述句子长度、关系短语形式、论元位置等，训练数据是少量人工判定正误的抽取。

把补充句译写为英语教学句“Haitang University hosts the reading group.”，
关系短语为“hosts”，左侧最近的名词短语为“Haitang University”，右侧为“the reading group”，
得到开放三元组(Haitang University, hosts, the reading group)。
封闭模式可以把“hosts”归一为“主办”；开放结果保留原措辞及其位置，归一需要另外的映射。
ReVerb的规则与词汇统计面向英语，中文需要另行设计词性模式与资源。
开放三元组较短，“does not host”中的否定、时间状语和条件从句若被丢弃，原意就会改变；
这些限定需要下一小节记录中的专门字段承接。

\subsection{关系记录的恢复与核验}
\label{sec:nlp-extraction-relation-records}

模型输出的是分数和词元位置，可用的记录还需要对象身份、规范关系、限定条件和来源。
教学记录可写为$\langle s,r,o,\text{polarity},\text{time},\text{condition},\text{source}\rangle$；
否定关系记为负极性，而不是删除后留下同样的正关系。
算法~\ref{alg:nlp-extraction-relation-records}给出从分数到记录的步骤，
多条记录之间的对齐、合并与更新见第~\ref{sec:nlp-extraction-records}节。
结构检查使用JSON Schema约束字段类型、必填项和枚举值\cite{wright2022jsonschema}，
用SHACL约束图中的类型、属性取值范围和基数\cite{knublauch2017shacl}；
它们检查的是格式与结构，内容是否为原文所支持，要靠证据与身份的逐项核对。

\begin{algorithm}[关系记录的恢复与核验]
  \label{alg:nlp-extraction-relation-records}
  \begin{algorithmic}[1]
    \Require 文档及其版本号；候选三元组及分数、阈值；本体（关系方向、类型约束、基数）；
      别名到规范关系的映射；记录模式
    \Ensure 通过检查的关系记录，以及带拒绝原因的候选
    \For{每个分数超过阈值的候选$(s,r,o)$}
      \State 由词元跨度恢复原文字符区间；重复字符串按主客体与上下文消歧
      \State 开放关系短语按映射归一，同时保存原措辞；被动等表述按映射调整主客体方向
      \State 把$s,o$链接到库中身份或标为未链接提及；检查类型是否满足$r$的定义域与值域
      \State 从证据范围读取否定、时间和条件，填入极性、时间与条件字段；时间按日期规范解析
      \State 记录来源：文档版本、证据句与字符区间
      \State 用记录模式检查字段，用图约束检查类型与基数；不通过则附原因并拒绝
      \State 以证据为前提、记录的文字表述为假设做支持判定，并单独复核身份、数值、否定与条件
    \EndFor
    \State 精确匹配按对象键、规范关系与条件字段比较；语义等价评价另声明允许的别名、时间表达与方向改写
  \end{algorithmic}
\end{algorithm}

算法中的几道检查各管一层：类型兼容能拒绝“日期主办组织”，却不能判断类型兼容的边是否真实；
格式通过也不能抵消内容错误。“未参与主办”与“参与主办”只差几个字却不等价，
而“负责举办”与“主办”是否等价要由预先声明的映射决定。
流水线的候选遗漏、联合模型的连接误组、生成模型的无依据补项，最终都落到同一组对象、
关系、条件与来源的检查上；给定正确实体条件下的关系分数，不能代替从原文开始的端到端结果。

\section{事件抽取}
\label{sec:nlp-extraction-events}

“海棠大学主办读书会”给出一条关系，“读书会从4月17日改至4月18日”则描述了一次改变。
\term{事件}（event）是本体要求记录的一次发生、计划或状态变化；
\term{事件提及}（event mention）是文中对它的一次表达。
同一事件可以被多句重复提及，同一类型的事件也可以在不同日期多次发生。
事件本体规定类型及每类的角色，例如“改期”带有活动、原时间、新时间等角色；
第~\ref{sec:nlp-analysis-srl}节的语义角色标注则描述任意谓词及其参与者。
两者的目标和角色集合不同：语义角色结构可以作特征，但事件角色仍需按本体标注与验收。

事件抽取的挑战按输出层次排列：触发词多义、多词或多类型，长尾类型样本稀少；
论元依赖事件类型，一个角色可有多个论元，论元可能跨句；流水线中触发错误传给论元，
局部分类器无法表达跨任务约束；文档中论元散布多句，多个同类事件共享部分论元；
多次提及要合并为同一事件，同类而不同次的事件又不能合并。
方法谱系沿“流水线$\to$联合$\to$生成”展开：流水线先检测触发、再以事件为条件填充论元；
联合方法在一个模型中同时预测并用全局特征约束组合；生成方法把事件结构写成目标序列。
表~\ref{tab:nlp-extraction-event-methods}汇总各方法解决的问题与新增代价。

\begin{table*}[htbp]
  \centering
  \small
  \begin{tabularx}{\linewidth}{@{}>{\raggedright\arraybackslash}p{22mm} XXX@{}}
    \toprule
    方法 & 解决的问题 & 核心机制 & 新的代价或局限\\
    \midrule
    触发BIO/跨度 & 定位触发并分类 & 序列标注或跨度枚举（第9章范式） & 多类型需改为多标签；不含论元 \\
    DMCNN & 整句池化混淆多个事件 & 按候选位置分段最大池化 & 仍需枚举候选；固定窗口卷积 \\
    查询式检测 & 类型以整数编号，难以扩展 & 类型定义作查询，逐位置边界判别 & 每个类型一次编码；依赖定义质量 \\
    论元跨度分类 & 论元依赖事件 & 事件--候选交互表示上的角色分类 & 依赖候选覆盖 \\
    EEQA & 需要实体候选；角色语义未利用 & 角色问题＋阅读理解起止预测 & 每个角色一次编码；阈值需调 \\
    JRNN/OneIE & 流水线误差传播，跨任务约束缺失 & 共享编码＋全局特征＋束搜索联合解码 & 全局特征需人工设计模板 \\
    Text2Event & 分阶段标注与多个任务头 & 事件树线性化＋前缀树约束解码 & 生成片段需回对齐原文 \\
    BART-Gen/DEGREE & 文档级论元、少样本 & 模板条件生成＋受限复制 & 每个事件一次生成；解析回跨度 \\
    Doc2EDAG & 论元分散、多事件共享论元 & 按角色顺序扩展有向无环图 & 依赖角色顺序与实体抽取 \\
    事件共指 & 多次提及合并为实例 & 提及对或提及排序＋簇级约束 & 改期等合法差异需单独处理 \\
    \bottomrule
  \end{tabularx}
  \caption{事件抽取方法对照。前三行检测触发，第四、五行填充论元，其后依次为联合解码、
  生成、文档级多事件与实例合并。}
  \label{tab:nlp-extraction-event-methods}
\end{table*}

\subsection{事件提及与类型识别}
\label{sec:nlp-extraction-event-detection}

\term{触发词}（trigger）是用于定位事件提及的词或短语。
在本节教学本体中，“会议延期”的“延期”触发改期事件；
“拟延期”中的触发仍是“延期”，但状态是计划，不是已确认发生。
事实性可以用发生、否定、假设、计划等属性表达，它与事件类型是两个独立的输出字段。

类型识别的困难首先来自覆盖范围。ACE 2005定义了33个事件子类型，
MAVEN从维基百科构造了168个事件类型\scite{wang2020maven}
类型越细，长尾样本越少，近义类型的边界也越模糊。下面三种方法依次改变候选的表示：
标注与跨度方法复用第9章的范式，DMCNN让表示随候选位置变化，查询式检测把类型定义放进输入。

\subsubsection{触发词BIO标注与事件跨度分类}
\label{sec:nlp-extraction-trigger-tagging}

平面触发结构可以把实体BIO标签换成事件类型标签，
沿用第~\ref{sec:nlp-analysis-paradigm-tagging}节的发射分数、合法转移掩码
（式~\eqref{eq:nlp-analysis-legal-mask}）与CRF负对数似然；解码后还原触发跨度及类型。
跨度方法枚举$(i,j)$，以式~\eqref{eq:nlp-analysis-span-representation}的跨度表示
在“事件类型或非触发”之间分类，多词触发因而不必拆成B、I两步。
训练时未标注的候选充当负例，这要求语料对目标类型充分标注。

教学句“会议／拟／延期”中，第三词元标为\code{B-RESCHEDULE}，恢复字符区间$[3,5)$；
“拟”的区间$[2,3)$提供计划状态证据。状态属性可以用一个附加头预测：
设触发表示与周围上下文汇聚为$\bm{g}\in\mathbb{R}^{d}$，
$\bm{p}_{\mathrm{status}}=\operatorname{softmax}(\bm{W}_{f}\bm{g}+\bm{b}_f)$，
$\bm{W}_f\in\mathbb{R}^{4\times d}$对应发生、否定、假设、计划四类，
在有状态标注的样本上加入交叉熵。
同一跨度若允许属于多个类型，单一BIO类别或跨度softmax无法表达，
此时为每个类型设独立的sigmoid并以多热标签训练。

\subsubsection{DMCNN动态多池化触发分类}
\label{sec:nlp-extraction-dmcnn}

同一句话可能包含几个事件。若对整句卷积特征取一次最大池化，
句中最强的特征无论属于哪个候选都会进入表示，模型难以判断它与当前候选的关系。
DMCNN在卷积之后按候选位置划分池化区域，使“特征出现在候选哪一侧”进入表示\scite{chen2015dmcnn}
设第$i$个词的词向量与相对候选位置嵌入拼为$\bm{x}_i\in\mathbb{R}^{d_x}$，
宽$w$的卷积核组$\bm{W}_c\in\mathbb{R}^{F\times wd_x}$产生
$\bm{c}_i=f(\bm{W}_c[\bm{x}_{i};\ldots;\bm{x}_{i+w-1}]+\bm{b}_c)\in\mathbb{R}^{F}$，
整句得到$F\times n$的特征图。触发分类只有一个锚点$t$，沿它分成两段逐通道池化：
\begin{equation}
  [\bm{p}_{L}(t)]_u=\max_{i\leq t}[\bm{c}_i]_u,\qquad
  [\bm{p}_{R}(t)]_u=\max_{i>t}[\bm{c}_i]_u,\qquad u=1,\ldots,F.
  \label{eq:nlp-extraction-dmcnn}
\end{equation}
两段拼成$2F$维的句子级特征，再与词汇级特征（候选词及左右各一词的词向量）拼接，
经线性层和softmax在事件类型与NONE之间分类，以交叉熵训练。
当$t=n$使右段为空时，教学实现约定$\bm{p}_R(t)=\bm{0}$，填充位置不参与比较。
论元分类阶段有触发$t$和候选论元$a$两个锚点，按两者位置把特征图分为三段池化，
得到$3F$维特征；这与第~\ref{sec:nlp-extraction-relations}节关系分类中PCNN按两个实体
分三段池化是同一构造，区别只在锚点是触发与论元还是两个实体。

教学句为“会议／延期／后／报名／暂停／一天”，设某一卷积通道在六个位置的响应为$(1,4,2,3,5,0)$。
候选“延期”（$t=2$）两段池化得$(4,5)$，候选“暂停”（$t=5$）得$(5,0)$；
整句最大池化对两者都只给出5。DMCNN的代价是每个候选重做一次池化并需要大量NONE样本；
预训练编码器后来取代了固定窗口卷积，“表示随候选变化”的思想保留在以触发为条件的论元模型中。

\subsubsection{查询式事件检测}
\label{sec:nlp-extraction-event-query}

前两种方法把事件类型当作输出层的一个整数编号，新增类型就要改输出层并重新标注。
查询式检测把类型的名称与定义写成查询，例如“哪些词表达了活动时间的更改”，
与原文一起输入编码器，预测触发的起止位置；换成“哪些词表达了活动取消”，
对章首通知应输出空集。

设查询$q_k$与原文联合编码得到$\bm{H}^{(k)}\in\mathbb{R}^{n\times d}$，
按第~\ref{sec:nlp-extraction-mrc-ner}节MRC式实体识别的做法，
对每个位置分别做起点、终点二分类，再用配对头判断哪些起止对构成一个触发。
没有该类型的查询全部标负；逐位置二分类允许一个查询返回多个触发。
收益是类型语义进入编码，相近类型可共享定义中的线索；代价是$K$个类型要编码$K$次，
且检测质量依赖定义与示例的写法。

\subsection{事件论元填充}
\label{sec:nlp-extraction-arguments}

确定一次事件后，才能问它的时间、地点和参与者是什么。
\term{事件论元}（event argument）是以特定角色参与该事件的对象或表达，
角色集合由本体定义，不由句法位置直接决定。为检查绑定，补充构造一份通知：
“4月18日14:00，机器学习读书会在海棠楼201举行。4月19日10:00，
数据讨论会在松柏楼102举行。两场活动均由海棠大学组织。”
两场活动可以共享组织者“海棠大学”，却不能共享日期或地点。

跨度分类以事件为条件直接给候选判角色；EEQA把每个角色写成问题，
利用阅读理解的预训练并自然支持“无答案”。

\subsubsection{事件--论元跨度分类}
\label{sec:nlp-extraction-argument-spans}

给定事件$e$，从实体提及或一般文本跨度形成候选集合$C_e$；
后者可覆盖“4月18日14:00”这类未被实体识别列出的表达。
设事件表示$\bm{t}_e$与候选跨度表示$\bm{a}_{i,j}$均为$d$维
（跨度表示沿用式~\eqref{eq:nlp-analysis-span-representation}），
事件类型、句距等附加嵌入为$\bm{u}_{e,i,j}\in\mathbb{R}^{k}$。分类头为
\begin{equation}
  \bm{z}_{e,i,j}
  =\bm{W}[\bm{t}_e;\bm{a}_{i,j};
             \bm{t}_e\odot\bm{a}_{i,j};\bm{u}_{e,i,j}]+\bm{b},
  \quad \bm{W}\in\mathbb{R}^{(|A|+1)\times(3d+k)}.
  \label{eq:nlp-extraction-argument-spans}
\end{equation}
$A$是角色集合，额外一类为无角色，$\odot$是逐元素乘积，用来表达事件与候选的交互。
单角色规范对每个事件--候选对作softmax交叉熵，多角色规范改用独立二分类；
类型模式掩码屏蔽该事件类型不允许的角色。

训练标签按事件构造。“海棠楼201”对读书会是地点正例，对数据讨论会是无角色；
“海棠大学”对两者都是组织者。若地点与无角色两类的分数对前者为$(2,0)$、
对后者为$(-1,1)$，地点概率分别为$\sigma(2)\approx0.881$与$\sigma(-2)\approx0.119$：
候选文本相同，事件条件不同，角色判断随之不同。
候选若被句子窗口裁掉，分类头无从补救，因此要另报候选覆盖率。

\subsubsection{EEQA角色问答式抽取}
\label{sec:nlp-extraction-eeqa}

跨度分类依赖候选集合，角色名称本身的语义也没有进入模型。
EEQA把论元抽取改写为角色条件的阅读理解：问题包含事件类型、触发及角色，
模型从原文直接预测答案的起止位置，不需要前置实体候选\scite{du2020eeqa}
例如针对第一场“举行”询问“机器学习读书会何时举行”，对第二场则明确询问数据讨论会；
只问“何时举行”而不指明触发，两场活动的日期都会显得合理。

设问题与原文联合编码为$\bm{H}\in\mathbb{R}^{n\times d}$，其中含一个CLS位置。
起止向量$\bm{w}_s,\bm{w}_e\in\mathbb{R}^{d}$对全部有效位置评分，并分别跨位置归一化：
\begin{equation}
  P_s(i)=\frac{\exp(\bm{h}_i^\top\bm{w}_s)}
               {\sum_{u=1}^{n}\exp(\bm{h}_u^\top\bm{w}_s)},
  \qquad
  P_e(j)=\frac{\exp(\bm{h}_j^\top\bm{w}_e)}
               {\sum_{u=1}^{n}\exp(\bm{h}_u^\top\bm{w}_e)}.
  \label{eq:nlp-extraction-eeqa}
\end{equation}
有答案时以$-\log P_s(i^*)-\log P_e(j^*)$训练，没有该角色时起止标签都指向CLS。
与MRC式实体识别的逐位置二分类相比，这里每个边界头在全部位置上和为1，
多个答案只能在推理阶段通过阈值恢复。

推理枚举正文中合法的$i\leq j$及允许长度，排除
$P_s(i)<P_s(\mathrm{CLS})$或$P_e(j)<P_e(\mathrm{CLS})$的候选，再计算
\begin{equation}
  \delta(i,j)=P_s(\mathrm{CLS})+P_e(\mathrm{CLS})-P_s(i)-P_e(j),
  \label{eq:nlp-extraction-eeqa-threshold}
\end{equation}
保留$\delta(i,j)$不大于开发集选定阈值的候选，因此一个问题可以返回多个论元。
设有效位置只有“CLS／4月18日／14:00／海棠楼201”，
起点概率为$(0.10,0.55,0.05,0.30)$，终点概率为$(0.10,0.10,0.65,0.15)$，
最大答案长度为2。日期时间跨度$(2,3)$的$\delta=0.2-0.55-0.65=-1.00$，
地点跨度$(4,4)$的$\delta=0.2-0.30-0.15=-0.25$。
取教学阈值$-0.60$只保留前者，恢复“4月18日14:00”，
其边界训练损失为$-\log0.55-\log0.65\approx1.029$。

EEQA能找出实体识别漏掉的时间表达，代价是每个“事件、角色”组合编码一次，
绑定仍靠问题中的触发信息。

\subsection{联合与生成式事件抽取}
\label{sec:nlp-extraction-joint-generative}

前两小节是流水线：触发错了，论元跟着错；局部分类器也无法利用
“死亡事件的受害者很可能也是同句袭击事件的受害者”这类跨任务、跨实例的规律。

\subsubsection{OneIE全局特征与联合解码}
\label{sec:nlp-extraction-oneie}

联合抽取有两个先例。Li等用结构化感知机在束搜索中同时决定触发与论元，
特征里加入跨论元的全局特征\scite{nlp-extraction-li2013joint}
JRNN用双向循环网络编码句子，在同一次从左到右的预测中先判触发、再判论元，
并用记忆向量记录已经预测的触发与论元之间的依赖\scite{nlp-extraction-nguyen2016jrnn}
OneIE把这两条思路合到预训练编码器上：局部打分用神经网络，全局一致性用可学习权重的特征模板，
解码用束搜索\scite{nlp-extraction-lin2020oneie}

OneIE的输出是一张\term{信息网络}（information network）：节点是实体提及与事件触发，
边是实体间关系与触发--实体间的论元角色。计算分三步。
\emph{识别}：BERT编码得到词表示$\bm{x}_i$，前馈层给出BIO标签分数，
再经CRF层（第~\ref{sec:nlp-analysis-bmes-crf}节）解码实体与触发的跨度，损失记为$L_I$；
这一步只定边界，不定类型。\emph{分类}：节点表示$\bm{v}_i$取跨度内词表示的平均，
各任务的前馈网络给出节点类型分数$\hat{\bm{y}}^{t}_i=\mathrm{FFN}^t(\bm{v}_i)$
与边标签分数$\hat{\bm{y}}^{t}_{i,j}=\mathrm{FFN}^t([\bm{v}_i;\bm{v}_j])$，
以交叉熵$L_t$训练。若每个元素取局部最高分，得到局部最优图$\hat{G}$，
其局部分$s'(G)$是图中所选标签分数之和。
\emph{全局打分}：一组特征模板把图映射为全局特征向量$\bm{f}_G\in\mathbb{R}^{M}$，
例如“图中是否有多个袭击事件”“某类事件带有多个同角色论元的次数”
“死亡事件、受害者角色、地名实体这一组合出现的次数”。全局分为
\begin{equation}
  s(G)=s'(G)+\bm{u}^{\top}\bm{f}_G,
  \qquad
  L=L_I+\sum_{t}L_t+\bigl(s(\hat{G})-s(G^{*})\bigr),
  \label{eq:nlp-extraction-oneie-score}
\end{equation}
其中$\bm{u}\in\mathbb{R}^{M}$是与网络一起训练的特征权重，$G^{*}$是金标准图。
最后一项要求金标准图的全局分不低于局部最优图：局部最优图含有而金标准没有的组合，其特征权重被压低。

以构造句“来自海棠市的一名志愿者在事故中遇难”为例，局部分类器把“志愿者”（人）
判为死亡事件受害者的分数为1.2、无角色为0.4；把“海棠市”（地名）判为受害者的分数为1.0、
无角色为0.8。局部最优图$G_1$把两者都标为受害者，局部分2.2；
把“海棠市”改为无角色的$G_2$局部分2.0。设“死亡—受害者—地名”特征的权重已学得$-1.5$，
则$s(G_1)=2.2-1.5=0.7$，$s(G_2)=2.0$，联合解码选择$G_2$。
该权重若初始为0而金标准是$G_2$，最后一项为0.2，对权重的梯度为$f_{G_1}-f_{G_2}=1$，
梯度下降使权重变负。以上分数均为教学示例。

全局特征使搜索空间不再能逐元素分解，OneIE用算法~\ref{alg:nlp-extraction-oneie-beam}
的束搜索近似求解：节点按出现顺序逐个加入，每个节点只考虑局部分最高的$\beta_v$个类型，
每条新边只考虑局部分最高的$\beta_e$个标签，每完成一个节点的全部边后按全局分保留$\theta$张候选图。
原文取$\beta_v=\beta_e=2$、$\theta=10$。束搜索的一般形式见第~\ref{sec:nlp-generation-beam}节，
这里被扩展的不是词元序列，而是逐步长大的图。

\Needspace{14\baselineskip}
\begin{algorithm}[OneIE的束搜索联合解码]
  \label{alg:nlp-extraction-oneie-beam}
  \begin{algorithmic}[1]
    \Require 识别出的节点$v_1,\ldots,v_m$（按出现顺序），各节点与节点对的标签分数，
      全局特征权重$\bm{u}$，候选数$\beta_v,\beta_e$，束宽$\theta$
    \Ensure 全局分最高的信息网络$G$
    \State $B\gets\{\varnothing\}$\Comment{只含空图的束}
    \For{$i=1,\ldots,m$}
      \State $V_i\gets v_i$局部分最高的$\beta_v$个类型
      \State $B\gets\{G+(v_i,\ell)\mid G\in B,\ \ell\in V_i\}$\Comment{节点步}
      \For{$j=1,\ldots,i-1$}\Comment{边步；两端都是触发时跳过}
        \State $E_{j,i}\gets$边$(v_j,v_i)$局部分最高的$\beta_e$个标签（含无边）
        \State $B\gets\{G+(v_j,v_i,\ell)\mid G\in B,\ \ell\in E_{j,i}\}$
        \If{$|B|>\theta$}
          \State 按$s(G)=s'(G)+\bm{u}^{\top}\bm{f}_G$排序，只保留前$\theta$个
        \EndIf
      \EndFor
    \EndFor
    \State \Return $\argmax_{G\in B}s(G)$
  \end{algorithmic}
\end{algorithm}

OneIE解决了流水线的误差传播与跨任务约束缺失，代价是全局特征模板需要人工设计，
且束搜索只给近似最优；节点边界仍由识别步决定，边界错误不会在联合解码中被修正。

\subsubsection{Text2Event序列到结构生成}
\label{sec:nlp-extraction-text2event}

OneIE仍需要触发标注、实体标注和多个任务头。Text2Event换一个方向：
把事件类型、触发与角色论元组织成一棵树，线性化为带括号的序列，
让编码器--解码器模型（原文用T5）直接从原文生成它\scite{lu2021text2event}
训练只需要“原文、事件记录”对，不需要逐词元的触发标注。
在本节教学改期模式下，章首通知的目标结构为
\begin{quote}
（（改期 改至（活动 机器学习读书会）（原时间 4月17日）\\
（新时间 4月18日（周五）14:00）））
\end{quote}
最外层括号包住全部事件，每个事件是“（类型 触发 （角色 论元）……）”；
没有事件时目标是空括号“（）”。实际序列用独立的结构标记表示括号，
日期中的普通括号按原文字符处理，避免与树边界混淆。
训练用教师强制的条件负对数似然（第~\ref{sec:nlp-sft-masked-loss}节）。

自由生成可能写出本体中没有的类型名、不允许的角色，或原文中不存在的词。
Text2Event在解码时用\term{前缀树}（trie）约束每一步的候选：前缀树把一组合法词元序列
按公共前缀合并，从根走到某个节点得到一个合法前缀，节点的子边就是下一步允许的词元。
解码器根据已生成的括号判断当前处在树的哪个位置：
刚打开一个事件时，候选限于事件类型名构成的前缀树；类型写完后进入触发位置，
候选限于输入原文中的词元；打开一个角色时，候选限于该事件类型允许的角色名；
角色名之后又回到原文词元。记当前位置的允许集合为$A(y_{<t})$，约束分布为
\begin{equation}
  \tilde{p}(y_t\mid y_{<t},x)\propto
  p_{\bm{\theta}}(y_t\mid y_{<t},x)\,\mathbb{I}[y_t\in A(y_{<t})],
  \label{eq:nlp-extraction-text2event-mask}
\end{equation}
它是式~\eqref{eq:nlp-inference-constrained}的特例：允许集合由事件模式与原文共同决定，
括号的配对则像第~\ref{sec:nlp-inference-grammar}节那样需要记录嵌套深度。
例如类型名有“改期”和“取消”时，前缀树在根下分出两条边，生成“改”之后只允许“期”。

训练顺序也是Text2Event的一部分。直接学完整事件树较难，原文先训练\term{子结构学习}
（substructure learning）阶段：目标只含“（类型 触发）”和“（角色 论元）”这样的独立片段，
让模型先学会从原文定位片段；再在完整结构上继续训练。这一课程只改变学习难度的次序，
不改变最终要恢复的目标。生成结束后，解析器把树转回记录，并把触发与论元字符串回对齐到原文位置。

Text2Event与第~\ref{sec:nlp-extraction-entities}节的UIE同属
第~\ref{sec:nlp-analysis-paradigm-generation}节的约束生成范式，后者侧重用统一结构语言覆盖多种任务。
其局限在回对齐：原文有两个相同日期字符串时，生成结果无法确定指哪一处；
树结构合法也不保证论元挂在正确的事件下。

\subsubsection{事件模板条件生成}
\label{sec:nlp-extraction-event-template}

Text2Event的目标序列是人工设计的括号语言，与预训练语料的自然语言相差较远。
模板条件生成改为让模型“填空”：每个事件类型配一句带角色占位符的自然语言模板，
模型在文档条件下生成填好的模板。BART-Gen以BART为生成器，处理文档级论元抽取
\scite{li2021wikievents}
输入是“\code{<s>}模板\code{<s></s>}文档\code{</s>}”，模板中各角色写成占位符\code{<arg>}，
文档中的当前触发两侧加标记\code{<tgr>}，以区分同一文档中的多个事件。
教学模板为“\code{<arg>}于\code{<arg>}在\code{<arg>}举行”，三个占位符按位置分别对应活动、
时间和地点；对读书会，目标是“机器学习读书会于4月18日14:00在海棠楼201举行”。
缺失的角色保留\code{<arg>}，同一角色的多个论元用“和”连接。

为防止生成原文中没有的论元，BART-Gen把输出分布限制在输入词元集合$V_c$
（模板与文档中出现过的词元）上：
\begin{equation}
  p(x_i=w\mid x_{<i},c)=
  \begin{cases}
    \dfrac{\exp(\bm{h}_i^{\top}\bm{e}_w)}{\sum_{w'\in V_c}\exp(\bm{h}_i^{\top}\bm{e}_{w'})}, & w\in V_c,\\[2ex]
    0, & w\notin V_c,
  \end{cases}
  \label{eq:nlp-extraction-bartgen-copy}
\end{equation}
其中$\bm{h}_i$是解码器第$i$步输出，$\bm{e}_w$是词元$w$的嵌入。
与式~\eqref{eq:nlp-analysis-copy}的复制机制相比，这里没有“生成或复制”的门控，
而是只允许复制；它减少了库外词，但输入词元的任意组合仍可能不是原文的一个连续跨度。
生成后按模板对齐解析：模板中的固定词（“于”“在”“举行”）作为分隔，
把各占位符位置上的文本切出并映射回角色，再在文档中定位字符跨度；
找不到连续匹配的输出标记为解析失败，而不是强行就近对齐。

DEGREE把模板生成扩展到触发检测\scite{nlp-extraction-hsu2022degree}
输入除了原文，还有事件类型的自然语言描述、若干示例触发词和模板；
输出一句“触发是……”加填好的模板，再确定性地解析出触发与论元。
每个事件类型各生成一次，类型描述和关键词作为弱监督进入输入，便于少样本学习。
模板生成的共同代价是每个“文档、事件或类型”都要生成一次，解析回跨度与论元归属仍需核验。

\subsection{文档级事件抽取：Doc2EDAG}
\label{sec:nlp-extraction-doc2edag}

句子级方法在金融公告上遇到两个问题。一是\emph{论元分散}：
一笔股权质押的质押人、股数、质权人、起止日期可能分别出现在不同句子中，
没有一句话包含完整记录。二是\emph{多事件}：一篇公告可以记录同一股东的多笔质押，
各笔共享质押人而股数、质权人不同，独立地对每个角色挑选论元会把不同笔的字段混在一起。
Doc2EDAG针对中文金融公告提出端到端的文档级框架，并采用无触发词设定：
标注只给事件表格，不要求标出触发词\scite{nlp-extraction-zheng2019doc2edag}

模型分三段。\emph{实体识别}：每句经Transformer编码为$\bm{h}_i\in\mathbb{R}^{N_w\times d_w}$，
接CRF做BIO标注，标签来自与表格匹配的论元，损失为$L_{\mathrm{er}}$。
\emph{文档级实体编码}：对每个实体提及的词表示做最大池化得到提及向量，
对每句的词表示做最大池化得到句向量$\bm{c}_i$；加上句子位置嵌入后，
全部提及向量与句向量一起送入第二个Transformer交换信息，
再把同名提及的输出最大池化合并，得到$N_e$个实体的文档感知表示
$\bm{E}^{d}\in\mathbb{R}^{N_e\times d_w}$与句表示$\bm{C}^{d}\in\mathbb{R}^{N_s\times d_w}$。
\emph{EDAG生成}：对$\bm{C}^{d}$最大池化得到文档向量，对每个事件类型做二分类判断是否出现
（损失$L_{\mathrm{tr}}$）；对出现的类型，按人工规定的角色顺序逐步扩展一张
\term{实体有向无环图}（entity-based directed acyclic graph, EDAG）。

EDAG由事件表格构造：每条记录按角色顺序写成论元链表，缺失角色用特殊节点NA，
再把所有链表的公共前缀合并。图~\ref{fig:nlp-extraction-doc2edag}的教学公告含两笔质押：
“股东甲将3000万股质押给乙银行”“甲另将1000万股补充质押给丙证券”。
按“质押人$\to$质押股数$\to$质权人”的顺序，两条链表共享前缀“甲”，
在第二个角色处分叉；每条从根到叶的完整路径对应表格的一行，恢复表格只需枚举路径。

\begin{figure}[htbp]
  \centering
  % Doc2EDAG按角色顺序扩展实体有向无环图：两笔股权质押共享质押人，在第二个角色处分叉。
\begin{tikzpicture}[
  x=1mm,y=1mm,figure picture,
  root/.style={figure node,minimum width=16mm,minimum height=8mm,
    draw=FigureInput,fill=FigureInputFill,font=\figurefont\scriptsize,align=center},
  ent/.style={figure node,minimum width=17mm,minimum height=7mm,
    draw=FigureModel,fill=FigureModelFill,font=\figurefont\scriptsize},
  leaf/.style={figure node,minimum width=17mm,minimum height=7mm,
    draw=FigureOutput,fill=FigureOutputFill,font=\figurefont\scriptsize},
  role/.style={font=\figurefont\scriptsize,text=FigureMuted,align=center},
  flow/.style={figure flow}
]
  \node[root] (r) at (0,0) {质押事件\\（已触发）};
  \node[ent] (a) at (27,0) {甲};
  \node[ent] (s1) at (54,8) {3000万股};
  \node[ent] (s2) at (54,-8) {1000万股};
  \node[leaf] (p1) at (81,8) {乙银行};
  \node[leaf] (p2) at (81,-8) {丙证券};
  \draw[flow] (r) -- (a);
  \draw[flow] (a) -- (s1);
  \draw[flow] (a) -- (s2);
  \draw[flow] (s1) -- (p1);
  \draw[flow] (s2) -- (p2);
  \node[role] at (27,17) {第1步\\质押人};
  \node[role] at (54,17) {第2步\\质押股数};
  \node[role] at (81,17) {第3步\\质权人};
  \node[role,anchor=west] at (92,8) {记录1};
  \node[role,anchor=west] at (92,-8) {记录2};
\end{tikzpicture}

  \caption{Doc2EDAG的路径扩展（教学示例）。左端是事件类型触发节点，自左向右依次决定
  质押人、质押股数、质权人三个角色。每一步对当前叶节点的路径，对全部候选实体做二分类，
  选中的实体成为新节点；两笔质押共享“甲”，在第二步分叉为两条路径，每条完整路径对应一条事件记录。}
  \label{fig:nlp-extraction-doc2edag}
\end{figure}

路径扩展把“填一张表”拆成一串子任务：对当前角色$k$与当前叶节点的路径，
判断每个实体是否应作为该角色的论元接到路径上；没有实体被选中时接一个NA节点。
为了让子任务知道路径上已经有哪些实体，模型为每条路径维护记忆张量
$\bm{M}\in\mathbb{R}^{(N_s+k-1)\times d_w}$：初始化为$\bm{C}^{d}$，
每扩展一步追加所选实体的表示（NA追加零向量）。扩展概率为
\begin{equation}
  \begin{aligned}
  \bm{E}^{r}&=\mathrm{Transformer}_3\bigl([\bm{M};\bm{E}^{d}]+\bm{r}_k\bigr)_{\text{实体部分}},\\
  p(\text{扩展实体}e)&=\sigma(\bm{w}^{\top}\bm{E}^{r}_{e,:}+b),
  \end{aligned}
  \label{eq:nlp-extraction-doc2edag-expand}
\end{equation}
其中$\bm{r}_k$是可训练的角色指示嵌入，加到每一行上；取输出中对应$N_e$个实体的行做二分类。
每个子任务的交叉熵相加为$L_{\mathrm{dag}}$，总损失为
$\lambda_1L_{\mathrm{er}}+\lambda_2L_{\mathrm{tr}}+\lambda_3L_{\mathrm{dag}}$。
在图示例子中，第三步在“3000万股”之后只有“乙银行”标为1，在“1000万股”之后只有“丙证券”标为1：
两条路径的记忆不同，同一实体在两个子任务中得到相反的标签，这正是区分两笔记录的机制。

Doc2EDAG用文档级编码解决论元分散，用共享前缀的图解决多事件共享论元；
代价是角色顺序需人工规定且前面选错会影响后续扩展，实体识别错误无法补回，推理逐角色串行。

\subsection{事件共指与合并}
\label{sec:nlp-extraction-event-coreference}

一篇通知和它的转发描述同一次读书会，连续两周的读书会却是不同实例。
\term{事件共指}（event coreference）把表达同一次事件的提及归为一簇。挑战是触发措辞不同、
论元缺失或冲突、改期使同一事件的提及日期不一致，以及簇级传递一致性。
方法沿用第~\ref{sec:nlp-analysis-coreference}节的实体共指框架，差别在于特征要比较论元。

\subsubsection{提及对与提及排序}
\label{sec:nlp-extraction-event-coref-models}

提及对模型对事件提及对$(a,b)$构造特征$\bm{\phi}_{a,b}$：触发语境的相似度、
已绑定论元的一致性（参与者是否共指、时间地点是否兼容）、事实性是否相同、来源关系，
预测$p_{a,b}=\sigma(\bm{w}^{\top}\bm{\phi}_{a,b})$，按式~\eqref{eq:nlp-analysis-mention-pair}的
二元交叉熵训练。编码器或双线性头可以替换线性评分，但正例必须指同一次事件，不能只是同类型。
提及排序模型则对每个事件提及在较早提及与“无先行项”之间做一次softmax
（式~\eqref{eq:nlp-analysis-antecedent-softmax}），并以正确先行项集合的边缘似然训练
（式~\eqref{eq:nlp-analysis-coreference-marginal}），从而比较多个候选而非逐个过阈值。

事件共指比实体共指更依赖簇级约束。构造三条记录：$a$为编号\code{ML-0418}的通知，
$b$为省略编号但明确转引该通知的报道，$c$为编号\code{ML-0425}的下周读书会。
教学分数取$p_{a,b}=0.92$、$p_{b,c}=0.86$、$p_{a,c}=0.10$，阈值0.8。
对过阈值的边取传递闭包会把三者合为一簇；受约束的凝聚聚类按分数从高到低合并，
先合并$a,b$，再因可靠编号冲突拒绝把$c$并入，得到$\{a,b\}$与$\{c\}$。
日期不同并不总是硬冲突：同一活动改期后，旧通知与新通知的日期不同却应共指，
因此要先辨认“原定”“改至”等角色，再决定日期比较的是哪一个字段。

\subsubsection{事件共指的评价}
\label{sec:nlp-extraction-event-coref-metrics}

事件共指沿用第~\ref{sec:nlp-analysis-coreference-metrics}节的MUC、B³、CEAF$_e$与CoNLL F1。
事件语料中单例很多：大部分事件只被提及一次。MUC不计单例，
B³与CEAF$_e$却会因为大量单例天然正确而偏高，三者对同一系统可能给出差别很大的结论。
\term{BLANC}（bilateral assessment of noun-phrase coreference）把兰德指数的思路用于共指，
对共指链接与非共指链接分别计算F1再平均\scite{nlp-extraction-recasens2011blanc}
设参考答案中的共指提及对集合为$C_k$、非共指提及对集合为$N_k$，
系统输出相应为$C_r$、$N_r$，则
\begin{equation}
  F_c=\frac{2|C_k\cap C_r|}{|C_k|+|C_r|},\qquad
  F_n=\frac{2|N_k\cap N_r|}{|N_k|+|N_r|},\qquad
  \mathrm{BLANC}=\frac{F_c+F_n}{2}.
  \label{eq:nlp-extraction-blanc}
\end{equation}
用第9章的五提及例子核算：参考簇$\{a,b,c\},\{d,e\}$，系统簇$\{a,b\},\{c,d,e\}$。
共指对$C_k=\{ab,ac,bc,de\}$，$C_r=\{ab,cd,ce,de\}$，交集2，$F_c=0.5$；
其余各6对为非共指，交集为$\{ad,ae,bd,be\}$，$F_n=4/6\approx0.667$；BLANC约为0.583。
直接用兰德指数（全部10对中判对的比例）得0.6，与BLANC相近；
但单例增多时非共指对数按提及数的平方增长，兰德指数几乎只反映非共指链接，BLANC则让两类链接各占一半。

\section{事件关系抽取}
\label{sec:nlp-extraction-event-relations}

两次事件没有合并为一次，不意味着它们彼此无关。报名先于活动、改期由场地冲突引起、
讨论是会议的一个环节，分别表达时间、因果和组成联系。本节在事件已经确定的前提下，
抽取事件之间的这些有向关系。MAVEN-ERE在同一批文档的同一组事件上标注了共指、
时间、因果与子事件四类关系，含约5.8万条因果关系和1.6万条子事件关系，
使四类任务可以在共享输入上比较和联合训练\scite{wang2022mavenere}
EventRelBench则以选择题形式在句级与文档级两种粒度上评估大模型对事件关系的理解
\scite{gong2025eventrelbench}
给定事件对再选答案，与从原文恢复全部关系是两种任务设置。

事件关系的挑战有四项。一是规模：$m$个事件有$O(m^2)$个有序对，而绝大多数对没有关系，
标签极度稀疏。二是证据分布：许多关系由上下文隐含，或跨越多个句子。
三是全局一致性：时间关系要满足逆关系对称与传递，组成关系要无环，
局部分类器各自决策时会产生“甲先于乙、乙先于丙、丙又先于甲”的矛盾。
四是类型间的相互约束：原因不晚于结果，父事件在时间上包含子事件，共指事件共享全部关系。
方法谱系据此分三层：局部分类器对事件对打分；全局推断（区间约束传播、整数线性规划）
在局部分数上恢复一致结构；约束感知的学习（结构化学习、可微逻辑约束、多类联合训练）
把约束提前到训练阶段。表~\ref{tab:nlp-extraction-event-relation-methods}在节末汇总。

\subsection{时间表达与时间关系}
\label{sec:nlp-extraction-temporal}

时间信息有三层对象：“明天”是待归一化的时间表达，某个日期属于哪次事件是关联，
两事件是否先后或重叠是关系。三层彼此独立：开始时间都找对了，也未必知道结束时间和区间关系；
章首通知只给出活动开始时刻而没有时长，不能补出结束时间。
日期粒度的“4月18日”表示一个日历日范围，“4月16日18:00前”是严格早于该时刻的条件。
三层任务依次由规则归一化、事件对分类器和区间代数与整数线性规划处理。

\subsubsection{HeidelTime与SUTime时间归一化}
\label{sec:nlp-extraction-time-normalization}

时间表达的归一化以规则为主：先用模式识别日期、时刻、时长或周期表达，
再结合语言资源与参照时间输出TIMEX3标准值。HeidelTime把模式、归一化资源与文本类型
条件分开组织\scite{strotgen2010heideltime}
SUTime以规则和组合机制识别、解析时间表达，原始资源以英语为主\scite{chang2012sutime}
中文规则的教学适配是：匹配“$k$天后”，选出参照日期$d_0$，计算$d_0+k$个日历日。
设补充句发布于2025年4月15日：“明天18:00前报名，三天后14:00开始。”
“明天”以文档日期为参照，得到\code{2025-04-16T18:00:00+08:00}；
“三天后”从发布日期加三个日历日，得到\code{2025-04-18T14:00:00+08:00}，与章首通知一致。
参照并不总是文档日期：“读书会4月18日举行，三天后公布纪要”应以活动日为参照得到4月21日；
规则因此输出值、粒度与参照来源，缺少参照时保留$d_{\mathrm{ref}}+3\text{天}$。

\subsubsection{事件对时间关系分类}
\label{sec:nlp-extraction-temporal-classifier}

给定有序事件对$(e_i,e_j)$，编码两事件的触发、相关时间表达及上下文，
形成$\bm{g}_{i,j}\in\mathbb{R}^{d}$（例如两触发表示的拼接），关系头为
$p(r\mid e_i,e_j,x)=\operatorname{softmax}(\bm{W}\bm{g}_{i,j}+\bm{b})_r$，
$\bm{W}\in\mathbb{R}^{|T|\times d}$，以有向时间标签的交叉熵训练。
标签集$T$随语料而定。TimeBank-Dense要求对一定窗口内的全部事件对标注，
标签为先于、后于、包含、被包含、同时与模糊六类\scite{nlp-extraction-cassidy2014dense}
MATRES只比较事件的开始点，把标签减为先于、后于、相等、模糊四类，以提高标注一致性
\scite{nlp-extraction-ning2018matres}
标签集决定端点含义：同一对事件在两个语料中的正确答案可能不同。

对“报名截止时点、活动开始时点”这一有向对，两者相差44小时，应输出“先于”；
交换事件顺序时应输出逆关系“后于”。

神经分类器出现之前，CAEVO用\term{筛子架构}（sieve architecture）组织多个分类器
\scite{nlp-extraction-chambers2014caevo}
它把若干规则与机器学习分类器按精确率从高到低排成序列，每个筛子只给它有把握的事件对加标签；
每加一批标签就计算传递闭包，把推得的关系一并加入；后面的筛子不能推翻或违反前面的决定。
筛子架构的好处是天然保持一致，代价是早期错误会被传递闭包放大，且各筛子的次序需要人工调定。

\subsubsection{Allen区间约束与全局关系恢复}
\label{sec:nlp-extraction-allen}

局部分类器之间的矛盾需要全局方法消除。Allen用13种基本关系描述两个非退化时间区间的位置：
先于、相接、重叠、同始而先结束、严格位于内部、同终而晚开始，这六种及其逆关系，再加相等
\scite{allen1983intervals}
设区间$I=[a,b)$与$J=[c,d)$，$a<b$，$c<d$。先于要求$b<c$，相接要求$b=c$，
重叠要求$a<c<b<d$，严格位于内部要求$c<a<b<d$。事件只有开始时刻时，
不能直接套用完整区间关系，应保留未知终点或改用点时间关系。

不确定的关系用基本关系的集合$R_{i,j}$表示。组合$R_{i,j}\circ R_{j,k}$给出两条边允许的$(i,k)$关系，
路径一致性反复收缩
\begin{equation}
  R_{i,k}\leftarrow R_{i,k}\cap
      (R_{i,j}\circ R_{j,k}).
  \label{eq:nlp-extraction-allen-propagation}
\end{equation}
例如$\mathrm{before}\circ\mathrm{before}=\{\mathrm{before}\}$，
$\mathrm{meets}\circ\mathrm{meets}=\{\mathrm{before}\}$。
若$a$先于$b$、$b$先于$c$，而$R_{a,c}=\{\mathrm{before},\mathrm{overlaps}\}$，更新后只剩before；
若原集合只有after，交集为空，暴露冲突。对一般的Allen关系网络，路径一致性不是可满足性的充分条件。

要兼顾模型分数，可为每个有序对与基本关系设二元变量$x_{i,j}^{r}$，求
\[
  \max_{\bm{x}}\sum_{i\ne j}\sum_r s_{i,j}^{r}x_{i,j}^{r}
  \quad\text{s.t.}\quad
  \sum_r x_{i,j}^{r}=1,\qquad x_{i,j}^{r}=x_{j,i}^{r^{-1}},
\]
其中$s_{i,j}^{r}$是局部对数概率，$r^{-1}$是逆关系；三元组合约束为
\begin{equation}
  x_{i,j}^{r}+x_{j,k}^{s}
  -\sum_{t\in r\circ s}x_{i,k}^{t}\leq1,
  \label{eq:nlp-extraction-temporal-ilp}
\end{equation}
即前两条边都选中时，第三条必须落入组合允许的集合。
这种“局部分类加整数线性规划推断”与第~\ref{sec:nlp-analysis-srl}节语义角色标注中的
ILP恢复是同一构造（式~\eqref{eq:nlp-analysis-srl-ilp}）。
教学三事件的候选先后边权为$a\to b:9$、$b\to c:8$、$c\to a:6$，反方向均为0，
要求三者有确定先后。局部取最大会成环；枚举六种全序，$a<b<c$总分17最高，
它舍去权重6的边，并由约束得到$a<c$。

上述推断只在测试时使用约束，训练时局部分类器并不知道约束的存在。
Ning等把推断放进训练循环，用结构化感知机学习时间关系\scite{nlp-extraction-ning2017structured}
每篇文档的预测不再是逐对取最大，而是带约束的整体最优；
参数按“金标准图的特征减去约束推断图的特征”更新，如算法~\ref{alg:nlp-extraction-temporal-structured}。
原文还用这一推断步骤处理语料中缺失的关系标注。代价是每次参数更新都要求解一次整数线性规划，
文档中事件数增多时训练变慢。

\Needspace{12\baselineskip}
\begin{algorithm}[带全局约束推断的时间关系结构化学习]
  \label{alg:nlp-extraction-temporal-structured}
  \begin{algorithmic}[1]
    \Require 训练文档$\{(x^{(n)},\bm{y}^{(n)})\}$，事件对特征$\bm{\phi}(x,i,j,r)$，
      约束集合$\mathcal{C}$（逆关系对称、三元组合），轮数$E$
    \Ensure 权重$\bm{w}$（取各步平均）
    \State $\bm{w}\gets\bm{0}$
    \For{$e=1,\ldots,E$}
      \For{每篇文档$(x,\bm{y})$}
        \State $\hat{\bm{y}}\gets\argmax_{\bm{y}'\in\mathcal{C}}
          \sum_{i\ne j}\bm{w}^{\top}\bm{\phi}(x,i,j,y'_{i,j})$\Comment{求解式~\eqref{eq:nlp-extraction-temporal-ilp}}
        \If{$\hat{\bm{y}}\ne\bm{y}$}
          \State $\bm{w}\gets\bm{w}+\sum_{i\ne j}\bigl(\bm{\phi}(x,i,j,y_{i,j})-\bm{\phi}(x,i,j,\hat{y}_{i,j})\bigr)$
        \EndIf
      \EndFor
    \EndFor
    \State \Return 平均后的$\bm{w}$
  \end{algorithmic}
\end{algorithm}

\subsection{因果关系}
\label{sec:nlp-extraction-causal}

“天气变差后会议取消”表达先后，“因天气变差会议取消”表达文本中的因果主张。
因果抽取识别的是这样的主张及其方向，不是在观测数据上估计干预效应。
挑战有三项：多数因果是隐式的，没有“因为”“导致”等连接词；
标注数据少，常用的EventStoryLine只有258篇文档\scite{nlp-extraction-caselli2017esc}
Causal-TimeBank只有318对因果事件\scite{nlp-extraction-mirza2014causal}
跨句的因果对远多于句内，在EventStoryLine中句间因果约3900对，句内约1800对
\scite{nlp-extraction-chen2022ergo}
三项挑战分别对应下面三类方法：事件对分类处理显式与隐式情形，
知识增强扩充训练数据，文档级因果图聚合跨句证据。

\subsubsection{显式与隐式因果识别}
\label{sec:nlp-extraction-causal-explicit}

显式因果可由连接词、因果谓词和依存路径产生候选，再把原因和结果绑定到已识别的事件；
隐式因果没有这些标志，只能依靠事件对的上下文表示。两种情形共用一个有向分类头。
设两事件的上下文表示为$\bm{u},\bm{v}\in\mathbb{R}^{d}$，关系$r$的双仿射分数为
$z_r=\bm{u}^{\top}\bm{W}_r\bm{v}+\bm{w}_r^{\top}[\bm{u};\bm{v}]+b_r$，
$\bm{W}_r\in\mathbb{R}^{d\times d}$；标签可为原因指向结果、结果指向原因、无因果，
以有向事件对标注作交叉熵训练；连接词可作为特征或输入标记加入，但“后”“随后”这类时间连接词不表示因果。
“天气变差使交通中断，也导致会议取消”中，交通中断与会议取消只有共同原因，不能补出二者之间的因果。

因果与时间相互约束：在一般的因果语义下，原因不晚于结果。
Ning等把这一点写进联合推断\scite{nlp-extraction-ning2018joint}
时间变量$x^{r}_{i,j}$沿用式~\eqref{eq:nlp-extraction-temporal-ilp}，
再加因果变量$c_{i,j}\in\{0,1\}$与约束$c_{i,j}\leq x^{\mathrm{before}}_{i,j}$，
目标是两类局部分数之和。于是有把握的“后于”会否决弱的因果判断，强的因果判断也能纠正时间关系；
持续过程或预期原因需要放宽这一约束。

\subsubsection{知识增强的因果识别}
\label{sec:nlp-extraction-causal-knowledge}

标注稀少时，一个自然的想法是用外部知识自动造出更多训练句。
难点在于：自动找到的句子要同时满足两条，事件对确有因果，句子确实表达了这种因果。
KnowDis用远程监督解决前者、用常识过滤解决后者\scite{nlp-extraction-zuo2020knowdis}
第一步扩充事件对：对每个已标注的因果对，用WordNet取两事件中心词的同义词与上位词，
用VerbNet取同一动词类中的词，交叉组合成新的候选因果对；
再以TransE式的间隔损失（第~\ref{sec:nlp-extraction-kge}节）学习一个因果表示空间，
按两事件的距离排序，保留最近的10\%。第二步远程标注：
在纽约时报语料中随机抽取5\%的句子，凡含有这些事件对的句子标为因果候选，
这与第~\ref{sec:nlp-extraction-relations}节关系抽取的远程监督同理，也带来同样的噪声。

第三步用常识过滤噪声。KnowDis从COPA这类“前提--原因或结果”的常识推理数据和已标注数据中，
统计原因侧词$i$与结果侧词$j$的共现，从必要性$p(i\mid j)$与充分性$p(j\mid i)$两个方向度量因果强度，
并用指数$\alpha$惩罚高频词：
\begin{equation}
  \mathrm{CS}(i,j)=\Bigl(\frac{p(i\mid j)}{p(i)^{\alpha}}\Bigr)^{\lambda}
                   \Bigl(\frac{p(j\mid i)}{p(j)^{\alpha}}\Bigr)^{1-\lambda}.
  \label{eq:nlp-extraction-knowdis-cs}
\end{equation}
$\alpha=1$时两个因子都等于点互信息的指数形式$p(i,j)/(p(i)p(j))$。
候选句按两事件之间文本片段的平均因果强度排序：两事件由因果连接词相连的保留前50\%，其余只保留前10\%。
最后用人工数据预训练的检测器重新标注这些句子，只保留仍判为因果的部分，并在训练中逐步加入。

KnowDis的数据生成与识别模型相互独立。LearnDA把二者放进\term{对偶学习}（dual learning）框架：
识别器$I$做“句子$\to$关系”，生成器$G$做“事件对与关系$\to$句子”，两个方向互为对偶任务
\scite{nlp-extraction-zuo2021learnda}
事件对仍由WordNet、VerbNet扩充，并从FrameNet、PDTB2的因果连接词在KBP语料中找出新的对；
生成器分三步造句：为事件分配与原句同类型且BERT表示最相近的实体，
用掩码语言模型在事件与实体之间填充连接词语，再按填充词元的平均概率与和已标注句子的
表示相似度加权打分，保留前一部分。两个模型通过两种奖励交互：因果奖励$R_c$取识别器对句子判断的概率，
判对为正、判错为负；语义对齐奖励$R_s$取生成器在给定关系下生成该句词元的平均概率。
原始循环中生成器造句、识别器判断，生成器按$R_s$与$R_c$的加权和做策略梯度更新；
对偶循环中识别器判断真实句子、生成器重建，识别器按同样两种奖励更新。
与KnowDis相比，LearnDA让生成随识别任务调整，代价是同时训练两个模型并承受强化学习的方差。

\subsubsection{跨句与文档级因果图}
\label{sec:nlp-extraction-causal-document}

句间因果缺少局部线索，但常能通过中间事件推出：若“断裂”导致“停电”，
“停电”又导致“抢修”，就可以推出“断裂”与“抢修”的因果。
ERGO把这种传递性直接建进图结构\scite{nlp-extraction-chen2022ergo}
它不以单个事件为节点，而以\emph{事件对}为节点：对文档中$m$个事件的每个有序对$(e_i,e_j)$
建一个节点，初始表示为两事件上下文表示的拼接
$\bm{v}^{(0)}_{i,j}=[\bm{h}_{e_i};\bm{h}_{e_j}]\in\mathbb{R}^{2d}$；
只有共享至少一个事件的两个节点之间才连边，因为没有共同事件的两对之间不存在中介链条。
文档编码用BERT分段处理或直接用Longformer。

在这张关系图上，ERGO用关系图Transformer更新节点：节点$a$对邻居$b$的注意力为
\begin{equation}
  \alpha_{a,b}=\operatorname*{softmax}_{b\in\mathcal{N}(a)}
  \frac{(\bm{v}_a\bm{W}_q)(\bm{v}_b\bm{W}_k)^{\top}}{\sqrt{d_k}},\qquad
  \bm{v}'_a=\sum_{b\in\mathcal{N}(a)}\alpha_{a,b}\,\bm{v}_b\bm{W}_v,
  \label{eq:nlp-extraction-ergo-attention}
\end{equation}
多头输出拼接后经$\bm{W}_o$投影，堆叠若干层后，每个节点表示与文档的CLS表示拼接，
做因果与否的二分类。由于因果对远少于非因果对，ERGO用$\alpha$平衡的\term{焦点损失}（focal loss）
\begin{equation}
  L_{\mathrm{FL}}=-\sum_{(i,j)}\alpha_{i,j}(1-p_{i,j})^{\gamma}\log p_{i,j},
  \label{eq:nlp-extraction-focal}
\end{equation}
其中$p_{i,j}$是模型给真实类别的概率，$\gamma\geq0$；判得越准的样本权重$(1-p)^{\gamma}$越小，
训练集中到难分的样本上，$\gamma=0$时退化为加权交叉熵。

事件对节点的代价是规模：$m$个事件有$O(m^2)$个节点、约$O(m^3)$条边，长文档需要剪枝候选对。
外生因素可以切断因果链条，因此ERGO只把传递性作为图结构的归纳偏置，而非硬约束。

\subsection{组成与子事件关系}
\label{sec:nlp-extraction-subevents}

补充教学安排“本次会议由专题报告和分组讨论两个环节组成”，支持报告、讨论是会议的子事件。
它不同于事件共指（报告不是整场会议），也不同于类型上下位（“学术会议是一类会议”谈的是类别）；
午休即使在会议时间范围内也未必是会议的一部分，时间包含不是充分条件。
挑战是父子方向、层次的无环性，以及子事件关系与共指、时间关系的相互约束。
方法先做层次关系分类再施加无环约束，进一步把这些约束与其他事件关系联合学习。

\subsubsection{层次关系分类与无环约束}
\label{sec:nlp-extraction-subevent-acyclic}

HiEve把新闻中的事件层次标为父子、子父、共指与无关四类\scite{nlp-extraction-glavas2014hieve}
给定事件对$(p,c)$，把两者锚点、参与者及上下文编码为$\bm{g}_{p,c}$，
用线性或双仿射头在这四类上分类，以有向实例对的交叉熵训练。
若所选定义要求组成关系无环，可对选边变量$z_{p,c}$作整数优化：
最大化$\sum_{p,c}w_{p,c}z_{p,c}$，并用拓扑次序变量$u_p$约束
\begin{equation}
  u_p+1\leq u_c+M(1-z_{p,c}),\qquad
  0\leq u_p\leq N-1,\quad M=N,
  \label{eq:nlp-extraction-subevent-ilp}
\end{equation}
其中$N$是事件节点数。选中边时强制父在子之前；这里的次序是图的拓扑次序，不是时间先后。
允许一个环节属于多个上层活动时，结构是有向无环图，而非单父树。

教学候选为会议$\to$报告得分3.0、报告$\to$讨论得分2.0、讨论$\to$会议得分0.2。
局部阈值0.1会选出环，无环最大权子图舍弃0.2的边，保留总分5.0。
但原文把报告和讨论列为并列环节，“报告$\to$讨论”缺乏语义支持：
去环只清除结构矛盾，最终仍应只保留会议到两个环节的边。

整数规划只在推断时起作用。Wang等把时间与子事件关系的逻辑约束改写成可微的训练损失
\scite{nlp-extraction-wang2020jcl}
记模型给关系$\alpha$的概率为$\alpha(e_1,e_2)$。对称性要求一个关系与其逆关系同真同假，
损失为$|\log\alpha(e_1,e_2)-\log\bar{\alpha}(e_2,e_1)|$；
合取规则“$\alpha(e_1,e_2)\wedge\beta(e_2,e_3)\to\gamma(e_1,e_3)$”在乘积t-范数下转为
\begin{equation}
  L_{t}=\bigl|\log\alpha(e_1,e_2)+\log\beta(e_2,e_3)-\log\gamma(e_1,e_3)\bigr|,
  \label{eq:nlp-extraction-transitivity-loss}
\end{equation}
“推出$\delta$不成立”的规则把$\log\gamma$换成$\log(1-\delta)$。
例如“甲先于乙”与“乙是丙的父事件”推出“甲先于丙”，并排除“甲是丙的子事件”，
这条规则同时涉及时间与子事件两类关系。总损失为标注交叉熵加
$\lambda_S$倍的对称损失与$\lambda_C$倍的合取损失。
与算法~\ref{alg:nlp-extraction-temporal-structured}相比，这里不在训练中求解整数规划，
而是把每条规则写成损失项，梯度直接把违反约束的概率组合推开；
约束不再严格成立，测试时仍可再加一次整数线性规划。

\subsubsection{多类事件关系的联合抽取}
\label{sec:nlp-extraction-joint-event-relations}

MAVEN-ERE的标注本身体现了类型间的约束：因果关系只在已标为先于或重叠的事件对上标注，
子事件关系只在已标为包含的事件对上标注，共指的事件共享全部关系。
其时间关系采用先于、包含、重叠、同时开始、同时结束、同时六类。
基线模型用RoBERTa编码整篇文档，取每个事件触发位置的上下文表示，
对事件对设分类头；四个任务可以分别训练，也可以共享编码器、把四个损失相加联合训练
\scite{wang2022mavenere}
原文报告联合训练对数据较少的因果与子事件任务帮助较明显。

联合训练只靠共享参数隐式传递约束，可微逻辑约束则把规则显式写成损失。
需要严格一致时，仍可在四类局部分数上联合求解整数线性规划：
共指簇内事件取相同的对外关系，子事件边蕴含包含，因果边蕴含先于或重叠。

\begin{table*}[htbp]
  \centering
  \small
  \begin{tabularx}{\linewidth}{@{}>{\raggedright\arraybackslash}p{24mm} XXX@{}}
    \toprule
    方法 & 解决的问题 & 核心机制 & 新的代价或局限\\
    \midrule
    HeidelTime\slash SUTime & 时间表达归一化 & 模式规则＋参照时间的日历计算 & 依赖语言资源；参照选择易错 \\
    事件对分类 & 局部时间、因果、组成判断 & 事件对表示上的有向分类头 & 结果可能自相矛盾 \\
    CAEVO & 多个分类器的一致组合 & 按精确率排序的筛子＋逐步传递闭包 & 早期错误被传播；次序需人工调定 \\
    Allen约束与ILP & 局部结果不一致 & 区间组合传播；带对称、传递约束的整数规划 & 只在推断时用约束；求解代价 \\
    结构化学习 & 训练不知道约束 & 训练循环内做约束推断的感知机更新 & 每步求解一次整数规划 \\
    时间--因果联合ILP & 因果与时间互相矛盾 & 原因先于结果的跨类约束 & 约束对持续过程需放宽 \\
    KnowDis\slash LearnDA & 因果标注稀少 & 词汇知识扩充事件对；常识过滤或对偶生成 & 增广数据带噪；过滤阈值需调 \\
    ERGO & 跨句因果缺少局部线索 & 事件对为节点的关系图Transformer＋焦点损失 & 节点$O(m^2)$；需剪枝 \\
    无环ILP & 组成关系成环 & 拓扑次序变量约束选边 & 去环不等于有证据 \\
    可微逻辑约束 & 约束只在推断时生效 & 对称与合取规则改写为对数空间损失 & 约束只被软性满足 \\
    MAVEN-ERE联合 & 各类关系分开训练 & 共享文档编码器，四个关系头损失相加 & 类型间约束只隐式传递 \\
    \bottomrule
  \end{tabularx}
  \caption{事件关系抽取方法对照。前四行处理时间关系，其后依次为约束感知学习、因果关系
  和组成关系；“新的代价”指相对于逐对局部分类新增的计算、依赖或局限。}
  \label{tab:nlp-extraction-event-relation-methods}
\end{table*}

\section{知识记录的构建}
\label{sec:nlp-extraction-records}

抽取系统可能从每份通知都正确读出了字段，汇总后仍然出错：把两场活动误并，
把一条转发当成新的独立证据，或用旧通知覆盖更正后的日期。
知识构建因此不只是把三元组装进图里，还要回答四个问题：哪些记录指同一对象（对齐）；
指同一对象的记录在值上冲突时信哪一个（合并与冲突处理）；图中缺失的事实能否由已有结构推测
（知识图谱补全）；新材料到来时改变了哪一条已有主张（核验与更新）。
贯穿的例子是同一活动的更正：旧排期与新排期是同一活动的两个版本，不是两场活动。

图~\ref{fig:nlp-extraction-records}把本章各层产出连成处理流程。
其中旧通知是额外的教学构造，现通知仍是章首2025年4月15日的材料。
\begin{figure*}[htbp]
  \centering
  % Chapter-local TikZ diagram; final width about 168 mm, no external assets.
% Solid arrows carry records. The dashed return arrow requests evidence checks.
\begin{tikzpicture}[
  x=1mm,y=1mm,
  font=\figurefont\fontsize{9}{11}\selectfont,
  box/.style={rectangle,align=left,inner sep=2mm,line width=0.85pt,
    draw=BookRule,fill=white,text width=42mm,minimum height=19mm},
  input/.style={box,draw=FigureInput,fill=FigureInput!8},
  model/.style={box,draw=FigureModel,fill=FigureModel!8},
  output/.style={box,draw=FigureOutput,fill=FigureOutput!8},
  flow/.style={line width=1pt,draw=BookInk,-{Latex[length=2mm]}},
  feedback/.style={line width=0.85pt,draw=FigureModel,dashed,
    -{Latex[length=2mm]}},
  heading/.style={font=\figurefont\bfseries\fontsize{9}{11}\selectfont,
    text=BookInk,anchor=west}
]
  \node[heading] at (0,17) {A\quad 每条事实保留可回查的依据};
  \node[input] (text) at (23,0) {
    \textbf{原始文本}\\
    文档ID、版本、发布日期\\
    保留未经覆盖的原文
  };
  \node[input] (mention) at (84,0) {
    \textbf{提及定位}\\
    字符跨度、类型\\
    词元到原文的映射
  };
  \node[model] (binding) at (145,0) {
    \textbf{对象绑定}\\
    共指簇、实体身份\\
    活动与角色归属
  };
  \node[model] (candidate) at (145,-34) {
    \textbf{关系或事件候选}\\
    对象、角色、极性\\
    时间、条件、证据句
  };
  \node[model] (check) at (84,-34) {
    \textbf{记录核验}\\
    模式、身份、原文支持\\
    限定条件与来源
  };
  \node[output] (record) at (23,-34) {
    \textbf{带来源记录}\\
    规范键、字段与值\\
    有效时间、写入时间
  };
  \draw[flow] (text.east) -- (mention.west);
  \draw[flow] (mention.east) -- (binding.west);
  \draw[flow] (binding.south) -- (candidate.north);
  \draw[flow] (candidate.west) -- (check.east);
  \draw[flow] (check.west) -- (record.east);
  \draw[feedback] (check.north) -- node[right,text=FigureModel] {回查} (mention.south);
  \draw[BookRule,line width=0.7pt] (0,-52) -- (168,-52);
  \node[heading] at (0,-60) {B\quad 同一活动改期：保留旧值，不把两次主张当成两场活动};
  \node[input,text width=75mm,minimum height=21mm] (old) at (40,-80) {
    \textbf{旧通知：4月14日发布}\\
    活动编号 ML-0418\\
    原排期为4月17日，未给时刻
  };
  \node[input,text width=75mm,minimum height=21mm] (new) at (127,-80) {
    \textbf{更正通知：4月15日发布}\\
    同一活动改至4月18日14:00\\
    “原定”与“现改至”连接新旧值
  };
  \node[output,text width=163mm,minimum height=25mm] (versions) at (84,-119) {
    \textbf{同一活动键，两个排期版本（年份均为2025年）}\\
    旧值：4月17日；排期主张有效期 $[\text{4月14日},\text{4月15日})$；保留旧来源\\
    新值：4月18日14:00；排期主张自4月15日起有效；关联更正来源\\
    系统若4月16日才读到更正，写入日为4月16日，有效起日仍为4月15日
  };
  \draw[BookInk,line width=1pt] (old.south) -- (40,-98) -- (84,-98);
  \draw[BookInk,line width=1pt] (new.south) -- (127,-98) -- (84,-98);
  \draw[flow] (84,-98) -- (versions.north);
\end{tikzpicture}

  \caption{从文本证据到知识记录。上部保存定位、身份、关系或事件以及来源字段，
  核验失败后返回相关定位或绑定步骤。下部两条通知均为教学构造；
  旧通知发布于4月14日，新通知发布于4月15日，本例约定更正从发布日生效。
  排期主张的有效期与活动开始时刻分开记录，系统写入日期也不决定事实的新旧。}
  \label{fig:nlp-extraction-records}
\end{figure*}

\subsection{记录对齐}
\label{sec:nlp-extraction-alignment}

对齐先解决“字段说的是不是同一种东西”，再解决“记录说的是不是同一个对象”。
“举行时间”与“start”可以通过模式映射对应，“报名截止”却不能因为同为日期就映射到start。
单位、时区、粒度与关系方向也要规范化并保留原值：“14:00”与“下午2点”在同一日、同一时区可对应，
“4月18日”与“4月18日14:00”则精度不同。

身份比较的挑战有三项。候选对数量为两个来源记录数之积，需要\term{阻塞}（blocking）：
按可靠编号、日期块或名称词项只比较落入同一块的记录，代价是可能漏掉别名与改期记录，
候选覆盖要单独评价。字段格式异构，同一信息在两边可能分布在不同字段中。
对应关系也未必一对一，一条全年汇总记录可能对应多次活动。
方法谱系依次是：Fellegi--Sunter用概率模型给出无监督的决策规则；
监督匹配用分类器学习字段交互，并以受约束分配保证一对一；
基于预训练模型的匹配把记录序列化后直接判断，减少字段对齐与特征设计。

\subsubsection{Fellegi--Sunter记录匹配}
\label{sec:nlp-extraction-fellegi-sunter}

Fellegi--Sunter把两条记录是否同指视作匹配$M$与非匹配$U$之间的判别，
以比较向量在两种条件下出现的概率形成证据\scite{fellegi1969linkage}
比较向量$\bm{\gamma}$由逐字段的比较结果构成，例如名称编辑距离的分档、地点与日期是否一致。
核心量是似然比；若给定匹配状态时各字段的比较结果条件独立，对数似然比可以分项相加：
\begin{equation}
  \Lambda(\bm{\gamma})=
  \frac{P(\bm{\gamma}\mid M)}{P(\bm{\gamma}\mid U)},
  \qquad
  \log\Lambda(\bm{\gamma})
    =\sum_{f\in O}
       \log\frac{m_f(\gamma_f)}{u_f(\gamma_f)}.
  \label{eq:nlp-extraction-fellegi}
\end{equation}
$O$为可比较的字段，$m_f,u_f$可由已核验的记录对估计，也可在混合模型下用EM从未标数据估计。
决策规则设两个阈值：高于上阈值接受，低于下阈值拒绝，其间交人工复核。
名称与简称高度相关时，分项相加会重复计分。

教学参数取名称相似的$m/u=0.9/0.1=9$，地点一致的$m/u=0.8/0.2=4$。
比较“机器学习读书会”与别名表中的“ML读书会”，地点相同而日期缺失，
缺失字段贡献1（假定缺失本身不携带身份信息），则$\Lambda=9\times4=36$。
规定$\Lambda<3$拒绝、$\Lambda>20$接受，则36接受；若日期确实冲突，乘上日期不一致的似然比
$0.05/0.95=1/19$，总比约1.89而拒绝。似然比不是后验概率：给定先验匹配率$\pi$，
后验为$\pi\Lambda/(\pi\Lambda+1-\pi)$。改期不构成身份不同的证据，比较前应先判明是原定日期还是现改日期。

\subsubsection{监督匹配与受约束分配}
\label{sec:nlp-extraction-supervised-matching}

Fellegi--Sunter的字段独立假设无法表达字段交互，例如“名称相近但学校不同”通常不是同一活动。
监督分类器输入规范字段、缺失掩码与相似度特征$\bm{\phi}_{i,j}$，输出
$p_{i,j}=\sigma(f_{\bm{\theta}}(\bm{\phi}_{i,j}))$，以已核验同指对与非同指对的二元交叉熵训练。
负例既要有随机的不同记录，也要有同名、同场所的难负例；
训练与测试按对象或时间划分，避免同一活动的不同转发跨集合泄漏。
DeepMatcher等神经匹配器把每个对齐字段的取值编码为向量，逐字段比较后再汇总分类
\scite{nlp-extraction-mudgal2018deepmatcher}
它们省去了手工相似度特征，但仍要求两边字段先对齐。

逐对概率还不够：若两个名单都“一条记录对应一次活动”，多条记录可能同时选中同一个目标。
在候选边效用$s_{i,j}$（例如校准后的匹配对数几率减去接受成本）上求
\begin{equation}
  \max_{\bm{x}}\sum_{i,j}s_{i,j}x_{i,j},
  \qquad
  \sum_jx_{i,j}\leq1,\quad
  \sum_ix_{i,j}\leq1,\quad
  x_{i,j}\in\{0,1\},
  \label{eq:nlp-extraction-assignment}
\end{equation}
未召回的边禁止选择，允许不匹配时加虚拟节点。这是二分图上的分配问题，
可用匈牙利算法在多项式时间内精确求解\scite{nlp-extraction-kuhn1955hungarian}

教学效用矩阵为
$\bm{S}=\left(\begin{smallmatrix}0.90&0.80\\0.85&0.20\end{smallmatrix}\right)$，要求完整一对一分配。
两行的局部首选都是第一列；对角分配总分1.10，交叉分配总分1.65。
匈牙利算法把效用换成成本$\bm{C}=\bm{1}-\bm{S}$，逐行、逐列减去最小值，
这不改变完整分配的最优解：
\begin{equation}
  \begin{pmatrix}0.10&0.20\\0.15&0.80\end{pmatrix}
  \ \longrightarrow\
  \begin{pmatrix}0&0.10\\0&0.65\end{pmatrix}
  \ \longrightarrow\
  \begin{pmatrix}0&0\\0&0.55\end{pmatrix}.
  \label{eq:nlp-extraction-hungarian}
\end{equation}
此时可选不同行列的零位$(1,2),(2,1)$，即效用1.65的交叉分配。
一般情形下零位不足以构成完整匹配时，算法用最少的行列覆盖零位、调整未覆盖元素后沿交替路径增广。
汇总记录对应多期活动时，一对一假设不成立，不应强制分配。

\subsubsection{基于预训练模型的实体匹配}
\label{sec:nlp-extraction-ditto}

前两种方法都依赖字段对齐，而真实来源的模式常常不一致：一边把时间写在“start”里，
另一边写在描述文本中。\term{实体匹配}（entity matching）的预训练模型方法Ditto
把记录对的判断改写为句对分类，由预训练编码器直接比较两条序列化后的记录
\scite{nlp-extraction-li2020ditto}
记录$e=\{(\mathrm{attr}_k,\mathrm{val}_k)\}_{k=1}^{K}$序列化为
\begin{equation}
  \begin{aligned}
  \mathrm{ser}(e)&=\texttt{[COL]}\ \mathrm{attr}_1\ \texttt{[VAL]}\ \mathrm{val}_1\ \cdots\
                  \texttt{[COL]}\ \mathrm{attr}_K\ \texttt{[VAL]}\ \mathrm{val}_K,\\
  \mathrm{ser}(e,e')&=\texttt{[CLS]}\ \mathrm{ser}(e)\ \texttt{[SEP]}\ \mathrm{ser}(e')\ \texttt{[SEP]},
  \end{aligned}
  \label{eq:nlp-extraction-ditto-serialize}
\end{equation}
\texttt{[COL]}与\texttt{[VAL]}标记字段名和值的开始。拼接后的序列送入BERT类编码器，
CLS位置的向量接全连接层做二分类，以交叉熵微调。
两条记录的字段名不必相同、数量也不必相同，自注意力会在两段之间自行对齐相关的词元。
教学记录A为“\texttt{[COL]}名称\texttt{[VAL]}机器学习读书会\texttt{[COL]}时间\texttt{[VAL]}2025-04-18 14:00
\texttt{[COL]}地点\texttt{[VAL]}海棠楼201”，记录B为“\texttt{[COL]}title\texttt{[VAL]}ML读书会
\texttt{[COL]}start\texttt{[VAL]}4月18日下午2点\texttt{[COL]}room\texttt{[VAL]}201”，
二者无需先建立“时间—start”的映射即可拼成一个输入。

Ditto在基本做法上加了三项技术。\emph{领域知识注入}有两种：
跨度类型标注用命名实体识别或正则表达式找出产品编号、门牌号、电话末四位等关键片段，
在两侧加上类型标记，提示模型对齐同类片段而不把门牌号与年份比较；
跨度规范化把写法不同但等价的片段改写成同一字符串，例如把“下午2点”改写为“14:00”、
把数字统一保留两位小数。\emph{长字段摘要}：编码器输入长度有限，直接截断会丢掉位于后部的关键信息，
Ditto按TF-IDF只保留得分高的非停用词。\emph{数据增强}：对训练样本做删除或打乱一段词元、
删除或打乱字段、交换两条记录的次序等操作，生成更难的样本；
为避免操作改变标签（例如删掉了决定匹配的字段），Ditto不直接用增强样本，
而采用MixDA在表示层插值：
\begin{equation}
  \lambda\sim\mathrm{Beta}(a,a),\qquad
  \bm{h}''=\lambda\,\mathrm{Enc}(s)+(1-\lambda)\,\mathrm{Enc}(\mathrm{aug}(s)),
  \label{eq:nlp-extraction-mixda}
\end{equation}
$\bm{h}''$代替原表示送入分类层，标签沿用原样本，原文取$a=0.8$。

Ditto省去了字段对齐与手工特征，代价是每个候选对一次完整的编码器前向，只适合阻塞后的候选；
领域知识规则仍需人工编写。表~\ref{tab:nlp-extraction-alignment-methods}对比三种方法。

\begin{table}[htbp]
  \centering
  \small
  \begin{tabularx}{\linewidth}{@{}>{\raggedright\arraybackslash}p{16mm} XXX@{}}
    \toprule
    方法 & 监督需求 & 字段对齐 & 主要局限\\
    \midrule
    Fellegi--Sunter & 可无监督（EM） & 需要 & 字段条件独立；相关字段重复计分 \\
    监督匹配＋分配 & 已核验记录对 & 需要 & 特征工程；一对一假设 \\
    Ditto & 已核验记录对 & 不需要 & 每对一次编码器前向；规则需人工 \\
    \bottomrule
  \end{tabularx}
  \caption{记录对齐方法对照。三者都只在阻塞后的候选对上运行，候选覆盖需另行评价。}
  \label{tab:nlp-extraction-alignment-methods}
\end{table}

\subsection{记录合并与冲突处理}
\label{sec:nlp-extraction-merge}

对齐后，完全相同的事实按“身份、字段、值、条件、有效时间”的规范键去重，并合并来源列表；
近似重复的记录先用受约束的凝聚聚类确认身份，再合并互补字段。
并查集适合维护已确认的等价关系，却不适合对任意高相似边直接取闭包：
“同名且同地点”可能串起多周的活动。

冲突值不能总靠多数票解决。对同一对象、同一字段、同一适用条件下的类别值，
来源加权投票为$\widehat{v}=\argmax_v\sum_jw_j\mathbb{I}[y_j=v]$，
权重应由来源核验或开发集估计。转载同一通知的十个页面并非十次独立观察，需要记录依赖来源；
日期更正则是时间上的状态改变，应先拆开适用期，而不是把新旧排期交给投票。

当同类记录较多、来源有各自的系统误差时，可以用Dawid--Skene模型同时估计潜在真值与来源的
混淆矩阵\scite{dawid1979observer}
设记录$i$的潜在值$z_i\in\{1,\ldots,K\}$，先验$\pi_k$，来源$j$的观测为$y_{i,j}$，
混淆矩阵$C^{(j)}_{k,\ell}=P(y_{i,j}=\ell\mid z_i=k)$。在给定真值后来源条件独立的假设下，
EM的E步与M步为
\begin{equation}
  q_{i,k}=
  \frac{\pi_k\prod_{j\in O_i}C^{(j)}_{k,y_{i,j}}}
       {\sum_{k'}\pi_{k'}\prod_{j\in O_i}C^{(j)}_{k',y_{i,j}}},
  \label{eq:nlp-extraction-ds-e}
\end{equation}
\begin{equation}
  \pi_k^{\mathrm{new}}=\frac{1}{N}\sum_iq_{i,k},
  \qquad
  C^{(j),\mathrm{new}}_{k,\ell}=
  \frac{\sum_{i:j\in O_i}q_{i,k}\mathbb{I}[y_{i,j}=\ell]}
       {\sum_{i:j\in O_i}q_{i,k}},
  \label{eq:nlp-extraction-ds-m}
\end{equation}
$O_i$只含实际给出观测的来源，缺失不当作反对票。
该模型的可识别条件、类别置换与来源相关时的修正，可参考
\BookChapterRef{chap:knowledge-fusion}{卷八的知识融合章}；这里只用到上面两步更新及其独立性前提。

\begin{example}[一次来源后验]
\label{ex:nlp-extraction-ds}
两条二分类记录，先验各0.5，两个来源的对称准确率为0.9与0.8。记录1收到来源1的A与来源2的B，
A、B的未归一化权重为$0.5\times0.9\times0.2=0.09$与$0.5\times0.1\times0.8=0.04$，
故$q_{1,A}=9/13\approx0.692$。若只有这两条记录，M步会把只报过B的来源2的混淆矩阵推到0或1，
实际应加Dirichlet平滑。
\end{example}

转载违反条件独立，此时后验会高估置信度；无法解释的冲突应并存保存并待核验。

\subsection{知识图谱嵌入与补全}
\label{sec:nlp-extraction-kge}

合并后的记录可以写成\term{知识图谱}（knowledge graph）：实体集合$E$、关系集合$R$
与三元组集合$\mathcal{G}\subseteq E\times R\times E$，三元组$(h,r,t)$表示头实体$h$与尾实体$t$
之间有关系$r$。抽取得到的图总是不完整的：通知说了读书会在海棠楼201举行、海棠楼201属于海棠大学，
却未必写出“读书会由海棠大学承办”。\term{知识图谱补全}（knowledge graph completion）
根据已有三元组推测缺失的三元组，常用的任务形式是\term{链接预测}（link prediction）：
给定$(h,r,?)$或$(?,r,t)$，对全部候选实体排序。

\term{知识图谱嵌入}（knowledge graph embedding）把每个实体和关系表示为向量，
用打分函数$f(h,r,t)$衡量三元组成立的可能性，分数越高越可能成立。
补全的挑战来自关系的结构：一对多、多对一关系（一所大学有多栋楼）
要求同一个$(h,r)$对应多个尾实体；关系还有对称（“同场举行”）、反对称（“先于”）、
互逆（“承办”与“由……承办”）、组合（“位于”接“隶属”）等模式，
不同的打分函数能表达的模式不同。下面先讲训练与评价的共同框架，
再讲平移距离与双线性两类打分函数。

训练只有正例，负例靠\term{负采样}（negative sampling）构造：对正例$(h,r,t)$，
随机替换头或尾得到$(h',r,t)$或$(h,r,t')$。被替换出的三元组可能恰好是图中的另一个真三元组，
这类假负例在一对多关系上尤其多。两种常用目标是\term{间隔排序损失}（margin ranking loss）
与逻辑斯谛损失：
\begin{equation}
  \begin{aligned}
  L_{\mathrm{margin}}&=\sum_{(h,r,t)\in\mathcal{G}}\ \sum_{(h',r,t')\in\mathcal{N}(h,r,t)}
  \bigl[\gamma-f(h,r,t)+f(h',r,t')\bigr]_+,\\
  L_{\mathrm{log}}&=\sum_{(h,r,t)}\log\bigl(1+\exp(-y_{h,r,t}\,f(h,r,t))\bigr),
  \end{aligned}
  \label{eq:nlp-extraction-kge-margin}
\end{equation}
其中$\mathcal{N}(h,r,t)$是该正例的负例集合，$[\cdot]_+=\max(0,\cdot)$，
$\gamma>0$是间隔，$y\in\{+1,-1\}$标记正负例。间隔损失只要求正例比负例高出$\gamma$，
超过之后不再产生梯度；逻辑斯谛损失则把每个三元组当作独立的二分类样本。

评价按排名进行。对测试三元组$(h,r,t)$，用模型对所有$(h,r,e)$打分，
记真实尾实体$t$的名次为$\mathrm{rank}$。\emph{过滤设定}在排序前删去训练、验证、测试集中
其他已知为真的$(h,r,e)$，避免模型因把另一个正确答案排在前面而被扣分。
在查询集合$Q$上，\term{平均倒数排名}（mean reciprocal rank, MRR）与Hits@$k$为
\begin{equation}
  \mathrm{MRR}=\frac{1}{|Q|}\sum_{q\in Q}\frac{1}{\mathrm{rank}_q},
  \qquad
  \mathrm{Hits@}k=\frac{1}{|Q|}\sum_{q\in Q}\mathbb{I}[\mathrm{rank}_q\leq k].
  \label{eq:nlp-extraction-mrr}
\end{equation}
教学示例有三个查询，原始名次为1、4、12；第二个查询中排在前面的三个实体里有两个也是正确答案，
过滤后名次为2。原始MRR为$(1+1/4+1/12)/3\approx0.444$，过滤MRR为$(1+1/2+1/12)/3\approx0.528$；
Hits@3从$1/3$升为$2/3$。头实体与尾实体两个方向通常都要预测并取平均。
数据集构造也会影响结论：FB15k中大量测试三元组的互逆三元组出现在训练集里，
简单规则即可得高分，FB15k-237删去了这类互逆关系\scite{nlp-extraction-toutanova2015observed}

\subsubsection{平移距离模型}
\label{sec:nlp-extraction-translational}

TransE把关系看作嵌入空间中的平移：若$(h,r,t)$成立，应有$\bm{h}+\bm{r}\approx\bm{t}$
\scite{nlp-extraction-bordes2013transe}
打分函数为
\begin{equation}
  f(h,r,t)=-\lVert\bm{h}+\bm{r}-\bm{t}\rVert_{p},\qquad p\in\{1,2\},
  \label{eq:nlp-extraction-transe}
\end{equation}
$\bm{h},\bm{r},\bm{t}\in\mathbb{R}^{d}$。训练用式~\eqref{eq:nlp-extraction-kge-margin}的间隔损失，
每个正例替换头或尾之一得到负例，并在每批更新前把实体向量归一化为单位长度，
防止模型靠放大实体范数来降低损失，过程见算法~\ref{alg:nlp-extraction-transe}。
每个实体和关系只有$d$个参数，模型规模随实体数与关系数线性增长。

\Needspace{12\baselineskip}
\begin{algorithm}[TransE的随机梯度训练]
  \label{alg:nlp-extraction-transe}
  \begin{algorithmic}[1]
    \Require 三元组集合$\mathcal{G}$，实体集合$E$，关系集合$R$，维度$d$，间隔$\gamma$，学习率$\eta$
    \Ensure 实体与关系嵌入
    \State 对每个$r\in R$随机初始化$\bm{r}$并归一化；对每个$e\in E$随机初始化$\bm{e}$
    \Repeat
      \State 对每个$e\in E$：$\bm{e}\gets\bm{e}/\lVert\bm{e}\rVert_2$
      \State 从$\mathcal{G}$采样一批正例$B$
      \For{每个$(h,r,t)\in B$}
        \State 以等概率替换$h$或$t$为随机实体，得到负例$(h',r,t')$
        \State 累加损失$\bigl[\gamma+\lVert\bm{h}+\bm{r}-\bm{t}\rVert-\lVert\bm{h}'+\bm{r}-\bm{t}'\rVert\bigr]_+$
      \EndFor
      \State 对涉及的实体与关系嵌入做一步梯度下降，学习率$\eta$
    \Until{验证集上的过滤MRR不再提高}
  \end{algorithmic}
\end{algorithm}

二维教学示例：$\bm{h}=(0,0)$，$\bm{r}=(1,0)$，$\bm{t}=(1,0.2)$，正例距离为0.2。
负例$\bm{t}'=(0,1)$的距离为$\lVert(1,-1)\rVert_2\approx1.414$，取$\gamma=1$时
$[1+0.2-1.414]_+=0$，这个负例已足够远而不产生梯度；
负例$\bm{t}''=(1,0.5)$的距离为0.5，损失为$1+0.2-0.5=0.7$，梯度把$\bm{t}''$推离$\bm{h}+\bm{r}$，
同时把$\bm{t}$拉近。

平移假设决定了TransE能表达什么。若$r$是一对多关系，$(h,r,t_1)$与$(h,r,t_2)$都完美成立，
则$\bm{t}_1=\bm{h}+\bm{r}=\bm{t}_2$，两个不同的尾实体被迫重合；若$r$对称，
由$\bm{h}+\bm{r}=\bm{t}$与$\bm{t}+\bm{r}=\bm{h}$得$\bm{r}=\bm{0}$且$\bm{h}=\bm{t}$。
另一方面，互逆关系取$\bm{r}'=-\bm{r}$即可，组合关系满足$\bm{r}_3=\bm{r}_1+\bm{r}_2$，
这两类模式TransE能自然表达。

TransH让实体在每个关系上有不同的投影\scite{nlp-extraction-wang2014transh}
每个关系有一个超平面法向量$\bm{w}_r$（$\lVert\bm{w}_r\rVert_2=1$）与平移向量$\bm{d}_r$，
\begin{equation}
  \bm{h}_{\perp}=\bm{h}-(\bm{w}_r^{\top}\bm{h})\bm{w}_r,\qquad
  \bm{t}_{\perp}=\bm{t}-(\bm{w}_r^{\top}\bm{t})\bm{w}_r,\qquad
  f(h,r,t)=-\lVert\bm{h}_{\perp}+\bm{d}_r-\bm{t}_{\perp}\rVert_2^2.
  \label{eq:nlp-extraction-transh}
\end{equation}
一对多关系的多个尾实体只需投影重合，在法向量方向上仍可不同；
TransH还按关系的“每头平均尾数”与“每尾平均头数”决定替换头或尾，以减少假负例。TransR进一步给每个关系一个投影矩阵
$\bm{M}_r\in\mathbb{R}^{k\times d}$，在关系空间中平移\scite{nlp-extraction-lin2015transr}
\begin{equation}
  f(h,r,t)=-\lVert\bm{M}_r\bm{h}+\bm{r}-\bm{M}_r\bm{t}\rVert_2^2,\qquad \bm{r}\in\mathbb{R}^{k},
  \label{eq:nlp-extraction-transr}
\end{equation}
每个关系的参数从$O(d)$增至$O(kd)$，表达力提高，关系多时训练代价也随之增加。
RotatE把实体与关系放在复数空间，令$\bm{t}=\bm{h}\circ\bm{r}$且每个分量$|r_i|=1$，
关系是逐分量的旋转\scite{nlp-extraction-sun2019rotate}
旋转角取0或$\pi$时关系对称，取相反角得到互逆，角度相加得到组合，
于是对称、反对称、互逆与组合都能表达；它还按当前模型对负例的打分加权负例，使训练集中于难负例。

\subsubsection{双线性语义匹配模型}
\label{sec:nlp-extraction-bilinear}

平移模型比较距离，双线性模型比较语义相似度。RESCAL为每个关系配一个满矩阵
\scite{nlp-extraction-nickel2011rescal}
$f(h,r,t)=\bm{h}^{\top}\bm{M}_r\bm{t}$，$\bm{M}_r\in\mathbb{R}^{d\times d}$
可以表达头尾每一对维度之间的交互，但每个关系需要$d^2$个参数，关系多时容易过拟合。
DistMult把$\bm{M}_r$限制为对角矩阵\scite{nlp-extraction-yang2015distmult}
\begin{equation}
  f(h,r,t)=\bm{h}^{\top}\operatorname{diag}(\bm{r})\,\bm{t}=\sum_{i=1}^{d}h_ir_it_i,
  \label{eq:nlp-extraction-distmult}
\end{equation}
参数降回$O(d)$，但三个向量的逐元素乘积对$h$与$t$对称，$f(h,r,t)=f(t,r,h)$，
因此只能表示对称关系，无法区分“甲先于乙”与“乙先于甲”。
ComplEx把嵌入放到复数域$\mathbb{C}^{d}$，对尾实体取共轭
\scite{nlp-extraction-trouillon2016complex}
\begin{equation}
  f(h,r,t)=\operatorname{Re}\Bigl(\sum_{i=1}^{d}h_ir_i\overline{t_i}\Bigr),
  \label{eq:nlp-extraction-complex}
\end{equation}
交换$h$与$t$后，$\operatorname{Re}(t_ir_i\overline{h_i})=\operatorname{Re}(h_i\overline{r_i}\,\overline{t_i})$，
即交换头尾等价于把关系向量取共轭。因此关系向量的实部贡献对称部分、虚部贡献反对称部分，
实部为零时$f(t,r,h)=-f(h,r,t)$。ComplEx用式~\eqref{eq:nlp-extraction-kge-margin}的逻辑斯谛损失
加L2正则训练，参数量是DistMult的两倍。

用最小的例子核对三者对互逆三元组的分数。RESCAL取$d=2$，
$\bm{M}_r=\left(\begin{smallmatrix}0&1\\-1&0\end{smallmatrix}\right)$，
$\bm{h}=(1,0)$，$\bm{t}=(0,1)$：$f(h,r,t)=\bm{h}^{\top}\bm{M}_r\bm{t}=1$，
$f(t,r,h)=\bm{t}^{\top}\bm{M}_r\bm{h}=-1$，反对称矩阵表达了方向。
DistMult无论取什么对角元素，都有$f(h,r,t)=f(t,r,h)$。
ComplEx取$d=1$，$h=1$，$t=\mathrm{i}$，$r=\mathrm{i}$：
$f(h,r,t)=\operatorname{Re}(1\cdot\mathrm{i}\cdot\overline{\mathrm{i}})=\operatorname{Re}(1)=1$，
$f(t,r,h)=\operatorname{Re}(\mathrm{i}\cdot\mathrm{i}\cdot1)=-1$，
只用$O(d)$个参数就恢复了RESCAL例子中的方向性。表~\ref{tab:nlp-extraction-kge-models}汇总各模型。

\begin{table*}[htbp]
  \centering
  \small
  \begin{tabularx}{\linewidth}{@{}l>{\raggedright\arraybackslash}p{38mm}>{\raggedright\arraybackslash}p{20mm}XXX@{}}
    \toprule
    模型 & 打分函数 & 每关系参数 & 对称 & 反对称与互逆 & 一对多 \\
    \midrule
    TransE & $-\lVert\bm{h}+\bm{r}-\bm{t}\rVert$ & $d$ & 否（需$\bm{r}=\bm{0}$） & 能 & 尾实体被迫重合 \\
    TransH & 超平面投影后平移 & $2d$ & 否 & 能 & 投影重合即可 \\
    TransR & 关系空间投影后平移 & $kd+k$ & 否 & 能 & 投影重合即可 \\
    RotatE & $-\lVert\bm{h}\circ\bm{r}-\bm{t}\rVert$，$|r_i|=1$ & $d$（复数） & 能 & 能 & 仍受限 \\
    RESCAL & $\bm{h}^{\top}\bm{M}_r\bm{t}$ & $d^2$ & 能 & 能 & 能 \\
    DistMult & $\sum_ih_ir_it_i$ & $d$ & 能 & 否 & 能 \\
    ComplEx & $\operatorname{Re}\sum_ih_ir_i\overline{t_i}$ & $d$（复数） & 能 & 能 & 能 \\
    \bottomrule
  \end{tabularx}
  \caption{知识图谱嵌入模型对照。“能”指打分函数存在参数取值表达该模式，不表示训练一定学到；
  组合关系在平移与旋转模型中由向量相加或角度相加表达，双线性模型对此没有对应的简单结构。}
  \label{tab:nlp-extraction-kge-models}
\end{table*}

补全给出的是带分数的候选：链接预测排名靠前的三元组应以“模型推测”为来源类型写成待核验主张，
找到原文支持或人工确认后才转为事实。

\subsection{记录核验与更新}
\label{sec:nlp-extraction-update}

每次写入前都要检查模式、原文支持、对象身份、限定条件与来源完整性，
并区分两种时间。\term{有效时间}（valid time）是主张在所建模世界中适用的时间，
\term{写入时间}（transaction time）是系统获知或写入这一版本的时间。
结构检查复用第~\ref{sec:nlp-extraction-relation-records}节的JSON Schema与SHACL；
支持模型判为蕴含后，数字、否定与身份仍要逐项比较，这类错误常只差一个词或一位数字。
算法~\ref{alg:nlp-extraction-record-update}给出一份新材料到来时的处理流程。

\Needspace{18\baselineskip}
\begin{algorithm}[知识记录的核验与增量更新]
  \label{alg:nlp-extraction-record-update}
  \begin{algorithmic}[1]
    \Require 新来源$d$及其版本与发布时间，已有记录库$\mathcal{K}$，模式$S$，
      依赖图$D$（来源$\to$主张$\to$合并记录$\to$派生关系）
    \Ensure 更新后的$\mathcal{K}$与操作日志
    \State 从$d$抽取主张$c=(\text{身份},\text{字段},\text{值},\text{条件},\text{有效时间},\text{证据位置},\text{来源})$
    \For{每个主张$c$}
      \If{$c$不符合$S$，或证据位置与$d$的当前版本不对应}
        \State 标记为待核验并记录原因；\textbf{continue}
      \EndIf
      \State 核对支持：证据蕴含$c$，且数字、否定、日期与原文逐项一致；否则待核验
      \State 解析身份：可靠键直接匹配，否则按第~\ref{sec:nlp-extraction-alignment}节对齐；
        有歧义则待核验
      \State 取同一身份、字段、条件下的已有主张集合$P$
      \If{$P=\varnothing$}
        \State 新增$c$\Comment{补充}
      \ElsIf{$c$与$P$中某主张值相同}
        \State 合并来源列表，并检查新来源是否只是转载\Comment{重复}
      \ElsIf{$d$明确更正$P$中的主张$p$}
        \State 把$p$的有效期在更正生效时刻结束；新增$c$并连接“替代$p$”\Comment{修改}
      \ElsIf{$d$明确否认$p$}
        \State 把$p$标为撤回，保留历史\Comment{撤回}
      \Else
        \State 并存保存$c$与$P$，标记冲突待核验
      \EndIf
      \State 记录写入时间为当前时间，与$c$的有效时间分开保存
    \EndFor
    \If{$d$撤回了某个旧来源$d_0$}
      \State 沿$D$找出依赖$d_0$的主张与派生记录，重算支持；
        无独立支持者降级或撤回
    \EndIf
    \State 检查到期的有效期，将过期主张标为过期\Comment{过期}
  \end{algorithmic}
\end{algorithm}

以图~\ref{fig:nlp-extraction-records}的教学通知为例：4月14日旧材料把\code{ML-0418}排在4月17日，
4月15日更正为4月18日14:00。算法先按可靠活动编号确认是同一活动，再解析“原定”与“现改至”
指向同一排期字段，于是走“修改”分支：在“更正自发布日生效”的约定下，
旧排期的有效期为$[\text{4月14日},\text{4月15日})$，新排期为$[\text{4月15日},\text{尚未失效})$，
并保留改期事件与两份原文证据；原通知没有给出旧时刻，不应反填14:00。
活动开始时刻4月18日14:00是字段的值，不是主张有效期的起点。
若系统到4月16日才读到更正，写入时间为4月16日，有效时间仍从4月15日起；
4月17日才导入的旧通知走“并存待核验”分支，不会因写入较晚而恢复旧日期。

新材料未提某字段不等于撤回该字段，因此只有明确否认才进入撤回分支；
更正只涉及活动日期时，报名截止仍为4月16日18:00前，不自动顺延。
评价应同时报告事实正确性、覆盖、重复率、未决冲突与证据位置正确性。


\section*{文献说明}

本章使用两类资料。数据资料界定任务与监督对象：ACE定义了实体、关系与事件的标注体系；Few-NERD扩展了细粒度实体类型并提供少样本划分；DocRED把关系抽取推进到文档级；MAVEN扩大了事件类型覆盖，MAVEN-ERE在同一语料上标注了共指、时间、因果与子事件四类事件关系；HiEve标注了子事件层级，TimeBank-Dense与MATRES提供了稠密或按起点比较的时间关系。方法资料说明具体的预测机制：实体识别部分的Lattice-LSTM、FLAT、超图、W2NER、BARTNER、UIE与少样本近邻标注，实体链接部分的BLINK、全局一致性与GENRE，关系抽取部分的PURE、PCNN、CasRel、TPLinker、ATLOP与EIDER，事件抽取部分的OneIE、Text2Event与Doc2EDAG，以及知识构建部分的Ditto与TransE系列。正文中的教学模式、数值例子与额外的核验规则是本书为说明机制而设，以正文为准，不代表原论文的设置或结果。以下给出主要资料的完整出处。
\begin{fullwidth}
  \small
  \noindent\fullcite{doddington2004ace}\par\medskip
  \noindent\fullcite{ding2021fewnerd}\par\medskip
  \noindent\fullcite{nlp-analysis-zhang2018lattice}\par\medskip
  \noindent\fullcite{li2020flat}\par\medskip
  \noindent\fullcite{lu2015hypergraph}\par\medskip
  \noindent\fullcite{li2022w2ner}\par\medskip
  \noindent\fullcite{yan2021bartner}\par\medskip
  \noindent\fullcite{lu2022uie}\par\medskip
  \noindent\fullcite{yang2020structshot}\par\medskip
  \noindent\fullcite{wu2020blink}\par\medskip
  \noindent\fullcite{ganea2017deep}\par\medskip
  \noindent\fullcite{decao2021genre}\par\medskip
  \noindent\fullcite{yao2019docred}\par\medskip
  \noindent\fullcite{zhong2021pure}\par\medskip
  \noindent\fullcite{zeng2015pcnn}\par\medskip
  \noindent\fullcite{wei2020casrel}\par\medskip
  \noindent\fullcite{wang2020tplinker}\par\medskip
  \noindent\fullcite{zhou2021atlop}\par\medskip
  \noindent\fullcite{xie2022eider}\par\medskip
  \noindent\fullcite{wang2020maven}\par\medskip
  \noindent\fullcite{nlp-extraction-lin2020oneie}\par\medskip
  \noindent\fullcite{lu2021text2event}\par\medskip
  \noindent\fullcite{nlp-extraction-zheng2019doc2edag}\par\medskip
  \noindent\fullcite{wang2022mavenere}\par\medskip
  \noindent\fullcite{nlp-extraction-glavas2014hieve}\par\medskip
  \noindent\fullcite{nlp-extraction-cassidy2014dense}\par\medskip
  \noindent\fullcite{nlp-extraction-ning2018matres}\par\medskip
  \noindent\fullcite{nlp-extraction-li2020ditto}\par\medskip
  \noindent\fullcite{nlp-extraction-bordes2013transe}
\end{fullwidth}

\section{本章小结}
\label{sec:nlp-extraction-summary}

章首的挑战清单在各节得到了逐项回答。实体结构方面，标注、跨度、超图、词对网格与生成五条路线依次扩大可表达的结构：BIO--CRF只表示平面实体，跨度与阅读理解式方法加入嵌套，超图以线性代价编码重叠组合但带有结构歧义，W2NER用词对关系统一平面、嵌套与不连续实体，指针与结构化生成则用一个序列模型覆盖全部结构，代价是串行解码与原文对齐。中文词边界由Lattice-LSTM与FLAT把词典词作为可被语境否决的候选接入字符模型；新类型的标注稀缺由原型与近邻方法、大语言模型提示部分缓解，后者仍需格式核验与原文对齐。实体身份由别名先验、双编码召回与交叉编码重排、文档全局一致性和生成式名称预测逐步确定，库外实体则需要单独的拒绝判断。

关系与事件方面，联合抽取从统一标注、复制生成到级联指针和词对表格，逐步解决了三元组重叠；远程监督以多示例学习容忍噪声；文档级方法用图或上下文池化聚合分散的提及；事件抽取从流水线走向联合打分与结构生成，文档级方法处理多事件共享论元；事件之间的时间、因果与组成关系，以局部分类加全局约束保证一致。知识构建把这些结果对齐、合并、补全，并按来源与时间更新。

回到章首通知，4月17日应作为旧排期保留，4月18日14:00是当前活动开始时刻，4月16日18:00则限定校外报名。三项日期并不竞争同一个字段，只有先明确对象、角色与状态，才能回答何时举行、如何参加，以及答案依据是什么。下一章转向整篇文本的属性、文本之间的关系与文本集合的组织；本章得到的实体、关系和带来源的记录可以作为其中的结构化输入，其不确定性也应一并传递。
